% Managing processes — Computer Science Upper Sixth — Cours signé OnBuch+
% Official MINESEC syllabus. Compiler avec : tectonic CSC10-managing-processes.tex
\newcommand{\DOCMATIERE}{Computer Science}
\newcommand{\DOCNIVEAU}{Upper Sixth}
\newcommand{\DOCMODULE}{Upper Sixth — Computer Science}
\newcommand{\DOCLECON}{10}
\newcommand{\DOCTITRE}{Managing processes}
% OnBuch+ — préambule « cours signés » ENRICHI (compilé avec tectonic / XeLaTeX)
% Système visuel « Pop » d'OnBuch+ : fond crème (jamais blanc), bordures encre
% épaisses, ombres « dures » décalées, titres Archivo Black en capitales.
% Version MATHÉMATIQUES : graphes (pgfplots/tikz), boîtes pédagogiques
% (définition, propriété, méthode, attention, le savais-tu, exercice résolu,
% exercice, TP, bilan), page de garde et signature OnBuch+ ancrée en bas.
%
% Macros attendues dans main.tex AVANT \input{preamble} :
%   \DOCMATIERE  \DOCNIVEAU  \DOCMODULE  \DOCLECON (numéro)  \DOCTITRE
\documentclass[11pt]{article}

\usepackage{fontspec}
\usepackage[english]{babel}
\usepackage[a4paper,margin=2cm,top=3.4cm,bottom=2.8cm,headheight=1.4cm,footskip=1.3cm]{geometry}
\usepackage{amsmath,amssymb,amsfonts}
\usepackage{mathtools} % \xmapsto, \coloneqq... souvent utilisés par les modèles
\usepackage{cancel} % simplifications barrées (bilans de forces)
% \abs{z} (module, valeur absolue) et \norm{u} (norme), utilisés par les modèles
\providecommand{\abs}{}\let\abs\relax\DeclarePairedDelimiter{\abs}{\lvert}{\rvert}
\providecommand{\norm}{}\let\norm\relax\DeclarePairedDelimiter{\norm}{\lVert}{\rVert}
\usepackage[table]{xcolor} % [table] : fournit aussi \rowcolors (alternance de lignes)
\usepackage{pagecolor}
\usepackage{tikz}
\usetikzlibrary{arrows.meta,positioning,calc,shapes.geometric,shapes.misc,shapes.arrows,decorations.pathmorphing,decorations.pathreplacing,decorations.markings,patterns,shadows,mindmap,trees,fit,backgrounds,matrix,intersections,angles,quotes}
\usepackage{pgfplots}
\usepackage[version=4]{mhchem} % équations nucléaires \ce{^{235}_{92}U}
\usepackage[european,siunitx]{circuitikz} % schémas électriques
\pgfplotsset{compat=1.18}
\usepgfplotslibrary{fillbetween,groupplots} % aires entre courbes (calcul intégral)
\usepackage{enumitem}
\usepackage{fancyhdr}
% [most] charge toutes les bibliothèques tcolorbox (listings, minted, raster,
% documentation...) et gonfle inutilement la mémoire TeX sur un document long
% (~300 boîtes) : on ne charge que ce qui est réellement utilisé.
\usepackage[skins,breakable]{tcolorbox}
\usepackage{graphicx}
\usepackage{colortbl}
\usepackage{booktabs}
\usepackage{tabularx}
\usepackage{diagbox} % case d'en-tête diagonale des tableaux à double entrée
% Colonnes extensibles centrée / gauche / droite (C, L, R), souvent utilisées par les modèles
\newcolumntype{C}{>{\centering\arraybackslash}X}
\newcolumntype{L}{>{\raggedright\arraybackslash}X}
\newcolumntype{R}{>{\raggedleft\arraybackslash}X}
\usepackage{array}
\usepackage{multicol}
\usepackage{siunitx}
% parse-numbers=false : accepte aussi \SI{2,5 \times 10^{-3}}{...}
\sisetup{locale=UK,output-decimal-marker={.},group-minimum-digits=4,parse-numbers=false}
\DeclareSIUnit\ohm{\ensuremath{\symup{\Omega}}} % U+2126 absent de la police
\DeclareSIUnit\division{div} % réglages d'oscilloscope (ms/div, V/div)
\DeclareSIUnit\tr{tr} % tours (vitesse de rotation, ex. \SI{1500}{\tr\per\minute})
\usepackage{qrcode}
% \degree : commande fréquemment utilisée par les modèles pour un angle ou une
% température en dehors de \SI{}{} (non fournie par les paquets chargés).
\providecommand{\degree}{^{\circ}}
% \qed (fin de démonstration), utilisé par les modèles sans amsthm
\providecommand{\qed}{\hfill$\square$}
% \barre : double barre des valeurs interdites dans un tableau de variations (tkz-tab)
\providecommand{\barre}{\ensuremath{\|}}
% environnement proof (amsthm non chargé) utilisé par les modèles
\ifdefined\proof\else\newenvironment{proof}[1][Proof]{\par\noindent\textit{#1.}\ }{\qed\par}\fi
\usepackage{lastpage}
\usepackage{hyperref}
\hypersetup{hidelinks}

% ============================================================
% Palette « Pop » OnBuch+ — mêmes hex que la classe P (plus_ui.dart)
% ============================================================
\definecolor{poporange}{HTML}{F59321}
\definecolor{popdark}{HTML}{E07A0C}
\definecolor{popcream}{HTML}{FFF6E8}
\definecolor{popink}{HTML}{1C1714}
\definecolor{poppurple}{HTML}{7A5AE0}
\definecolor{popgreen}{HTML}{2E9E6A}
\definecolor{poppink}{HTML}{E0457B}
\definecolor{popblue}{HTML}{2F7DE1}
\definecolor{popgold}{HTML}{C98A12}
\definecolor{popmuted}{HTML}{8A7F76}
\definecolor{popline}{HTML}{EADFD2}
\definecolor{popgray}{HTML}{8A7F76}
\colorlet{popgrey}{popgray}
% Alias tolérants (le modèle nomme parfois les couleurs différemment)
\colorlet{popwhite}{white}
\colorlet{popblack}{popink}
\colorlet{popred}{poppink}
\colorlet{popyellow}{popgold}
% Teintes claires pour fonds de tableaux / zones
\colorlet{popblueL}{popblue!10!white}
\colorlet{poporangeL}{poporange!14!white}
\colorlet{popgreenL}{popgreen!12!white}
\colorlet{poppurpleL}{poppurple!12!white}
\colorlet{poppinkL}{poppink!12!white}
\colorlet{popgoldL}{popgold!14!white}

\pagecolor{popcream}

% ============================================================
% Typographie
% ============================================================
\defaultfontfeatures{Ligatures=TeX}
\setmainfont{PlusJakartaSans-Medium.ttf}[Path=../../../fonts/,BoldFont=PlusJakartaSans-Bold.ttf]
% Maths en sans-serif (Fira Math) avec chiffres et lettres droites de Plus
% Jakarta, pour que les formules se fondent dans le texte.
\usepackage[math-style=ISO,bold-style=ISO]{unicode-math}
\setmathfont{FiraMath-Regular.otf}[Path=../../../fonts/]
\setmathfont{PlusJakartaSans-Medium.ttf}[Path=../../../fonts/,range={up/num,up/{Latin,latin},\mathup}]
% (icomma retiré : décimales avec point en anglais)
% Caractères mathématiques Unicode absents des polices de texte (Plus Jakarta,
% Archivo Black) : rendus par leur équivalent mathématique, en texte comme en
% formule — sinon ils s'affichent en carré vide (ex. « Arithmétique dans ℤ »).
\usepackage{newunicodechar}
\newunicodechar{ℝ}{\ensuremath{\mathbb{R}}}
\newunicodechar{ℤ}{\ensuremath{\mathbb{Z}}}
\newunicodechar{ℂ}{\ensuremath{\mathbb{C}}}
\newunicodechar{ℕ}{\ensuremath{\mathbb{N}}}
\newunicodechar{ℚ}{\ensuremath{\mathbb{Q}}}
\newunicodechar{𝔻}{\ensuremath{\mathbb{D}}}
\newunicodechar{α}{\ensuremath{\alpha}}
\newunicodechar{β}{\ensuremath{\beta}}
\newunicodechar{γ}{\ensuremath{\gamma}}
\newunicodechar{δ}{\ensuremath{\delta}}
\newunicodechar{ε}{\ensuremath{\varepsilon}}
\newunicodechar{θ}{\ensuremath{\theta}}
\newunicodechar{λ}{\ensuremath{\lambda}}
\newunicodechar{μ}{\ensuremath{\mu}}
\newunicodechar{σ}{\ensuremath{\sigma}}
\newunicodechar{τ}{\ensuremath{\tau}}
\newunicodechar{φ}{\ensuremath{\varphi}}
\newunicodechar{ψ}{\ensuremath{\psi}}
\newunicodechar{ω}{\ensuremath{\omega}}
\newunicodechar{Σ}{\ensuremath{\Sigma}}
% majuscules produites par \MakeUppercase dans les titres (α-aminés -> Α)
\newunicodechar{Α}{\ensuremath{\alpha}}
\newunicodechar{Β}{\ensuremath{\beta}}
\newunicodechar{Γ}{\ensuremath{\Gamma}}
\newunicodechar{Δ}{\ensuremath{\Delta}}
\newunicodechar{Ε}{\ensuremath{\varepsilon}}
\newunicodechar{Θ}{\ensuremath{\Theta}}
\newunicodechar{Λ}{\ensuremath{\Lambda}}
\newunicodechar{Μ}{\ensuremath{\mu}}
\newunicodechar{Τ}{\ensuremath{\tau}}
\newunicodechar{Φ}{\ensuremath{\Phi}}
\newunicodechar{Ψ}{\ensuremath{\Psi}}
\newunicodechar{Ω}{\ensuremath{\Omega}}
\newunicodechar{↦}{\ensuremath{\mapsto}}
\newunicodechar{∈}{\ensuremath{\in}}
\newunicodechar{∉}{\ensuremath{\notin}}
\newunicodechar{∧}{\ensuremath{\wedge}}
\newunicodechar{⇔}{\ensuremath{\Leftrightarrow}}
\newunicodechar{⟺}{\ensuremath{\Longleftrightarrow}}
\newunicodechar{⇒}{\ensuremath{\Rightarrow}}
\newunicodechar{⟹}{\ensuremath{\Longrightarrow}}
\newunicodechar{∘}{\ensuremath{\circ}}
\newunicodechar{∪}{\ensuremath{\cup}}
\newunicodechar{∩}{\ensuremath{\cap}}
\newunicodechar{≡}{\ensuremath{\equiv}}
\newunicodechar{⊂}{\ensuremath{\subset}}
\newunicodechar{⊥}{\ensuremath{\perp}}
\newunicodechar{∃}{\ensuremath{\exists}}
\newunicodechar{∀}{\ensuremath{\forall}}
\newunicodechar{∛}{\ensuremath{\sqrt[3]{\,}}}
\newunicodechar{∜}{\ensuremath{\sqrt[4]{\,}}}
\newunicodechar{⃗}{\ensuremath{{}^{\to}}}
\newunicodechar{ⁿ}{\ifmmode^{n}\else\textsuperscript{\ensuremath{n}}\fi}
\newunicodechar{ˣ}{\ifmmode^{x}\else\textsuperscript{\ensuremath{x}}\fi}
\newunicodechar{ᵇ}{\ifmmode^{b}\else\textsuperscript{\ensuremath{b}}\fi}
\newunicodechar{ʳ}{\ifmmode^{r}\else\textsuperscript{\ensuremath{r}}\fi}
\newunicodechar{ᵘ}{\ifmmode^{u}\else\textsuperscript{\ensuremath{u}}\fi}
\newunicodechar{ʸ}{\ifmmode^{y}\else\textsuperscript{\ensuremath{y}}\fi}
\newunicodechar{ᵃ}{\ifmmode^{a}\else\textsuperscript{\ensuremath{a}}\fi}
\newunicodechar{ᵉ}{\ifmmode^{e}\else\textsuperscript{\ensuremath{e}}\fi}
\newunicodechar{ᵏ}{\ifmmode^{k}\else\textsuperscript{\ensuremath{k}}\fi}
\newunicodechar{ᵖ}{\ifmmode^{p}\else\textsuperscript{\ensuremath{p}}\fi}
\newunicodechar{ᴺ}{\ifmmode^{N}\else\textsuperscript{\ensuremath{N}}\fi}
\newunicodechar{ᶜ}{\ifmmode^{c}\else\textsuperscript{\ensuremath{c}}\fi}
\newunicodechar{ʰ}{\ifmmode^{h}\else\textsuperscript{\ensuremath{h}}\fi}
\newunicodechar{ᵠ}{\ifmmode^{q}\else\textsuperscript{\ensuremath{q}}\fi}
\newunicodechar{ₙ}{\ifmmode_{n}\else\textsubscript{\ensuremath{n}}\fi}
\newunicodechar{ₐ}{\ifmmode_{a}\else\textsubscript{\ensuremath{a}}\fi}
\newunicodechar{ᵢ}{\ifmmode_{i}\else\textsubscript{\ensuremath{i}}\fi}
\newunicodechar{ᵣ}{\ifmmode_{r}\else\textsubscript{\ensuremath{r}}\fi}
\newunicodechar{ₖ}{\ifmmode_{k}\else\textsubscript{\ensuremath{k}}\fi}
\newunicodechar{ᵦ}{\ifmmode_{\beta}\else\textsubscript{\ensuremath{\beta}}\fi}
\newunicodechar{ₓ}{\ifmmode_{x}\else\textsubscript{\ensuremath{x}}\fi}
\newfontfamily\popdisplay{ArchivoBlack-Regular.ttf}[Path=../../../fonts/]
\newfontfamily\poplogo{SpaceGrotesk-Bold.ttf}[Path=../../../fonts/]
\newfontfamily\popsemi{PlusJakartaSans-SemiBold.ttf}[Path=../../../fonts/]
\newfontfamily\popxbold{PlusJakartaSans-ExtraBold.ttf}[Path=../../../fonts/]
\newfontfamily\popxboldls{PlusJakartaSans-ExtraBold.ttf}[Path=../../../fonts/,LetterSpace=3.0]

\setlist[enumerate]{leftmargin=1.7em,itemsep=5pt,topsep=4pt}
\setlist[itemize]{leftmargin=1.7em,itemsep=4pt,topsep=4pt,label={\color{poporange}\textbullet}}
\setlength{\parindent}{0pt}
\setlength{\parskip}{5pt}
\renewcommand{\arraystretch}{1.25}

% pgfplots : style Pop par défaut
\pgfplotsset{
  every axis/.append style={axis line style={popink,line width=1pt},tick style={popink},
    grid style={popline,line width=0.5pt},label style={font=\small},tick label style={font=\footnotesize}},
  popcurve/.style={poporange,line width=1.8pt,no markers,smooth},
}
% Même style disponible dans un \draw TikZ simple (hors pgfplots)
\tikzset{popcurve/.style={poporange,line width=1.8pt}}
\tikzset{popgrid/.style={popline,very thin}}
\colorlet{popcurve}{poporange} % le nom du style est parfois utilisé comme couleur

% ============================================================
% Pill / squiggle
% ============================================================
\newcommand{\poppill}[3]{%
  \tikz[baseline=-0.55ex]\node[draw=popink,line width=1.2pt,rounded corners=2.6pt,%
    fill=#2,inner xsep=6.5pt,inner ysep=3pt]{%
    \popxboldls\fontsize{7}{7}\selectfont\color{#3}\MakeUppercase{#1}};%
}
\newcommand{\popsquiggle}[1]{%
  \tikz[baseline=-0.4ex]{
    \draw[#1,line width=2.6pt,line cap=round]
      plot [smooth] coordinates {(0,0.09) (0.34,0.01) (0.68,0.21) (1.02,0.01) (1.36,0.10)};
  }%
}
% Mot-clé surligné (marqueur orange sous le texte)
\newcommand{\cle}[1]{{\popxbold\color{popdark}#1}}

% ============================================================
% En-tête / pied
% ============================================================
\pagestyle{fancy}
\fancyhf{}
\renewcommand{\headrulewidth}{0pt}
\renewcommand{\footrulewidth}{0pt}
\lhead{\raisebox{-0.15ex}{\poplogo\fontsize{13}{13}\selectfont\color{popink}OnBuch\color{poporange}+}}
\rhead{\poppill{\DOCMATIERE\ \textbullet\ \DOCNIVEAU\ \textbullet\ Chapter \DOCLECON}{white}{popink}}
\lfoot{\popxbold\fontsize{7.6}{7.6}\selectfont\color{popmuted}\MakeUppercase{Signed course OnBuch+}\ \textbullet\ {\popsemi\DOCTITRE}}
\rfoot{\tikz[baseline=-0.6ex]\node[rounded corners=8.5pt,fill=popink,minimum height=17pt,minimum width=17pt,inner xsep=5pt,inner sep=0]{\popxbold\fontsize{8}{8}\selectfont\color{popcream}\thepage\,/\,\pageref*{LastPage}};}

% ============================================================
% Titres
% ============================================================
\newcommand{\chaptitre}[2]{%
  \begin{tcolorbox}[enhanced,colback=poporange,colframe=popink,boxrule=2.6pt,arc=11pt,
      drop shadow=popink,left=15pt,right=15pt,top=12pt,bottom=11pt,before skip=3pt,after skip=18pt]
    \poppill{#2}{popink}{popcream}\par\vspace{9pt}
    {\raggedright\hyphenpenalty=10000\exhyphenpenalty=10000\popdisplay\fontsize{22}{24}\selectfont\color{popcream}\MakeUppercase{#1}\par}
  \end{tcolorbox}
}
% Section numérotée : pastille orange avec le numéro + titre Archivo
\newcounter{popsec}
\newcommand{\coursec}[1]{%
  \stepcounter{popsec}\setcounter{popsub}{0}%
  \par\vspace{16pt}\noindent
  \begin{minipage}{\linewidth}
  \tikz[baseline=-0.7ex]\node[draw=popink,line width=1.6pt,fill=poporange,rounded corners=4pt,minimum size=22pt,inner sep=0]{\popdisplay\fontsize{12}{12}\selectfont\color{popcream}\thepopsec};%
  \hspace{8pt}{\popdisplay\fontsize{15}{17}\selectfont\color{popink}\MakeUppercase{#1}}\par
  \vspace{4pt}\noindent{\color{poporange}\rule{36pt}{3.6pt}}
  \end{minipage}\par\nopagebreak\vspace{8pt}
}
\newcounter{popsub}
\newcommand{\courssub}[1]{%
  \stepcounter{popsub}%
  \par\vspace{9pt}\noindent{\popxbold\fontsize{11.5}{13}\selectfont\color{popink}{\color{poporange}\thepopsec.\thepopsub}\hspace{6pt}#1}\par\nopagebreak\vspace{4pt}
}

% ============================================================
% Boîtes pédagogiques — même recette : fond blanc, bordure épaisse colorée,
% ombre dure encre, badge plein. Titre optionnel [#1].
% ============================================================
\tcbset{popbase/.style={breakable,enhanced,colback=white,boxrule=2.3pt,arc=9pt,
  shadow={3pt}{-3pt}{0pt}{popink},left=14pt,right=14pt,top=11pt,bottom=11pt,before skip=11pt,after skip=15pt,
  fontupper=\popsemi\fontsize{10.6}{14.5}\selectfont\color{popink},
  fontlower=\popsemi\fontsize{10.6}{14.5}\selectfont\color{popink}}}
% Coche dessinée (le glyphe ✓ n'existe pas dans les polices du document)
\newcommand{\popcheck}[1]{\tikz[baseline=-0.5ex]\draw[#1,line width=1.6pt,line cap=round,line join=round](-0.13,0.0)--(-0.04,-0.09)--(0.13,0.1);}
\providecommand{\coche}{\popcheck{popgreen}}
\providecommand{\resultat}[1]{\par\smallskip\noindent\fcolorbox{popgreen}{popgreenL}{\parbox{\dimexpr\linewidth-2\fboxsep-2\fboxrule\relax}{\textbf{Result.} #1}}\par\smallskip}
\newcommand{\popboxhead}[3]{\poppill{#1}{#2}{white}\ifx\relax#3\relax\else\hspace{6pt}{\popxbold\fontsize{10.6}{12}\selectfont\color{popink}#3}\fi\par\vspace{7pt}}

\newtcolorbox{aretenir}[1][]{popbase,colframe=popgreen,before upper={\popboxhead{Key points}{popgreen}{#1}}}
\newtcolorbox{exemplebox}[1][]{popbase,colframe=poporange,before upper={\popboxhead{Exemple}{poporange}{#1}}}
\newtcolorbox{definition}[1][]{popbase,colframe=popblue,before upper={\popboxhead{Definition}{popblue}{#1}}}
\newtcolorbox{propriete}[1][]{popbase,colframe=poppurple,before upper={\popboxhead{Property}{poppurple}{#1}}}
\newtcolorbox{methode}[1][]{popbase,colframe=popgold,before upper={\popboxhead{Method}{popgold}{#1}}}
\newtcolorbox{attention}[1][]{popbase,colframe=poppink,before upper={\popboxhead{Warning}{poppink}{#1}}}
\newtcolorbox{experience}[1][]{popbase,colframe=popblue,colback=popblueL,before upper={\popboxhead{Experiment}{popblue}{#1}}}
% « Le savais-tu ? » : carte violette pleine (contexte, culture, Cameroun)
\newtcolorbox{savaistu}[1][]{popbase,colback=poppurple,colframe=popink,
  fontupper=\popsemi\fontsize{10.4}{14.2}\selectfont\color{white},
  before upper={\poppill{Did you know?}{white}{poppurple}\ifx\relax#1\relax\else\hspace{6pt}{\popxbold\color{white}#1}\fi\par\vspace{7pt}}}
% Exercice résolu : énoncé en haut, \tcblower puis la solution sur fond crème
\newcounter{popexo}\newcounter{popexr}
\newtcolorbox{exoresolu}[1][]{popbase,colframe=poporange,colbacklower=popcream,
  segmentation style={popink,line width=1.2pt,dashed},
  before upper={\stepcounter{popexr}\popboxhead{Worked example \thepopexr}{poporange}{#1}},
  before lower={\poppill{Solution}{popink}{popcream}\par\vspace{6pt}}}
% Exercice (non résolu en place) — la correction est dans la partie Corrigés
\newtcolorbox{exercice}[1][]{popbase,colframe=popink,
  before upper={\stepcounter{popexo}\popboxhead{Exercise \thepopexo}{popink}{#1}}}
% Corrigé
\newtcolorbox{corrige}[1][]{popbase,colframe=popgreen,colback=popgreenL,
  before upper={\popboxhead{Solution}{popgreen}{#1}}}
% Objectifs / prérequis / bilan
\newtcolorbox{objectifs}{popbase,colframe=popink,colback=poporangeL,
  before upper={\popboxhead{Chapter objectives}{popink}{}}}
\newtcolorbox{prerequis}{popbase,colframe=popink,colback=popblueL,
  before upper={\popboxhead{Prerequisites}{popblue}{}}}
\newtcolorbox{bilan}{popbase,colframe=popink,colback=white,boxrule=2.8pt,
  before upper={\popboxhead{Fiche bilan}{poporange}{L'essentiel à connaître pour l'examen}}}
% Figure encadrée avec légende
\newtcolorbox{popfigure}[1][]{enhanced,breakable=false,colback=white,colframe=popline,boxrule=1.4pt,arc=8pt,
  left=10pt,right=10pt,top=10pt,bottom=8pt,before skip=10pt,after skip=12pt,halign=center,
  lower separated=false,halign lower=center,
  fontlower=\popsemi\fontsize{9}{11.5}\selectfont\color{popmuted},
  #1}
\newcommand{\legende}[1]{\par\vspace{6pt}{\popsemi\fontsize{9}{11.5}\selectfont\color{popmuted}#1}}

% ============================================================
% Signature OnBuch+
% ============================================================
\newcommand{\poplogoinline}[1]{{\poplogo\fontsize{#1}{#1}\selectfont\color{popink}OnBuch\color{poporange}+}}
\newcommand{\signatureonbuch}{%
  \begin{tcolorbox}[enhanced,colback=popink,colframe=popink,boxrule=2.2pt,arc=10pt,
      drop shadow=poporange,left=14pt,right=14pt,top=10pt,bottom=10pt,before skip=0pt,after skip=0pt]
    \tikz[baseline=-0.7ex]\node[circle,fill=popgreen,draw=popcream,line width=1.3pt,minimum size=22pt,inner sep=0]{\popcheck{white}};%
    \hspace{9pt}\begin{minipage}[c]{0.62\linewidth}
      {\popxbold\fontsize{12}{13}\selectfont\color{popcream}Signed\ \textbullet\ {\poplogo OnBuch\color{poporange}+}}\par\vspace{2pt}
      {\raggedright\popsemi\fontsize{8}{10.5}\selectfont\color{popline}\MakeUppercase{OnBuch+ teaching team}\\\MakeUppercase{Follows the official MINESEC syllabus}\par}
    \end{minipage}\hfill
    {\poplogo\fontsize{22}{22}\selectfont\color{popcream}OnBuch\color{poporange}+}
  \end{tcolorbox}%
}

% ============================================================
% Page de garde
% #1 titre · #2 matière · #3 niveau · #4 module · #5 n° leçon · #6 durée ·
% #7 description / objectifs (texte ou itemize)
% ============================================================
\newcommand{\pagegarde}[7]{%
  \thispagestyle{empty}
  \newgeometry{a4paper,margin=1.7cm,top=1.6cm,bottom=1.4cm}%
  \begin{tikzpicture}[remember picture,overlay]
    \fill[poporange!18] ([xshift=-3.2cm,yshift=-3.6cm]current page.north east) circle (4.6cm);
    \fill[poppurple!14] ([xshift=2.4cm,yshift=7.6cm]current page.south west) circle (3.8cm);
  \end{tikzpicture}%
  {\poplogo\fontsize{30}{30}\selectfont\color{popink}OnBuch\color{poporange}+}\hfill
  \raisebox{0.6ex}{\poppill{Signed course \textbullet\ Premium}{popink}{popcream}}\par
  \vspace{1pt}\popsquiggle{poppurple}\par
  \vspace{8pt}
  \begin{tcolorbox}[enhanced,colback=poporange,colframe=popink,boxrule=2.8pt,arc=13pt,
      drop shadow=popink,left=18pt,right=18pt,top=12pt,bottom=13pt,after skip=10pt]
    \poppill{#2 \textbullet\ #3}{popink}{popcream}\hspace{5pt}\poppill{#4}{white}{popink}\par\vspace{8pt}
    {\popdisplay\fontsize{14}{14}\selectfont\color{popink}CHAPTER #5}\par\vspace{4pt}
    {\raggedright\hyphenpenalty=10000\exhyphenpenalty=10000\popdisplay\fontsize{25}{27}\selectfont\color{popcream}\MakeUppercase{#1}\par}
    \vspace{8pt}
    {\popxbold\fontsize{9.5}{9.5}\selectfont\color{popink}\MakeUppercase{Suggested time: #6}}
  \end{tcolorbox}
  \begin{tcolorbox}[enhanced,colback=white,colframe=popink,boxrule=2.2pt,arc=10pt,
      drop shadow=popink,left=16pt,right=16pt,top=10pt,bottom=10pt,after skip=8pt]
    \poppill{In this chapter}{poppurple}{white}\par\vspace{5pt}
    \popsemi\fontsize{9.3}{12.6}\selectfont\color{popink}#7
  \end{tcolorbox}
  \noindent\begin{minipage}[t]{0.487\linewidth}
    \begin{tcolorbox}[enhanced,colback=popgreen,colframe=popink,boxrule=2.2pt,arc=10pt,
        drop shadow=popink,left=11pt,right=11pt,top=10pt,bottom=10pt,height=3.1cm,valign=center]
      \tcbox[colback=white,colframe=popink,boxrule=1.3pt,arc=5pt,boxsep=0pt,nobeforeafter,
        left=3pt,right=3pt,top=3pt,bottom=3pt,valign=center]{\qrcode[height=1.9cm]{https://play.google.com/store/apps/details?id=com.luvvix.onbuch&hl=en}}%
      \hspace{8pt}\begin{minipage}[c]{0.5\linewidth}
        {\popdisplay\fontsize{11}{12.5}\selectfont\color{white}\MakeUppercase{Download OnBuch}}\par\vspace{4pt}
        {\raggedright\popsemi\fontsize{8.4}{11}\selectfont\color{white}Free \textbullet\ Google Play\\AI tutor, past papers, results\par}
      \end{minipage}
    \end{tcolorbox}
  \end{minipage}\hfill
  \begin{minipage}[t]{0.487\linewidth}
    \begin{tcolorbox}[enhanced,colback=poppurple,colframe=popink,boxrule=2.2pt,arc=10pt,
        drop shadow=popink,left=11pt,right=11pt,top=10pt,bottom=10pt,height=3.1cm,valign=center]
      \tikz[baseline=-0.7ex]\node[circle,fill=white,draw=popink,line width=1.3pt,minimum size=1.6cm,inner sep=0]{\popdisplay\fontsize{24}{24}\selectfont\color{poppurple}+};%
      \hspace{8pt}\begin{minipage}[c]{0.6\linewidth}
        {\popdisplay\fontsize{11}{12.5}\selectfont\color{white}\MakeUppercase{Go OnBuch+}}\par\vspace{4pt}
        {\raggedright\popsemi\fontsize{8.4}{11}\selectfont\color{white}All signed courses, unlimited AI tutor, PDF sheets — OnBuch+ tab in the app\par}
      \end{minipage}
    \end{tcolorbox}
  \end{minipage}
  \vspace{8pt}
  \popcontenu
  \vfill
  \signatureonbuch
  \restoregeometry
  \newpage
}

% Rangée « Ce cours contient » (page de garde)
\newcommand{\poptile}[3]{% #1 couleur #2 gros texte #3 légende
  \begin{tcolorbox}[enhanced,colback=#1,colframe=popink,boxrule=2pt,arc=9pt,shadow={2.5pt}{-2.5pt}{0pt}{popink},
      left=6pt,right=6pt,top=9pt,bottom=9pt,halign=center,valign=center,height=1.75cm,width=\linewidth,nobeforeafter]
    {\popdisplay\fontsize{12.5}{13}\selectfont\color{white}\MakeUppercase{#2}}\par\vspace{3pt}
    {\centering\popsemi\fontsize{7.8}{9.5}\selectfont\color{white}#3\par}
  \end{tcolorbox}}
\newcommand{\popcontenu}{%
  \par\noindent{\popxboldls\fontsize{7.5}{7.5}\selectfont\color{popmuted}THIS SIGNED COURSE CONTAINS}\par\vspace{7pt}
  \noindent\begin{minipage}[t]{0.235\linewidth}\poptile{popblue}{Course}{complete, illustrated, step by step}\end{minipage}\hfill
  \begin{minipage}[t]{0.235\linewidth}\poptile{poporange}{Methods}{and worked examples}\end{minipage}\hfill
  \begin{minipage}[t]{0.235\linewidth}\poptile{poppink}{Exercises}{exam style + solutions}\end{minipage}\hfill
  \begin{minipage}[t]{0.235\linewidth}\poptile{popgreen}{Summary sheet}{the essentials to remember}\end{minipage}\par}

% Sommaire
\newtcolorbox{sommaire}{enhanced,colback=white,colframe=popink,boxrule=2.2pt,arc=10pt,
  drop shadow=popink,left=18pt,right=18pt,top=13pt,bottom=13pt,before skip=6pt,after skip=20pt,
  before upper={\poppill{Contents}{poppurple}{white}\par\vspace{9pt}\popsemi\fontsize{10.6}{14.5}\selectfont\color{popink}}}

% Signature de fin de cours — poussée tout en bas de la dernière page
\newcommand{\findecours}{%
  \par\vfill\nopagebreak
  \begin{center}\popsquiggle{poporange}\end{center}\vspace{4pt}
  \signatureonbuch
}
\AtBeginDocument{\def\llbracket{\mathopen{[\mkern-3.5mu[}}\def\rrbracket{\mathclose{]\mkern-3.5mu]}}} % intervalles d'entiers (glyphe ⟦ absent de la police)
% Glyphes absents de Fira Math : redéfinis à partir de glyphes disponibles
\AtBeginDocument{%
  \def\setminus{\mathbin{\backslash}}\let\smallsetminus\setminus
  \def\vdots{\mathord{\vbox{\baselineskip3.2pt\lineskiplimit0pt\kern4pt\hbox{.}\hbox{.}\hbox{.}}}}%
  \def\ddots{\mathinner{\mkern1mu\raise6.5pt\hbox{.}\mkern2mu\raise3.6pt\hbox{.}\mkern2mu\raise0.7pt\hbox{.}\mkern1mu}}%
  \def\triangle{\mathord{\scalebox{0.9}{$\symup{\Delta}$}}}%
  \def\nabla{\mathord{\raisebox{\depth}{\scalebox{1}[-1]{$\symup{\Delta}$}}}}%
  \def\checkmark{\popcheck{popdark}}%
  \def\star{\mathbin{\ast}}\def\bigstar{\mathord{\ast}}%
  \def\diamond{\mathbin{\scalebox{0.6}{$\Diamond$}}}%
  \def\bigcup{\mathop{\mathchoice{\raisebox{-0.3em}{\scalebox{1.7}{$\cup$}}}{\scalebox{1.25}{$\cup$}}{\cup}{\cup}}\displaylimits}%
  \def\bigcap{\mathop{\mathchoice{\raisebox{-0.3em}{\scalebox{1.7}{$\cap$}}}{\scalebox{1.25}{$\cap$}}{\cap}{\cap}}\displaylimits}%
  \def\lesseqqgtr{\mathrel{\lessgtr}}\def\bigoplus{\mathop{\oplus}\displaylimits}%
}
\providecommand{\ssi}{if and only if} % abréviation fréquente
\AtBeginDocument{\providecommand{\wideparen}{\overparen}} % arc de cercle

% Fonctions usuelles que les modèles utilisent sans que amsmath les fournisse
\providecommand{\arsinh}{\operatorname{arsinh}}\providecommand{\arcosh}{\operatorname{arcosh}}
\providecommand{\artanh}{\operatorname{artanh}}\providecommand{\arsech}{\operatorname{arsech}}
\providecommand{\arcsch}{\operatorname{arcsch}}\providecommand{\arcoth}{\operatorname{arcoth}}
\providecommand{\arcsinh}{\operatorname{arsinh}}\providecommand{\arccosh}{\operatorname{arcosh}}
\providecommand{\arctanh}{\operatorname{artanh}}\providecommand{\sech}{\operatorname{sech}}
\providecommand{\csch}{\operatorname{csch}}\providecommand{\arcsec}{\operatorname{arcsec}}
\providecommand{\arccsc}{\operatorname{arccsc}}\providecommand{\arccot}{\operatorname{arccot}}
\providecommand{\sgn}{\operatorname{sgn}}\providecommand{\lcm}{\operatorname{lcm}}
\providecommand{\Var}{\operatorname{Var}}\providecommand{\Cov}{\operatorname{Cov}}

%% FIGFIX-BEGIN v5 — correctifs globaux des figures (docs/onbuch-generation/tools/figfix)
\usepackage{adjustbox}\usepackage{etoolbox}\tcbuselibrary{raster}
% toute figure TikZ/pgfplots plus large que la ligne est réduite à la largeur disponible
\BeforeBeginEnvironment{tikzpicture}{\begin{adjustbox}{max width=\linewidth}}
\AfterEndEnvironment{tikzpicture}{\end{adjustbox}}
% légendes pgfplots lisibles par-dessus les courbes
\pgfplotsset{every axis/.append style={legend style={fill=white,fill opacity=0.92,text opacity=1,draw=black!15,inner sep=3pt}}}
% étiquettes d'arêtes (syntaxe edge["..."]) avec fond blanc
\tikzset{every edge quotes/.append style={fill=white,inner sep=1.2pt}}
%% FIGFIX-END
\begin{document}

\pagegarde{Managing processes}{Computer Science}{Upper Sixth}{Upper Sixth — Computer Science}{10}{7 periods}{This chapter explains how an operating system creates, schedules and coordinates processes so that many programs appear to run simultaneously on a single processor. Students master scheduling algorithms, synchronisation tools and deadlock handling — core topics of the GCE Advanced Level Computer Science paper.
\vspace{4pt}
\begin{multicols}{2}\small
\begin{itemize}
\item By the end of this chapter I can define a process and distinguish it from a program
\item By the end of this chapter I can describe the process states and draw the state transition diagram
\item By the end of this chapter I can explain the role of the Process Control Block (PCB) and of context switching
\item By the end of this chapter I can distinguish preemptive from non-preemptive scheduling
\item By the end of this chapter I can compute turnaround time, waiting time and draw Gantt charts for FCFS, SJF, SRTF, Round Robin and priority scheduling
\item By the end of this chapter I can compare scheduling algorithms using performance criteria and choose the appropriate one for a given situation
\end{itemize}\end{multicols}}

\begin{sommaire}
\begin{multicols}{2}
\begin{enumerate}
\item Section 1: Process Management Basics (Lesson 44)
\item Section 2: Process Scheduling and Non-preemptive Algorithms (Lessons 45–46)
\item Section 3: Preemptive Scheduling (Lesson 47)
\item Section 4: Process Synchronisation (Lesson 48)
\item Section 5: Deadlock Handling and Integration (Lessons 49–50)
\item Integration activity
\item Methods and worked examples
\item Exercises
\item Solutions to the exercises
\item Summary sheet
\end{enumerate}
\end{multicols}
\end{sommaire}

\begin{objectifs}
\begin{itemize}
\item By the end of this chapter I can define a process and distinguish it from a program
\item By the end of this chapter I can describe the process states and draw the state transition diagram
\item By the end of this chapter I can explain the role of the Process Control Block (PCB) and of context switching
\item By the end of this chapter I can distinguish preemptive from non-preemptive scheduling
\item By the end of this chapter I can compute turnaround time, waiting time and draw Gantt charts for FCFS, SJF, SRTF, Round Robin and priority scheduling
\item By the end of this chapter I can compare scheduling algorithms using performance criteria and choose the appropriate one for a given situation
\item By the end of this chapter I can explain race conditions, critical sections and mutual exclusion
\item By the end of this chapter I can use semaphores to solve classical synchronisation problems (producer–consumer)
\item By the end of this chapter I can state the four deadlock conditions and apply prevention, avoidance (Banker's algorithm) and detection strategies
\end{itemize}
\end{objectifs}

\begin{prerequis}
\begin{itemize}
\item Role and functions of an operating system (Lower Sixth: introduction to system software)
\item Basic hardware structure: CPU, registers, main memory, interrupts
\item Multiprogramming and time-sharing concepts seen in the OS overview chapter
\item Elementary algorithmics: queues (FIFO), averages and simple table computations
\item Notion of concurrent users from the databases/networks chapters
\end{itemize}
\end{prerequis}

\newpage

\coursec{Section 1: Process Management Basics (Lesson 44)}

\courssub{The Hook: A Mobile Money Kiosk in Douala}

Imagine a busy mobile money kiosk in Douala. One agent serves dozens of customers: some transactions are quick airtime top-ups, others are long money transfers requiring multiple confirmations. If the agent serves strictly in arrival order, a single long transfer can make twenty customers wait; if the agent keeps interrupting customers to be ``fair'', some may never finish because of the constant restarting. Worse, imagine two customers each holding one of the two forms the other needs — nobody moves. A computer's processor faces exactly the same dilemma every millisecond with hundreds of competing \cle{processes}. How does the operating system decide who runs, for how long, and how does it prevent total gridlock? This section builds the fundamental vocabulary and mechanisms the OS uses to answer those questions.

\courssub{What Is a Process?}

\begin{definition}[Program vs. Process]
A \cle{program} is a passive, static entity: a file containing executable code and data stored on secondary storage (e.g., \texttt{/bin/ls} or \texttt{chrome.exe}). A \cle{process} is an active, dynamic entity: a program in execution. It includes the program code (text section), current values of variables (data section), the execution stack (function calls, local variables), the heap (dynamic memory), and the \cle{process context} — the content of the CPU registers (program counter, stack pointer, general-purpose registers, status flags) at a given instant.
\end{definition}

\begin{attention}[Program $\neq$ Process]
Never confuse the executable file on disk with the running instance in memory. One program can give rise to many processes (e.g., opening three terminal windows runs three \texttt{bash} processes). Conversely, a process may load several shared libraries, but it remains a single execution flow.
\end{attention}

\begin{savaistu}[Why Processes? Multiprogramming and Time-Sharing]
Early computers ran one program at a time (batch processing). The CPU often sat idle waiting for slow I/O devices. \cle{Multiprogramming} keeps several processes in memory simultaneously; when one process waits for I/O, the OS switches the CPU to another. \cle{Time-sharing} (multitasking) extends this by rapidly switching the CPU among processes, giving each user the illusion of a dedicated machine. Both techniques maximise CPU utilisation and throughput.
\end{savaistu}

\courssub{Process States and Transitions}

At any moment a process occupies exactly one of five states. The OS moves processes between states in response to specific events.

\begin{definition}[Five Process States]
\begin{itemize}
\item \cle{New}: The process is being created; the OS allocates a PCB and initial resources.
\item \cle{Ready}: The process resides in main memory, has all resources except the CPU, and waits only for the dispatcher to assign it a processor core.
\item \cle{Running}: The process's instructions are currently executing on a CPU core.
\item \cle{Blocked} (or \cle{Waiting}): The process cannot continue until an external event occurs (I/O completion, signal from another process, timer expiry, etc.). It voluntarily relinquishes the CPU.
\item \cle{Terminated} (or \cle{Exit}): The process has finished execution (normal exit or forced kill). The OS reclaims resources but keeps the PCB briefly so the parent can read the exit status.
\end{itemize}
\end{definition}

\begin{popfigure}
\centering
\begin{tikzpicture}[
    state/.style={circle, draw=popdark, thick, minimum size=2.2cm, fill=popblueL, text width=2.2cm, align=center, font=\small},
    arrow/.style={->, thick, >=stealth, popdark},
    label/.style={font=\footnotesize, text width=2.5cm, align=center}
]
% States
\node[state, align=center] (new) at (0,0){New\\Admission};
\node[state, align=center] (ready) at (4.5,0){Ready\\Dispatch};
\node[state, align=center] (running) at (9,0){Running\\Exit / I/O request / Preemption};
\node[state, align=center] (blocked) at (9,-4){Blocked\\I/O completion};
\node[state, align=center] (terminated) at (4.5,-4){Terminated\\Release resources};

% Transitions
\draw[arrow] (new) -- node[above,label] {admission} (ready);
\draw[arrow] (ready) -- node[above,label] {dispatch} (running);
\draw[arrow] (running) -- node[above,label] {exit} (terminated);
\draw[arrow] (running) .. controls (11,0) and (11,-2) .. node[right,label] {I/O request} (blocked);
\draw[arrow] (blocked) .. controls (11,-4) and (11,-2) .. node[right,label] {I/O completion} (ready);
\draw[arrow] (running) -- node[below,label] {preemption} (ready);
\end{tikzpicture}
\legende{Process state transition diagram. Arrows show the event that triggers the transition.}
\end{popfigure}

\begin{methode}[Identifying the State of a Process from a Scenario]
\begin{enumerate}
\item Locate the keyword: ``waiting for keyboard input'' $\to$ \textbf{Blocked}; ``selected by the scheduler'' $\to$ \textbf{Running}; ``in memory, ready to run'' $\to$ \textbf{Ready}.
\item Check whether the process holds the CPU. If yes $\to$ Running. If no, but it could run immediately if given CPU $\to$ Ready. If it needs an external event $\to$ Blocked.
\item Creation events (fork, exec) $\to$ New. Termination (exit, kill) $\to$ Terminated.
\end{enumerate}
\end{methode}

\begin{exemplebox}[State Identification]
A user opens a text editor. The editor reads a configuration file (disk I/O), then waits for a keystroke. Meanwhile the OS runs a background antivirus scan.
\begin{itemize}
\item Editor during disk read: \textbf{Blocked} (waiting for I/O completion).
\item Editor after file read, waiting for keystroke: \textbf{Blocked} (waiting for user input).
\item Antivirus scan while editor is blocked: \textbf{Running} (has CPU).
\item Editor after keystroke arrives, before scheduler picks it: \textbf{Ready}.
\end{itemize}
\end{exemplebox}

\courssub{The Process Control Block (PCB)}

Every process is represented in the OS by a \cle{Process Control Block} (PCB), sometimes called a \cle{Task Control Block}. The PCB is the OS's ``identity card'' for the process.

\begin{definition}[Process Control Block (PCB)]
A PCB is a kernel data structure that contains all information needed to suspend and later resume a process. Typical fields:
\begin{itemize}
\item \textbf{Process ID (PID)} — unique integer identifier.
\item \textbf{Process State} — Ready, Running, Blocked, etc.
\item \textbf{Program Counter (PC)} — address of the next instruction to execute.
\item \textbf{CPU Registers} — general-purpose registers, stack pointer, status flags.
\item \textbf{Memory Management Info} — page table pointer, segment table, memory limits.
\item \textbf{Open File Descriptors} — list of files, sockets, pipes the process uses.
\item \textbf{Scheduling Information} — priority, scheduling class, time quantum used.
\item \textbf{Accounting Information} — CPU time used, real time elapsed, user/group IDs.
\item \textbf{I/O Status Information} — devices allocated, pending I/O requests.
\end{itemize}
\end{definition}

\begin{popfigure}
\centering
\begin{tikzpicture}[
    pcb/.style={rectangle, draw=popdark, thick, minimum width=5cm, minimum height=0.7cm, fill=popblueL, font=\small},
    field/.style={rectangle, draw=popdark, minimum width=5cm, minimum height=0.55cm, font=\small},
    arrow/.style={->, thick, >=stealth, poporange}
]
% PCB 1 (Process A)
\node[pcb] (pcbA) at (0,2) {\textbf{PCB of Process A (PID 1024)}};
\foreach \i/\txt in {
    1/{State: Running},
    2/{PC: 0x400A3C},
    3/{Registers: RAX=12, RBX=0, ...},
    4/{Page Table: 0x7FF000},
    5/{Open Files: fd0=stdin, fd1=stdout},
    6/{Priority: 10},
    7/{CPU Time: 0.042 s}
} {
    \node[field, anchor=north] at ([yshift=-(\i-1)*0.6cm]pcbA.south) (fA\i) {\txt};
}
% PCB 2 (Process B)
\node[pcb] (pcbB) at (8,2) {\textbf{PCB of Process B (PID 1025)}};
\foreach \i/\txt in {
    1/{State: Ready},
    2/{PC: 0x800120},
    3/{Registers: RAX=0, RBX=5, ...},
    4/{Page Table: 0x810000},
    5/{Open Files: fd0=stdin, fd1=log.txt},
    6/{Priority: 8},
    7/{CPU Time: 0.015 s}
} {
    \node[field, anchor=north] at ([yshift=-(\i-1)*0.6cm]pcbB.south) (fB\i) {\txt};
}
% Context switch arrow
\draw[arrow] (4, -2.5) -- node[above, font=\small] {Context Switch} (4, -3.5);
\node[align=center, font=\footnotesize] at (4, -4.5) {OS saves A's registers into PCB A,\\loads B's registers from PCB B};
\end{tikzpicture}
\legende{Two PCBs during a context switch from Process A (Running) to Process B (Ready).}
\end{popfigure}

\courssub{Context Switching}

\begin{definition}[Context Switch]
A \cle{context switch} is the procedure of saving the context (CPU registers, PC, stack pointer, etc.) of the currently running process into its PCB and restoring the context of the next process to run from its PCB. It is triggered by interrupts (timer, I/O), system calls (blocking I/O, yield), or scheduler decisions.
\end{definition}

\begin{propriete}[Context Switch Overhead]
A context switch performs \textbf{no useful user work}; it is pure OS overhead. Typical costs on modern hardware: $1$–$5\ \mu\mathrm{s}$ (microseconds) for saving/restoring registers, flushing pipelines, invalidating TLB entries, and updating kernel data structures. Excessive switching (thrashing) wastes CPU cycles.
\end{propriete}

\begin{attention}[Context Switch Overhead]
Do not confuse context switching with the scheduler's decision-making. The scheduler \emph{chooses} the next process; the dispatcher \emph{performs} the context switch. Frequent preemption (e.g., a time quantum of $1\ \mathrm{ms}$) increases overhead; too long a quantum hurts responsiveness. The OS designer must balance both.
\end{attention}

\begin{popfigure}
\centering
\begin{tikzpicture}[
    proc/.style={rectangle, draw=popdark, thick, minimum height=0.8cm, fill=popgreenL, font=\small},
    cs/.style={rectangle, draw=poporange, thick, minimum height=0.8cm, fill=poporange!30, font=\small},
    timeaxis/.style={->, thick, >=stealth}
]
% Time axis
\draw[timeaxis] (0,0) -- (12,0) node[right] {Time ($\mu$s)};
\foreach \x/\lab in {0/0, 1.5/20, 3/40, 4.5/60, 6/80, 7.5/100, 9/120, 10.5/140}
    \draw (\x,0.1) -- (\x,-0.1) node[below] {\lab};

% Process A runs (0-20)
\node[proc, minimum width=1.5cm] at (0.75, -1.5) {Process A};
% Context switch (20-40)
\node[cs, minimum width=1.5cm] at (2.25, -1.5) {CS};
% Process B runs (40-80)
\node[proc, minimum width=3cm] at (4.5, -1.5) {Process B};
% Context switch (80-100)
\node[cs, minimum width=1.5cm] at (6.75, -1.5) {CS};
% Process A runs again (100-140)
\node[proc, minimum width=3cm] at (9.75, -1.5) {Process A};

% Labels
\node[align=center, font=\footnotesize] at (0.75, -2.5) {User code};
\node[align=center, font=\footnotesize] at (2.25, -2.5) {Save A / Load B};
\node[align=center, font=\footnotesize] at (4.5, -2.5) {User code};
\node[align=center, font=\footnotesize] at (6.75, -2.5) {Save B / Load A};
\node[align=center, font=\footnotesize] at (9.75, -2.5) {User code};
\end{tikzpicture}
\legende{Timeline showing CPU alternating between two processes with context-switch (CS) gaps. Time scale: each unit = 20 $\mu$s.}
\end{popfigure}

\courssub{Process Creation and Termination}

\begin{definition}[Process Creation]
In Unix-like systems, a new process is created by the \cle{fork()} system call. The parent process is cloned: the child receives a copy of the parent's address space (copy-on-write optimises this), file descriptors, and register state. The child gets a new PID. The return value of \texttt{fork()} distinguishes parent ($>0$, child's PID) from child ($=0$). The child often calls \cle{exec()} to replace its memory image with a new program.
\end{definition}

\begin{definition}[Process Termination]
A process terminates via \cle{exit()} (normal) or a signal (abnormal). The OS changes its state to \textbf{Terminated}, releases memory, closes files, and notifies the parent (via \texttt{SIGCHLD}). The parent can call \cle{wait()} to read the child's exit status.
\end{definition}

\begin{attention}[Orphan and Zombie Processes]
\begin{itemize}
\item \textbf{Orphan}: A child whose parent has terminated. The OS re-parents it to \texttt{init} (PID 1), which periodically calls \texttt{wait()}.
\item \textbf{Zombie}: A terminated child whose parent has not yet called \texttt{wait()}. The PCB remains in the process table, consuming a PID slot. Many zombies can exhaust the PID space.
\end{itemize}
\end{attention}

\courssub{Processes vs. Threads (GCE-Useful Extra)}

\begin{definition}[Thread]
A \cle{thread} (lightweight process) is a single sequential flow of control within a process. Threads share the process's code, data, heap, and open files, but each has its own stack, program counter, and register set. Creating and switching threads is faster than processes because no memory-map changes are needed.
\end{definition}

\begin{aretenir}[Process vs. Thread Summary]
\begin{center}
\begin{tabularx}{\linewidth}{|l|X|X|}\hline
\textbf{Aspect} & \textbf{Process} & \textbf{Thread} \\ \hline
Memory & Separate address spaces & Shared address space \\ \hline
Creation cost & High (copy page tables, etc.) & Low (allocate stack, TCB) \\ \hline
Context switch & Slow (TLB flush, new page table) & Fast (register swap only) \\ \hline
Communication & IPC (pipes, sockets, shared mem) & Direct via shared variables \\ \hline
Fault isolation & Crash affects only that process & Crash kills whole process \\ \hline
\end{tabularx}
\end{center}
\end{aretenir}

\courssub{The Process Table}

The OS maintains a \cle{process table} (or task array), an array or linked list of pointers to all PCBs in the system. The index or pointer serves as the PID. The scheduler scans the process table to find Ready processes; the dispatcher uses the PCB pointer to load context. In Linux, the process table is the \texttt{task\_struct} array; each entry is a \texttt{task\_struct} (the PCB).

\begin{exoresolu}[Process Table Walkthrough]
\textbf{Statement:} A system has three processes: PID 1 (\texttt{init}), PID 1024 (text editor), PID 1025 (compiler). The editor is blocked on keyboard input; the compiler is running; \texttt{init} is ready. Show the process table entries and explain what happens when the user presses a key.

\textbf{Solution:}
\begin{center}
\begin{tabularx}{\linewidth}{|c|X|X|X|}\hline
\textbf{PID} & \textbf{State} & \textbf{PC} & \textbf{Notes} \\ \hline
1 & Ready & 0x8000 & Parent of all orphans \\ \hline
1024 & Blocked & 0x400A3C & Waiting for keyboard interrupt \\ \hline
1025 & Running & 0x800120 & Currently on CPU \\ \hline
\end{tabularx}
\end{center}
When the key is pressed:
\begin{enumerate}
\item Keyboard controller raises interrupt $\to$ CPU vectors to keyboard ISR.
\item ISR places scan code in kernel buffer, signals the editor's wait queue.
\item OS changes editor's state from \textbf{Blocked} to \textbf{Ready} (updates PCB).
\item Scheduler may preempt compiler if editor has higher priority; otherwise editor waits in Ready queue.
\item Next context switch loads editor's PCB.
\end{enumerate}
\end{exoresolu}

\courssub{Worked Exercise: Identifying States and Transitions}

\begin{exercice}[State Transitions in a Print Job]
A user prints a large PDF. The print spooler process reads the file from disk, sends data to the printer via USB, and waits for the printer to signal ``page done'' before sending the next page. The CPU runs a video player while the spooler waits. Identify the state of the spooler at each stage and the transitions.
\end{exercice}

\begin{corrige}[Exercise 1]
\begin{enumerate}
\item Spooler starts (fork+exec): \textbf{New} $\xrightarrow{\text{admission}}$ \textbf{Ready}.
\item Scheduler dispatches spooler: \textbf{Ready} $\xrightarrow{\text{dispatch}}$ \textbf{Running}.
\item Spooler issues disk read: \textbf{Running} $\xrightarrow{\text{I/O request}}$ \textbf{Blocked}.
\item Disk interrupt, data ready: \textbf{Blocked} $\xrightarrow{\text{I/O completion}}$ \textbf{Ready}.
\item Spooler dispatched again: \textbf{Ready} $\xrightarrow{\text{dispatch}}$ \textbf{Running}.
\item Spooler sends page to printer, waits for ``page done'': \textbf{Running} $\xrightarrow{\text{I/O request}}$ \textbf{Blocked}.
\item Video player runs during wait: \textbf{Running} (video player).
\item Printer interrupt ``page done'': \textbf{Blocked} $\xrightarrow{\text{I/O completion}}$ \textbf{Ready} (spooler).
\item ... cycle repeats until last page.
\item Spooler calls exit: \textbf{Running} $\xrightarrow{\text{exit}}$ \textbf{Terminated}.
\end{enumerate}
\end{corrige}

\courssub{Key Points to Retain}

\begin{aretenir}[Lesson 44 Essentials]
\begin{itemize}
\item A \textbf{process} is a program in execution; a \textbf{program} is a static file.
\item Five states: \textbf{New, Ready, Running, Blocked, Terminated}. Know the six transitions: admission, dispatch, exit, I/O request, I/O completion, preemption.
\item The \textbf{PCB} holds everything needed to stop and restart a process (PID, state, PC, registers, memory info, open files, scheduling info).
\item \textbf{Context switch} = save current PCB + load next PCB. Pure overhead; keep it fast and infrequent enough.
\item \textbf{fork()} creates a child (copy of parent); \textbf{exec()} loads a new program. \textbf{exit()} terminates; \textbf{wait()} reaps zombies.
\item \textbf{Threads} share memory, lighter than processes; useful for parallelism within one application.
\item The \textbf{process table} is the OS's directory of all PCBs.
\end{itemize}
\end{aretenir}

\coursec{Section 2: Process Scheduling and Non-preemptive Algorithms (Lessons 45–46)}

In the previous section we saw that a \cle{process} is a program in execution, characterised by its \cle{Process Control Block (PCB)} and moving through states such as \emph{New}, \emph{Ready}, \emph{Running}, \emph{Waiting} and \emph{Terminated}. We now address a central question: \textbf{when several processes are in the Ready state, which one should be given the CPU next?} This decision is the job of the \cle{CPU scheduler}, and the algorithms it uses determine the overall performance, responsiveness and fairness of the system.

\courssub{Goals of CPU Scheduling}

The scheduler aims to optimise several, often competing, criteria:

\begin{itemize}
  \item \textbf{CPU utilisation} — keep the CPU as busy as possible (ideally 100\%).
  \item \textbf{Throughput} — number of processes completed per unit time.
  \item \textbf{Turnaround time} — time from submission to completion (minimise).
  \item \textbf{Waiting time} — time spent in the ready queue (minimise).
  \item \textbf{Response time} — time from submission to first response (critical for interactive systems).
  \item \textbf{Fairness} — no process should starve indefinitely.
\end{itemize}

In batch systems, turnaround and throughput dominate; in interactive systems, response time is paramount. A good scheduler balances these metrics.

\courssub{Key Time Metrics — Definitions and Formulas}

Before studying algorithms, we must precisely define the time quantities used to evaluate them. Consider a process $P_i$ with \cle{arrival time} $AT_i$ (when it enters the ready queue) and \cle{burst time} $BT_i$ (total CPU time it needs).

\begin{definition}[Completion, Turnaround and Waiting Times]
For a process $P_i$:
\begin{itemize}
  \item \cle{Completion Time} $CT_i$ — the clock time at which $P_i$ finishes execution.
  \item \cle{Turnaround Time} $TAT_i = CT_i - AT_i$ — total time in the system.
  \item \cle{Waiting Time} $WT_i = TAT_i - BT_i$ — time spent waiting in the ready queue.
\end{itemize}
\end{definition}

\begin{attention}[Common Mistake: Wrong Waiting Time Formula]
Many candidates write $WT = CT - AT$. This is \textbf{incorrect} — that formula gives \emph{turnaround time}. Waiting time excludes the actual CPU execution: $WT = TAT - BT = (CT - AT) - BT$. Always compute $TAT$ first, then subtract $BT$.
\end{attention}

\courssub{The Scheduler, the Dispatcher and the Ready Queue}

\begin{itemize}
  \item \cle{Ready queue} — a data structure (often a linked list or priority queue) holding PCBs of processes in the Ready state.
  \item \cle{Short-term scheduler (CPU scheduler)} — selects a process from the ready queue and allocates the CPU. Invoked frequently (milliseconds), so must be fast.
  \item \cle{Dispatcher} — performs the \cle{context switch}: saves the state of the outgoing process, loads the state of the incoming process, and jumps to its program counter. Dispatcher latency adds to scheduling overhead.
  \item \cle{Long-term scheduler (job scheduler)} — decides which jobs are admitted to the ready queue from the job pool (controls degree of multiprogramming). Invoked rarely (seconds/minutes).
  \item \cle{Medium-term scheduler} — swaps processes out to disk (suspended) and back in to manage memory pressure or improve process mix.
\end{itemize}

\courssub{CPU-bound vs I/O-bound Processes and the CPU–I/O Burst Cycle}

Processes alternate between \cle{CPU bursts} (computation) and \cle{I/O bursts} (waiting for devices).

\begin{itemize}
  \item \cle{CPU-bound process} — long CPU bursts, infrequent I/O (e.g., scientific computation, compilation).
  \item \cle{I/O-bound process} — short CPU bursts, frequent I/O (e.g., text editor, web browser).
\end{itemize}

A good scheduler gives priority to I/O-bound processes: they release the CPU quickly, keeping I/O devices busy and improving response time for interactive users.

\begin{popfigure}
\begin{tikzpicture}[scale=0.8, every node/.style={font=\small}]
% CPU-I/O burst cycle diagram
\draw[thick, ->] (0,0) -- (12,0) node[right] {Time};
\draw[thick] (0,-0.3) -- (0,0.3) node[above] {0};
% Process P1: CPU-I/O-CPU-I/O-CPU
\draw[fill=popblue!30, draw=popblue, thick] (0.5,0.5) rectangle (2.5,1.2) node[midway] {CPU};
\draw[fill=poporange!30, draw=poporange, thick] (2.5,0.5) rectangle (4.5,1.2) node[midway] {I/O};
\draw[fill=popblue!30, draw=popblue, thick] (4.5,0.5) rectangle (6.0,1.2) node[midway] {CPU};
\draw[fill=poporange!30, draw=poporange, thick] (6.0,0.5) rectangle (7.5,1.2) node[midway] {I/O};
\draw[fill=popblue!30, draw=popblue, thick] (7.5,0.5) rectangle (9.0,1.2) node[midway] {CPU};
\node at (-0.8,0.85) {P1:};
% Process P2: short CPU bursts
\draw[fill=popblue!30, draw=popblue, thick] (0.5,-1.2) rectangle (1.2,-0.5) node[midway] {CPU};
\draw[fill=poporange!30, draw=poporange, thick] (1.2,-1.2) rectangle (2.5,-0.5) node[midway] {I/O};
\draw[fill=popblue!30, draw=popblue, thick] (2.5,-1.2) rectangle (3.0,-0.5) node[midway] {CPU};
\draw[fill=poporange!30, draw=poporange, thick] (3.0,-1.2) rectangle (4.5,-0.5) node[midway] {I/O};
\draw[fill=popblue!30, draw=popblue, thick] (4.5,-1.2) rectangle (5.0,-0.5) node[midway] {CPU};
\node at (-0.8,-0.85) {P2:};
% Legend
\node[draw, fill=popblue!30, minimum width=1cm, minimum height=0.5cm] at (10.5,1.0) {};
\node at (11.5,1.0) {CPU burst};
\node[draw, fill=poporange!30, minimum width=1cm, minimum height=0.5cm] at (10.5,0.3) {};
\node at (11.5,0.3) {I/O burst};
\end{tikzpicture}
\legende{CPU–I/O burst cycle: CPU-bound process (P1) has long CPU bursts; I/O-bound process (P2) has short CPU bursts and frequent I/O.}
\end{popfigure}

\courssub{Non-preemptive Scheduling — Definition}

\begin{definition}[Non-preemptive Scheduling]
In \cle{non-preemptive scheduling}, once a process is allocated the CPU, it keeps the CPU until it \textbf{terminates} or \textbf{blocks for I/O} (enters the Waiting state). The scheduler cannot forcibly take the CPU away.
\end{definition}

This simplifies the dispatcher (no forced context switches) but can lead to long waiting times for short processes if a long process is running — the \cle{convoy effect}.

\courssub{First-Come First-Served (FCFS) Scheduling}

\courssub{Rule and Intuition}

FCFS is the simplest non-preemptive algorithm: the process that requests the CPU \textbf{first} gets it \textbf{first}. The ready queue is a FIFO queue. When the CPU becomes free, the process at the head of the queue is dispatched.

\courssub{Worked Example — FCFS Gantt Chart and Metrics}

Consider four processes with the following arrival and burst times:

\begin{center}
\begin{tabularx}{\linewidth}{|c|X|X|}\hline
Process & Arrival Time $AT$ & Burst Time $BT$ \\ \hline
P1 & 0 & 8 \\
P2 & 1 & 4 \\
P3 & 2 & 9 \\
P4 & 3 & 5 \\ \hline
\end{tabularx}
\end{center}

\begin{methode}[Step-by-step: Solving Any Scheduling Exercise]
\begin{enumerate}
  \item \textbf{Draw the Gantt chart} — simulate the algorithm on a timeline, respecting arrival times.
  \item \textbf{Read completion times $CT_i$} from the Gantt chart (end of each process block).
  \item \textbf{Compute Turnaround Time} $TAT_i = CT_i - AT_i$.
  \item \textbf{Compute Waiting Time} $WT_i = TAT_i - BT_i$.
  \item \textbf{Calculate averages} $\overline{TAT}$ and $\overline{WT}$.
\end{enumerate}
\end{methode}

\begin{exoresolu}[FCFS Scheduling for the Four Processes]
\textbf{Step 1: Gantt Chart}

Processes arrive at 0, 1, 2, 3. FCFS serves in arrival order: P1, P2, P3, P4.

\begin{popfigure}
\begin{tikzpicture}[scale=0.7, every node/.style={font=\small}]
% Time axis
\draw[thick, ->] (0,0) -- (28,0);
\foreach \x in {0,2,4,6,8,10,12,14,16,18,20,22,24,26} {
  \draw (\x,0.1) -- (\x,-0.1) node[below] {\x};
}
% P1: 0-8
\draw[fill=popblue!40, draw=popblue, thick] (0,0.5) rectangle (8,1.5) node[midway] {P1 (8)};
% P2: 8-12
\draw[fill=popgreen!40, draw=popgreen, thick] (8,0.5) rectangle (12,1.5) node[midway] {P2 (4)};
% P3: 12-21
\draw[fill=poporange!40, draw=poporange, thick] (12,0.5) rectangle (21,1.5) node[midway] {P3 (9)};
% P4: 21-26
\draw[fill=poppurple!40, draw=poppurple, thick] (21,0.5) rectangle (26,1.5) node[midway] {P4 (5)};
% Arrival markers
\foreach \t/\lab in {0/P1, 1/P2, 2/P3, 3/P4} {
  \draw[dashed, thin] (\t,0) -- (\t,-1.2);
  \node[below] at (\t,-1.2) {\lab\ arrives};
}
\node at (-1.5,1) {FCFS Gantt Chart};
\end{tikzpicture}
\legende{FCFS Gantt chart: processes run in arrival order P1 $\to$ P2 $\to$ P3 $\to$ P4.}
\end{popfigure}

\textbf{Step 2: Completion Times}
$CT_1 = 8$, $CT_2 = 12$, $CT_3 = 21$, $CT_4 = 26$.

\textbf{Step 3: Turnaround Times}
$TAT_1 = 8 - 0 = 8$
$TAT_2 = 12 - 1 = 11$
$TAT_3 = 21 - 2 = 19$
$TAT_4 = 26 - 3 = 23$

\textbf{Step 4: Waiting Times}
$WT_1 = 8 - 8 = 0$
$WT_2 = 11 - 4 = 7$
$WT_3 = 19 - 9 = 10$
$WT_4 = 23 - 5 = 18$

\textbf{Step 5: Averages}
$\overline{TAT} = \frac{8+11+19+23}{4} = \frac{61}{4} = 15.25$
$\overline{WT} = \frac{0+7+10+18}{4} = \frac{35}{4} = 8.75$

\begin{center}
\begin{tabularx}{\linewidth}{|c|X|X|X|X|X|}\hline
Process & $AT$ & $BT$ & $CT$ & $TAT$ & $WT$ \\ \hline
P1 & 0 & 8 & 8 & 8 & 0 \\
P2 & 1 & 4 & 12 & 11 & 7 \\
P3 & 2 & 9 & 21 & 19 & 10 \\
P4 & 3 & 5 & 26 & 23 & 18 \\ \hline
\textbf{Average} & & & & \textbf{15.25} & \textbf{8.75} \\ \hline
\end{tabularx}
\end{center}
\end{exoresolu}

\courssub{The Convoy Effect}

\begin{attention}[Convoy Effect in FCFS]
If a long CPU-bound process (e.g., P1 with $BT=8$) arrives first, all subsequent short processes (P2, P3, P4) must wait for it to finish, even though they could have completed quickly. This \cle{convoy effect} inflates average waiting time and reduces CPU/I/O overlap. FCFS is therefore poor for interactive systems.
\end{attention}

\courssub{Shortest Job First (SJF) Scheduling}

\courssub{Rule and Intuition}

SJF selects the process with the \textbf{smallest next CPU burst time} from the ready queue. It is non-preemptive: once started, a process runs to completion. SJF is \textbf{provably optimal} for minimising average waiting time for a given set of processes.

\begin{propriete}[Optimality of SJF (Non-preemptive)]
For any set of processes with known burst times, non-preemptive SJF gives the \textbf{minimum possible average waiting time}.
\end{propriete}

\emph{Proof sketch (examinable):} Consider any schedule not in SJF order. There exists a pair of adjacent processes where a longer job precedes a shorter one. Swapping them reduces the waiting time of the shorter job by the burst time of the longer job, while increasing the waiting time of the longer job by the burst time of the shorter job. Since the shorter burst is smaller, the net change in total waiting time is negative. By repeated swaps, any schedule can be transformed into SJF order while strictly decreasing average waiting time. Hence SJF is optimal.

\courssub{Worked Example — SJF on the Same Four Processes}

\begin{exoresolu}[SJF Scheduling for the Four Processes]
\textbf{Step 1: Gantt Chart}

At time 0, only P1 is available $\to$ P1 runs 0–8.
At time 8, ready queue: P2 ($BT=4$), P3 ($BT=9$), P4 ($BT=5$). Shortest is P2 $\to$ runs 8–12.
At time 12, ready: P3 ($9$), P4 ($5$). Shortest is P4 $\to$ runs 12–17.
At time 17, only P3 remains $\to$ runs 17–26.

\begin{popfigure}
\begin{tikzpicture}[scale=0.7, every node/.style={font=\small}]
\draw[thick, ->] (0,0) -- (28,0);
\foreach \x in {0,2,4,6,8,10,12,14,16,18,20,22,24,26} {
  \draw (\x,0.1) -- (\x,-0.1) node[below] {\x};
}
% P1: 0-8
\draw[fill=popblue!40, draw=popblue, thick] (0,0.5) rectangle (8,1.5) node[midway] {P1 (8)};
% P2: 8-12
\draw[fill=popgreen!40, draw=popgreen, thick] (8,0.5) rectangle (12,1.5) node[midway] {P2 (4)};
% P4: 12-17
\draw[fill=poppurple!40, draw=poppurple, thick] (12,0.5) rectangle (17,1.5) node[midway] {P4 (5)};
% P3: 17-26
\draw[fill=poporange!40, draw=poporange, thick] (17,0.5) rectangle (26,1.5) node[midway] {P3 (9)};
% Arrival markers
\foreach \t/\lab in {0/P1, 1/P2, 2/P3, 3/P4} {
  \draw[dashed, thin] (\t,0) -- (\t,-1.2);
  \node[below] at (\t,-1.2) {\lab\ arrives};
}
\node at (-1.5,1) {SJF Gantt Chart};
\end{tikzpicture}
\legende{SJF Gantt chart: at each decision point, the shortest available burst is chosen. Order: P1 $\to$ P2 $\to$ P4 $\to$ P3.}
\end{popfigure}

\textbf{Step 2: Completion Times}
$CT_1 = 8$, $CT_2 = 12$, $CT_4 = 17$, $CT_3 = 26$.

\textbf{Step 3: Turnaround Times}
$TAT_1 = 8 - 0 = 8$
$TAT_2 = 12 - 1 = 11$
$TAT_4 = 17 - 3 = 14$
$TAT_3 = 26 - 2 = 24$

\textbf{Step 4: Waiting Times}
$WT_1 = 8 - 8 = 0$
$WT_2 = 11 - 4 = 7$
$WT_4 = 14 - 5 = 9$
$WT_3 = 24 - 9 = 15$

\textbf{Step 5: Averages}
$\overline{TAT} = \frac{8+11+14+24}{4} = \frac{57}{4} = 14.25$
$\overline{WT} = \frac{0+7+9+15}{4} = \frac{31}{4} = 7.75$

\begin{center}
\begin{tabularx}{\linewidth}{|c|X|X|X|X|X|}\hline
Process & $AT$ & $BT$ & $CT$ & $TAT$ & $WT$ \\ \hline
P1 & 0 & 8 & 8 & 8 & 0 \\
P2 & 1 & 4 & 12 & 11 & 7 \\
P4 & 3 & 5 & 17 & 14 & 9 \\
P3 & 2 & 9 & 26 & 24 & 15 \\ \hline
\textbf{Average} & & & & \textbf{14.25} & \textbf{7.75} \\ \hline
\end{tabularx}
\end{center}

\textbf{Comparison:} SJF reduces average waiting time from \textbf{8.75} (FCFS) to \textbf{7.75}, and average turnaround from 15.25 to 14.25.
\end{exoresolu}

\courssub{Problems with SJF}

\begin{itemize}
  \item \textbf{Predicting burst times:} The OS does not know $BT$ in advance. Common heuristic: \cle{exponential averaging} $\tau_{n+1} = \alpha \cdot t_n + (1-\alpha)\tau_n$, where $t_n$ is the last actual burst, $\tau_n$ the predicted burst, $\alpha \in [0,1]$ (typically 0.5).
  \item \textbf{Starvation of long jobs:} A steady stream of short jobs can indefinitely postpone a long job. Remedy: \cle{ageing} (increase priority of waiting processes over time).
\end{itemize}

\courssub{Priority Scheduling (Non-preemptive Version)}

\courssub{Rule}

Each process has a \cle{priority} (integer; lower number = higher priority or vice versa — convention must be stated). The scheduler picks the highest-priority process from the ready queue. Non-preemptive: runs to completion or block.

\courssub{Worked Example — Priority Scheduling}

Same four processes, with priorities (lower = higher priority):

\begin{center}
\begin{tabularx}{\linewidth}{|c|X|X|X|}\hline
Process & $AT$ & $BT$ & Priority \\ \hline
P1 & 0 & 8 & 3 \\
P2 & 1 & 4 & 1 \\
P3 & 2 & 9 & 4 \\
P4 & 3 & 5 & 2 \\ \hline
\end{tabularx}
\end{center}

\begin{exoresolu}[Non-preemptive Priority Scheduling]
At time 0: only P1 (pri 3) $\to$ runs 0–8.
At time 8: ready: P2 (pri 1), P3 (pri 4), P4 (pri 2). Highest priority = P2 $\to$ runs 8–12.
At time 12: ready: P3 (pri 4), P4 (pri 2). Highest = P4 $\to$ runs 12–17.
At time 17: P3 runs 17–26.

Order: P1 $\to$ P2 $\to$ P4 $\to$ P3 (same as SJF in this case, but for different reason).

Metrics identical to SJF example above: $\overline{WT} = 7.75$, $\overline{TAT} = 14.25$.
\end{exoresolu}

\courssub{Starvation and Ageing}

\begin{attention}[Starvation in Priority Scheduling]
Low-priority processes may never run if high-priority processes keep arriving. \cle{Ageing} gradually increases the priority of waiting processes (e.g., increment priority by 1 every fixed time interval). This guarantees eventual execution.
\end{attention}

\courssub{Comparative Visualisation: FCFS vs SJF Average Waiting Time}

\begin{popfigure}
\begin{tikzpicture}
\begin{axis}[
  ybar,
  bar width=1.2cm,
  width=12cm,
  height=7cm,
  ylabel={Average Waiting Time},
  xlabel={Scheduling Algorithm},
  xtick={1,2},
  xticklabels={FCFS, SJF},
  ymin=0,
  ymax=10,
  ymajorgrids=true,
  grid style=dashed,
  nodes near coords,
  nodes near coords align={vertical},
  every node near coord/.append style={font=\small, anchor=south},
  legend style={at={(0.5,-0.15)}, anchor=north, legend columns=-1},
  title={Comparison of Average Waiting Time},
  title style={font=\normalsize\bfseries}
]
\addplot[fill=popblue!60, draw=popblue] coordinates {(1,8.75)};
\addplot[fill=popgreen!60, draw=popgreen] coordinates {(2,7.75)};
\legend{FCFS, SJF}
\end{axis}
\end{tikzpicture}
\legende{Bar chart comparing average waiting time: SJF (7.75) outperforms FCFS (8.75) for the given process set.}
\end{popfigure}

\courssub{Summary of Non-preemptive Algorithms}

\begin{center}
\begin{tabularx}{\linewidth}{|l|X|X|X|}\hline
Algorithm & Decision Rule & Advantages & Disadvantages \\ \hline
FCFS & Arrival order & Simple, fair (no starvation) & Convoy effect, high avg waiting time \\
SJF & Shortest burst first & Optimal avg waiting time & Requires burst prediction; starvation of long jobs \\
Priority & Highest priority first & Flexible, matches business needs & Starvation of low priority; needs ageing \\ \hline
\end{tabularx}
\end{center}

\courssub{Key Points to Remember}

\begin{aretenir}[Key Points for Examination]
\begin{itemize}
  \item Turnaround Time $TAT = CT - AT$; Waiting Time $WT = TAT - BT = CT - AT - BT$.
  \item Non-preemptive: process keeps CPU until termination or I/O block.
  \item FCFS: FIFO queue; suffers from convoy effect.
  \item SJF: picks shortest burst; \textbf{provably minimises average waiting time} (proof by pairwise swap).
  \item SJF needs burst prediction (exponential averaging); can starve long jobs.
  \item Priority scheduling: highest priority first; starvation solved by ageing.
  \item Always draw Gantt chart first, then compute $CT$, $TAT$, $WT$, then averages.
\end{itemize}
\end{aretenir}

\courssub{Exercises}

\begin{exercice}[Practice: FCFS vs SJF]
Five processes arrive at times and have burst times as follows:
P1: $AT=0, BT=10$; P2: $AT=1, BT=1$; P3: $AT=2, BT=2$; P4: $AT=3, BT=1$; P5: $AT=4, BT=5$.
\begin{enumerate}
  \item Draw the FCFS Gantt chart and compute $\overline{WT}$ and $\overline{TAT}$.
  \item Draw the SJF Gantt chart and compute $\overline{WT}$ and $\overline{TAT}$.
  \item Explain the difference.
\end{enumerate}
\end{exercice}

\begin{exercice}[Priority with Ageing]
A system uses non-preemptive priority scheduling (lower number = higher priority). Processes:
P1: $AT=0, BT=8, Pri=3$; P2: $AT=1, BT=4, Pri=1$; P3: $AT=2, BT=9, Pri=4$; P4: $AT=3, BT=5, Pri=2$.
Ageing increases priority by 1 every 5 time units of waiting.
\begin{enumerate}
  \item Simulate the schedule up to time 20.
  \item Does ageing change the order compared to no ageing?
\end{enumerate}
\end{exercice}

\begin{corrige}[Exercise 1]
\textbf{FCFS:} Order P1(0–10), P2(10–11), P3(11–13), P4(13–14), P5(14–19).
$CT$: 10,11,13,14,19. $TAT$: 10,10,11,11,15. $WT$: 0,9,9,10,10. $\overline{WT}=7.6$, $\overline{TAT}=11.4$.

\textbf{SJF:} Time 0: P1 runs 0–10. Time 10: ready P2(1),P3(2),P4(1),P5(5). Tie P2/P4 (both 1) $\to$ FCFS tie-break: P2(10–11), P4(11–12), P3(12–14), P5(14–19).
$CT$: P1=10, P2=11, P4=12, P3=14, P5=19.
$TAT$: 10,10,9,12,15. $WT$: 0,9,8,10,10. $\overline{WT}=7.4$, $\overline{TAT}=11.2$.
SJF slightly better; convoy effect visible in FCFS (P1 blocks many short jobs).
\end{corrige}

\begin{corrige}[Exercise 2]
Without ageing: same as priority example above (P1,P2,P4,P3).
With ageing: at time 8, P3 waited 6 units $\to$ priority 4$\to$3; P4 waited 5 units $\to$ priority 2$\to$1. At time 8, ready: P2(pri1), P3(pri3), P4(pri1). Tie P2/P4 $\to$ FCFS: P2 runs. At time 12, P3 waited 10 units $\to$ priority 4$\to$2; P4 waited 9 units $\to$ priority still 1 (next increment at 10 units waited, i.e., time 13). Ready: P3(pri2), P4(pri1). Highest priority is P4 $\to$ P4 runs. Order unchanged here, but ageing can change order in other cases.
\end{corrige}

\begin{savaistu}[Real-world Scheduling in Linux]
Linux uses the \cle{Completely Fair Scheduler (CFS)}, a preemptive, priority-based scheduler that models an \emph{ideal multi-tasking processor} using virtual runtime. It approximates SJF for interactive tasks via dynamic priority adjustments. The concepts you learn here (FCFS, SJF, priority, ageing) are the building blocks of modern schedulers like CFS, Windows' Multilevel Feedback Queue, and real-time schedulers (Rate Monotonic, EDF) used in embedded systems controlling traffic lights in Yaoundé or Douala.
\end{savaistu}

\coursec{Section 3: Preemptive Scheduling (Lesson 47)}

In the previous section we studied \emph{non-preemptive} scheduling: once a process is given the CPU, it keeps it until it terminates or voluntarily enters the waiting state (e.g., for I/O). That policy is simple to implement, but it has a serious weakness: a single long process can monopolise the processor while many short, interactive processes wait in the ready queue. Imagine a cybercafé in Yaoundé where one customer is printing a 200-page document: if the attendant refuses to serve anyone else until the printing is over, every other customer leaves frustrated. A wiser attendant \emph{pauses} the big job, serves the quick requests, then resumes the big job. This is exactly the idea of \cle{preemptive scheduling}.

\courssub{The Principle of Preemption}

\begin{definition}[Preemptive Scheduling]
A scheduling policy is \cle{preemptive} if the operating system may \textbf{interrupt a running process} before it has finished its CPU burst, move it back to the ready queue, and allocate the CPU to another process. The interruption is usually triggered by a \cle{timer interrupt} (the hardware clock signals the kernel at regular intervals, e.g., every 10--100 ms) or by the arrival of a higher-priority process.
\end{definition}

When the CPU is taken away from a process, the operating system performs a \cle{context switch}: it saves the state of the interrupted process (program counter, CPU registers, memory management information) in its Process Control Block (PCB), then loads the state of the newly selected process. Context switches make preemption possible, but they are \emph{pure overhead}: during a context switch the CPU does no useful work for any user process.

\begin{propriete}[Advantages and Costs of Preemption]
\begin{itemize}
\item \textbf{Advantages:} good \cle{responsiveness} (interactive programs react quickly); \textbf{no process can monopolise the CPU}; essential for multi-user, time-sharing, and real-time systems.
\item \textbf{Costs:} \cle{context-switch overhead} at every preemption (saving/restoring registers, flushing pipelines, cache pollution); shared data may be left in an inconsistent state, so preemption creates the need for \cle{synchronisation} mechanisms (studied in Section 4).
\end{itemize}
\end{propriete}

\courssub{Shortest Remaining Time First (SRTF)}

\cle{SRTF} is simply the \emph{preemptive version} of Shortest Job First (SJF). The rule is:

\begin{quote}
Whenever a new process arrives (or a process finishes), compare the \textbf{remaining burst time} of the currently running process with the burst times of all processes in the ready queue. The process with the \textbf{shortest remaining time} gets the CPU --- even if this means interrupting the running process.
\end{quote}

\begin{attention}[The Classic SRTF Mistake]
The most frequent examination error is to choose the shortest job \emph{once} at the start and then let it run to completion. That is SJF, not SRTF! In SRTF you must \textbf{re-evaluate the shortest remaining time at every new arrival}. Draw a vertical line on your time axis at each arrival time and ask: ``Does the newcomer beat the current process's remaining time?''
\end{attention}

\begin{exemplebox}[SRTF on Four Processes]
Consider the following processes (arrival time, burst time in ms):
\begin{center}
\begin{tabular}{|c|c|c|}
\hline
\rowcolor{popblueL} Process & Arrival & Burst \\
\hline
$P_1$ & 0 & 8 \\
$P_2$ & 1 & 4 \\
$P_3$ & 2 & 9 \\
$P_4$ & 3 & 5 \\
\hline
\end{tabular}
\end{center}
Trace of the scheduler (decision points at every arrival and completion):
\begin{itemize}
\item $t=0$: only $P_1$ is present; it runs.
\item $t=1$: $P_2$ arrives (burst 4). $P_1$ has $8-1=7$ ms remaining. Since $4 < 7$, $P_1$ is \textbf{preempted} and $P_2$ runs.
\item $t=2$: $P_3$ arrives (9). $P_2$ has 3 remaining; $3 < 9$, so $P_2$ continues.
\item $t=3$: $P_4$ arrives (5). $P_2$ has 2 remaining; $2 < 5$, so $P_2$ continues.
\item $t=5$: $P_2$ finishes. Ready queue: $P_4$ (5), $P_1$ (7), $P_3$ (9). $P_4$ runs.
\item $t=10$: $P_4$ finishes. $P_1$ (7) beats $P_3$ (9); $P_1$ runs to completion at $t=17$, then $P_3$ runs from 17 to 26.
\end{itemize}
Turnaround times (completion $-$ arrival): $P_1: 17$, $P_2: 4$, $P_3: 24$, $P_4: 7$. Average turnaround $= \dfrac{17+4+24+7}{4} = \SI{13}{ms}$.
\end{exemplebox}

\begin{popfigure}
\begin{center}
\begin{tikzpicture}[x=0.42cm,y=0.8cm]
\draw[->] (0,0) -- (27,0) node[right] {\small $t$};
\foreach \t in {0,1,2,3,5,10,17,26} \draw (\t,0.08) -- (\t,-0.08) node[below] {\scriptsize \t};
\fill[popblue]   (0,0.25)    rectangle (1,0.85);
\fill[poporange] (1,0.25)    rectangle (5,0.85);
\fill[popgreen]  (5,0.25)    rectangle (10,0.85);
\fill[popblue]   (10,0.25)   rectangle (17,0.85);
\fill[poppurple] (17,0.25)   rectangle (26,0.85);
\draw (0,0.25) rectangle (26,0.85);
\node at (0.5,0.55) {\scriptsize $P_1$};
\node at (3,0.55) {\scriptsize $P_2$};
\node at (7.5,0.55) {\scriptsize $P_4$};
\node at (13.5,0.55) {\scriptsize $P_1$};
\node at (21.5,0.55) {\scriptsize $P_3$};
\foreach \t in {1,5,10,17} \draw[popdark, dashed] (\t,0.25) -- (\t,1.35);
\node[popdark] at (1,1.5) {\scriptsize preempt};
\node[popdark] at (5,1.5) {\scriptsize $P_2$ done};
\end{tikzpicture}
\end{center}
\legende{Gantt chart for SRTF. $P_1$ is preempted at $t=1$ when the shorter job $P_2$ arrives.}
\end{popfigure}

\courssub{Round Robin (RR)}

SRTF is optimal for average waiting time (provably minimal), but it requires knowing burst times in advance --- which is rarely possible --- and it can starve long processes. \cle{Round Robin} takes a completely different, beautifully fair approach: \emph{every process gets a small, equal slice of the CPU, in turn}.

\begin{definition}[Time Quantum]
In Round Robin scheduling, the \cle{time quantum} (or \emph{time slice}) $q$ is the maximum continuous time a process may hold the CPU before being preempted. The ready queue is organised as a \textbf{circular FIFO queue}: the scheduler takes the process at the front, lets it run for at most $q$ time units, and if it is not finished, puts it back \textbf{at the rear} of the queue.
\end{definition}

\begin{methode}[Applying Round Robin --- Step by Step]
\begin{enumerate}
\item \textbf{Initialisation:} Place all processes that have arrived at time 0 in the ready queue (FCFS order). Set current time $t=0$.
\item \textbf{Event Loop:} While the ready queue is not empty:
  \begin{enumerate}
  \item \textbf{Handle arrivals at time $t$:} Add any process arriving exactly at time $t$ to the \textbf{rear} of the ready queue.
  \item \textbf{Handle preemption at time $t$ (if $t>0$):} If the process that just finished its quantum (or blocked) is not terminated, move it to the \textbf{rear} of the ready queue.
  \item \textbf{Dispatch:} Remove the process at the \textbf{front} of the queue. Let it run for $\Delta t = \min(q,\ \text{remaining burst})$. Advance time: $t \leftarrow t + \Delta t$.
  \item \textbf{Record:} Add the interval to the Gantt chart. If the process finishes, record its completion time; otherwise it will be re-queued at the next iteration (step 2b).
  \end{enumerate}
\item \textbf{Metrics:} From the Gantt chart, compute for each process:
  \begin{itemize}
  \item Turnaround Time $TAT = \text{Completion} - \text{Arrival}$
  \item Waiting Time $WT = TAT - \text{Burst}$
  \end{itemize}
\end{enumerate}
\emph{Convention used in this course: at a given instant $t$, arrivals are enqueued \textbf{before} the preempted process is re-queued. This ensures a newly arrived process gets a chance slightly sooner than the process that just used the CPU.}
\end{methode}

\begin{exemplebox}[Round Robin with $q = 2$ on the Same Processes]
Using $P_1(0,8)$, $P_2(1,4)$, $P_3(2,9)$, $P_4(3,5)$ with quantum $q=2$:
\begin{itemize}
\item $t=0$: Queue $[P_1]$. $P_1$ runs $0\!-\!2$ (rem 6).
\item $t=2$: Arrivals: $P_3$ (arr 2). Queue $[P_2, P_3]$. Preempt $P_1$ $\to$ Queue $[P_2, P_3, P_1]$. Dispatch $P_2$.
\item $t=4$: $P_2$ runs $2\!-\!4$ (rem 2). Arrival: none. Preempt $P_2$ $\to$ Queue $[P_3, P_1, P_2]$. Dispatch $P_3$.
\item $t=6$: $P_3$ runs $4\!-\!6$ (rem 7). Arrival: $P_4$ (arr 3, already in queue since $t=3$? No, $P_4$ arrives at 3. At $t=3$, $P_2$ was running. $P_4$ added to queue $\to [P_3, P_1, P_4]$? Wait. Let's trace queue at $t=3$: Running $P_2$. Queue $[P_3, P_1]$. $P_4$ arrives $\to [P_3, P_1, P_4]$. At $t=4$, Preempt $P_2 \to [P_3, P_1, P_4, P_2]$. Dispatch $P_3$.)
\item $t=6$: $P_3$ preempted (rem 7). Queue $[P_1, P_4, P_2, P_3]$. Dispatch $P_1$.
\item $t=8$: $P_1$ runs $6\!-\!8$ (rem 4). Preempt $P_1 \to [P_4, P_2, P_3, P_1]$. Dispatch $P_4$.
\item $t=10$: $P_4$ runs $8\!-\!10$ (rem 3). Preempt $P_4 \to [P_2, P_3, P_1, P_4]$. Dispatch $P_2$.
\item $t=12$: $P_2$ runs $10\!-\!12$ and \textbf{finishes}. Queue $[P_3, P_1, P_4]$. Dispatch $P_3$.
\item $t=14$: $P_3$ runs $12\!-\!14$ (rem 5). Preempt $P_3 \to [P_1, P_4, P_3]$. Dispatch $P_1$.
\item $t=16$: $P_1$ runs $14\!-\!16$ (rem 2). Preempt $P_1 \to [P_4, P_3, P_1]$. Dispatch $P_4$.
\item $t=18$: $P_4$ runs $16\!-\!18$ (rem 1). Preempt $P_4 \to [P_3, P_1, P_4]$. Dispatch $P_3$.
\item $t=20$: $P_3$ runs $18\!-\!20$ (rem 3). Preempt $P_3 \to [P_1, P_4, P_3]$. Dispatch $P_1$.
\item $t=22$: $P_1$ runs $20\!-\!22$ and \textbf{finishes}. Queue $[P_4, P_3]$. Dispatch $P_4$.
\item $t=23$: $P_4$ runs $22\!-\!23$ and \textbf{finishes}. Queue $[P_3]$. Dispatch $P_3$.
\item $t=26$: $P_3$ runs $23\!-\!26$ and \textbf{finishes}.
\end{itemize}
Turnarounds: $P_1: 22$, $P_2: 11$, $P_3: 24$, $P_4: 20$; average $= \dfrac{77}{4} = \SI{19.25}{ms}$. The average is worse than SRTF, but notice that \emph{every} process received CPU time within the first 10 ms --- excellent responsiveness.
\end{exemplebox}

\begin{popfigure}
\begin{center}
\begin{tikzpicture}[x=0.42cm,y=0.8cm]
\draw[->] (0,0) -- (27,0) node[right] {\small $t$};
\foreach \t in {0,2,4,6,8,10,12,14,16,18,20,22,23,26} \draw (\t,0.08) -- (\t,-0.08) node[below] {\scriptsize \t};
\fill[popblue]   (0,0.25)  rectangle (2,0.85);
\fill[poporange] (2,0.25)  rectangle (4,0.85);
\fill[poppurple] (4,0.25)  rectangle (6,0.85);
\fill[popblue]   (6,0.25)  rectangle (8,0.85);
\fill[popgreen]  (8,0.25)  rectangle (10,0.85);
\fill[poporange] (10,0.25) rectangle (12,0.85);
\fill[poppurple] (12,0.25) rectangle (14,0.85);
\fill[popblue]   (14,0.25) rectangle (16,0.85);
\fill[popgreen]  (16,0.25) rectangle (18,0.85);
\fill[poppurple] (18,0.25) rectangle (20,0.85);
\fill[popblue]   (20,0.25) rectangle (22,0.85);
\fill[popgreen]  (22,0.25) rectangle (23,0.85);
\fill[poppurple] (23,0.25) rectangle (26,0.85);
\draw (0,0.25) rectangle (26,0.85);
\node at (1,0.55) {\scriptsize $P_1$}; \node at (3,0.55) {\scriptsize $P_2$};
\node at (5,0.55) {\scriptsize $P_3$}; \node at (7,0.55) {\scriptsize $P_1$};
\node at (9,0.55) {\scriptsize $P_4$}; \node at (11,0.55) {\scriptsize $P_2$};
\node at (13,0.55) {\scriptsize $P_3$}; \node at (15,0.55) {\scriptsize $P_1$};
\node at (17,0.55) {\scriptsize $P_4$}; \node at (19,0.55) {\scriptsize $P_3$};
\node at (21,0.55) {\scriptsize $P_1$}; \node at (22.5,0.55) {\scriptsize $P_4$};
\node at (24.5,0.55) {\scriptsize $P_3$};
\end{tikzpicture}
\end{center}
\legende{Gantt chart for Round Robin with quantum $q=2$ on the same four processes.}
\end{popfigure}

\begin{propriete}[Bounded Response Time in RR]
With $n$ processes in the ready queue and quantum $q$, a process waits \textbf{at most $(n-1)\times q$} time units before receiving the CPU again. Round Robin therefore guarantees a bounded response time: no process can starve. (This bound follows because at most $n-1$ other processes can run for at most $q$ each before your turn comes again.) The proof idea is examinable as a short justification, not a formal proof.
\end{propriete}

\courssub{Choosing the Quantum}

The performance of Round Robin depends entirely on the choice of $q$:

\begin{itemize}
\item \textbf{Very large quantum} ($q \geq$ longest CPU burst): no process is ever preempted, so RR degenerates into \cle{FCFS}. Responsiveness is lost.
\item \textbf{Very small quantum}: the scheduler context-switches constantly; if a context switch costs about the same as the quantum, most CPU time is wasted on administration instead of useful work (thrashing).
\item \textbf{Reasonable quantum}: typically chosen so that the large majority of interactive bursts complete within one slice (real systems use values around $10$--$100$ ms, e.g., Linux default time slice is dynamic but often in this range).
\end{itemize}

\begin{attention}[Quantum Too Small]
A classic trap: candidates write ``smaller quantum = fairer = better''. Wrong. If the quantum is too small, \cle{context-switch overhead dominates} and the effective CPU throughput collapses. In calculations, if the question gives a context-switch cost (e.g., 1 ms), you must add it to the timeline at every preemption (i.e., the clock advances by $q + \text{switch\_cost}$ for a full slice).
\end{attention}

\courssub{Preemptive Priority Scheduling, Starvation and Ageing}

In \cle{preemptive priority scheduling}, each process has a priority number (integer); the CPU always goes to the ready process with the \textbf{highest priority} (lowest number or highest number, convention must be stated), and a newly arrived high-priority process \emph{immediately preempts} the running one. This is ideal for urgent tasks, but it creates a danger:

\begin{definition}[Starvation and Ageing]
\cle{Starvation} (or indefinite blocking) occurs when a low-priority process waits forever because higher-priority processes keep arriving. The standard cure is \cle{ageing}: the priority of a waiting process is \textbf{gradually increased} (or its priority number decreased) as time passes, so that every process eventually reaches the top of the queue and runs.
\end{definition}

\begin{attention}[Exam Tip: State Your Tie-Breaking Rule]
Whenever two processes have equal priority, equal remaining time, or arrive simultaneously, a tie-breaking rule is needed. \textbf{Always state it explicitly} in your answer --- the standard convention is FCFS (the process that arrived first goes first). An unstated assumption can cost you marks even if your Gantt chart is otherwise correct.
\end{attention}

\courssub{Multilevel Queue Scheduling (GCE-Useful Extra)}

Real systems rarely use a single algorithm for all processes. Two hybrid schemes are worth knowing for the ``integration activity'' and essay questions:

\begin{itemize}
\item \textbf{Multilevel queue:} the ready queue is split into several separate queues (e.g., \emph{system processes}, \emph{interactive}, \emph{batch}), each with its own algorithm (RR for interactive, FCFS for batch). A process stays in its assigned queue forever, and queues are served by fixed priority (e.g., system $>$ interactive $>$ batch). Simple but rigid.
\item \textbf{Multilevel feedback queue:} processes may \emph{move between} queues. A new process starts in the top queue with a small quantum; if it uses its whole quantum (CPU-bound), it is demoted to a lower queue with a larger quantum; if it waits a long time, it is promoted (ageing). This adapts automatically to each process's behaviour and is the basis of many real schedulers (e.g., classic Unix, Windows NT/XP).
\end{itemize}

\begin{savaistu}[Real-World Schedulers]
Modern OS kernels use sophisticated variants of multilevel feedback queues. Linux's \emph{Completely Fair Scheduler (CFS)} uses a red-black tree keyed by \emph{virtual runtime} (effectively a dynamic priority) to approximate ideal fair sharing. Windows uses a priority-based preemptive scheduler with 32 priority levels and dynamic priority boosting for I/O completion and foreground windows. Both implement the core ideas you have learnt: preemption, time slices, priority, and ageing/boosting.
\end{savaistu}

\courssub{Comparing All the Scheduling Algorithms}

\begin{popfigure}
\begin{center}
\begin{tikzpicture}
\begin{axis}[
    ybar, bar width=14pt,
    width=0.85\linewidth, height=5.5cm,
    symbolic x coords={FCFS, SJF, SRTF, RR},
    xtick=data,
    ymin=0, ymax=22,
    ylabel={Average turnaround time (ms)},
    nodes near coords, every node near coord/.append style={font=\scriptsize},
    axis lines=left,
    ymajorgrids=true, grid style={popline},
]
\addplot[fill=popblue] coordinates {(FCFS,15.25) (SJF,14.25) (SRTF,13) (RR,19.25)};
\end{axis}
\end{tikzpicture}
\end{center}
\legende{Average turnaround time for the four algorithms on our example ($P_1$--$P_4$). SRTF wins on average, but RR gives the best responsiveness.}
\end{popfigure}

\begin{center}
\begin{tabularx}{\linewidth}{|l|X|X|X|X|}
\hline
\rowcolor{popblueL} Algorithm & Preemptive? & Starvation risk & Overhead & Best use \\
\hline
FCFS & No & No & Very low & Simple batch systems \\
\hline
SJF & No & Yes (long jobs) & Low & Batch with known bursts \\
\hline
SRTF & Yes & Yes (long jobs) & Moderate & Batch, optimal average wait \\
\hline
RR & Yes & No & High if $q$ small & Interactive / time-sharing \\
\hline
Priority & Both & Yes (cured by ageing) & Low & Real-time, mixed urgency \\
\hline
\end{tabularx}
\end{center}

\begin{aretenir}[Key Points of Lesson 47]
\begin{itemize}
\item Preemption = the OS may interrupt a running process (timer interrupt) and reassign the CPU; it costs a context switch each time.
\item SRTF = preemptive SJF: re-evaluate the shortest \emph{remaining} time at \textbf{every arrival}.
\item RR = circular FIFO with quantum $q$; unfinished processes go to the back of the queue; response time bounded by $(n-1)q$.
\item $q$ too large $\Rightarrow$ behaves like FCFS; $q$ too small $\Rightarrow$ overhead dominates (thrashing).
\item Preemptive priority can starve low-priority processes; \textbf{ageing} is the cure.
\item Multilevel feedback queues combine the advantages of RR (responsiveness) and priority (urgency) with automatic adaptation.
\end{itemize}
\end{aretenir}

\begin{exoresolu}[SRTF with a Late Arrival]
Processes: $A$ (arrival 0, burst 6), $B$ (arrival 2, burst 3), $C$ (arrival 4, burst 2). Give the Gantt chart under SRTF and the average turnaround time.
\tcblower
\begin{itemize}
\item $t=0$: only $A$; it runs.
\item $t=2$: $B$ arrives (3). $A$ has $6-2=4$ left; $3 < 4$, so $A$ is preempted, $B$ runs.
\item $t=4$: $C$ arrives (2). $B$ has $3-2=1$ left; $1 < 2$, so $B$ continues and finishes at $t=5$.
\item $t=5$: choose between $C$ (2) and $A$ (4): $C$ runs $5\!-\!7$, then $A$ runs $7\!-\!11$.
\end{itemize}
Turnarounds: $A: 11-0=11$, $B: 5-2=3$, $C: 7-4=3$. Average $= \dfrac{11+3+3}{3} = \dfrac{17}{3} \approx \SI{5.67}{ms}$.
\end{exoresolu}

\begin{experience}[Simulating a Scheduler (Guided Mini-Project)]
Write a short Python script (or pseudocode) that simulates the RR scheduler for a given list of processes (arrival, burst) and a quantum $q$. Use a list as a queue (\texttt{pop(0)} for dequeue, \texttt{append()} for enqueue). Print the Gantt chart as a list of tuples \texttt{(pid, start, end)}. Test it with the examples from this section. This reinforces the exact queue mechanics and helps avoid off-by-one errors in exams.
\end{experience}

\begin{exercice}[Practice]
Using $P_1(0,7)$, $P_2(2,4)$, $P_3(4,1)$, $P_4(5,4)$:
(a) Draw the SRTF Gantt chart and compute the average waiting time.
(b) Draw the Round Robin Gantt chart with $q=2$ (use the arrival-first convention from the Method box) and compute the average waiting time.
(c) State which algorithm gives the better average waiting time and which gives the better first response time for $P_3$.
\end{exercice}

\begin{corrige}[Exercise 1]
(a) \textbf{SRTF Trace:}
\begin{itemize}
\item $t=0$: $P_1$ runs (rem 7).
\item $t=2$: $P_2$ arrives (4). $P_1$ rem 5. $4<5$ $\to$ Preempt $P_1$, Run $P_2$.
\item $t=4$: $P_3$ arrives (1). $P_2$ rem 2. $1<2$ $\to$ Preempt $P_2$, Run $P_3$.
\item $t=5$: $P_3$ finishes. Ready: $P_2$(2), $P_1$(5), $P_4$(4) [arrives at 5]. Shortest $P_2$(2) runs.
\item $t=7$: $P_2$ finishes. Ready: $P_4$(4), $P_1$(5). $P_4$ runs.
\item $t=11$: $P_4$ finishes. $P_1$ runs $11\!-\!16$.
\end{itemize}
Gantt: $P_1(0-2)$, $P_2(2-4)$, $P_3(4-5)$, $P_2(5-7)$, $P_4(7-11)$, $P_1(11-16)$.
Waiting: $P_1=(2-0)+(11-7)=2+4=6$? No. $P_1$ runs 0-2, waits 2-11 (9), runs 11-16. Wait=9. $P_2$: waits 4-5 (1). Wait=1. $P_3$: wait=0. $P_4$: arrives 5, starts 7. Wait=2. Avg Wait $= (9+1+0+2)/4 = \mathbf{3\ ms}$.

(b) \textbf{RR $q=2$ Trace (Arrivals first, then Preemption):}
\begin{itemize}
\item $t=0$: Queue $[P_1]$. Run $P_1$ $0\!-\!2$ (rem 5).
\item $t=2$: Arr $P_2$. Queue $[P_1, P_2]$. Preempt $P_1 \to [P_2, P_1]$. Run $P_2$ $2\!-\!4$ (rem 2).
\item $t=4$: Arr $P_3$. Queue $[P_1, P_3]$. Preempt $P_2 \to [P_1, P_3, P_2]$. Run $P_1$ $4\!-\!6$ (rem 3).
\item $t=5$: Arr $P_4$. Running $P_1$. Queue $[P_3, P_2, P_4]$.
\item $t=6$: Preempt $P_1 \to [P_3, P_2, P_4, P_1]$. Run $P_3$ $6\!-\!7$ (finishes, burst 1).
\item $t=7$: Queue $[P_2, P_4, P_1]$. Run $P_2$ $7\!-\!9$ (finishes, rem 2).
\item $t=9$: Queue $[P_4, P_1]$. Run $P_4$ $9\!-\!11$ (rem 2).
\item $t=11$: Preempt $P_4 \to [P_1, P_4]$. Run $P_1$ $11\!-\!13$ (rem 1).
\item $t=13$: Preempt $P_1 \to [P_4, P_1]$. Run $P_4$ $13\!-\!15$ (finishes).
\item $t=15$: Queue $[P_1]$. Run $P_1$ $15\!-\!16$ (finishes).
\end{itemize}
Gantt: $P_1(0-2)$, $P_2(2-4)$, $P_1(4-6)$, $P_3(6-7)$, $P_2(7-9)$, $P_4(9-11)$, $P_1(11-13)$, $P_4(13-15)$, $P_1(15-16)$.
Turnaround: $P_1=16$, $P_2=7$, $P_3=3$, $P_4=10$. Waiting: $P_1=9$, $P_2=3$, $P_3=2$, $P_4=6$. Avg Wait $= 20/4 = \mathbf{5\ ms}$.

(c) \textbf{Comparison:} SRTF gives the better average waiting time (3 ms vs 5 ms). For $P_3$ (arrives 4), first response: SRTF starts at $t=4$ (immediate), RR starts at $t=6$. SRTF gives better first response for $P_3$ as well (0 ms wait vs 2 ms wait).
\end{corrige}

\coursec{Section 4: Process Synchronisation (Lesson 48)}

In the previous section we saw how the scheduler decides \emph{which} process runs next. But when several processes run concurrently and share data, the order in which their individual instructions interleave becomes critical. A tiny change in timing can turn a correct program into one that loses money, corrupts a file, or hangs forever. This section studies the techniques that guarantee correct cooperation: \cle{process synchronisation}.

\courssub{Why Synchronisation is Needed: Cooperating Processes and Interleaving}

Processes are often independent, but many useful programs require \emph{cooperating processes} that share data through shared memory, shared files, or message queues. Consider a mobile banking application used by a family in Douala: the father checks the balance while the mother transfers money to a sibling in Bafoussam. Both operations are carried out by processes that read and write the same \texttt{balance} variable in the bank's server memory.

The danger comes from the fact that a high-level statement such as \texttt{balance = balance + amount} is \emph{not} executed in one indivisible step. It is compiled into at least three machine instructions: \textbf{load} the value of \texttt{balance} into a register, \textbf{add} the amount, \textbf{store} the result back. The scheduler may switch from one process to another between any two of these instructions.

\begin{popfigure}
\begin{tikzpicture}[scale=0.9, every node/.style={font=\small}]
  % Process P1 (Father)
  \draw[popblue, thick] (0,4) node[left]{P1 (Father)} -- (10,4);
  \draw[popblue, fill=popblueL] (1,3.7) rectangle (2.5,4.3) node[midway]{read};
  \draw[popblue, fill=popblueL] (2.5,3.7) rectangle (4,4.3) node[midway]{add 10\,000};
  \draw[popblue, fill=popblueL] (4,3.7) rectangle (5.5,4.3) node[midway]{write};
  % Process P2 (Mother)
  \draw[poporange, thick] (0,2) node[left]{P2 (Mother)} -- (10,2);
  \draw[poporange, fill=popgold] (1.5,1.7) rectangle (3,2.3) node[midway]{read};
  \draw[poporange, fill=popgold] (3,1.7) rectangle (4.5,2.3) node[midway]{sub 5\,000};
  \draw[poporange, fill=popgold] (5.5,1.7) rectangle (7,2.3) node[midway]{write};
  % Context switch markers
  \draw[->, popmuted, dashed] (2.5,3.7) -- (2.5,2.5);
  \draw[->, popmuted, dashed] (5.5,2.5) -- (5.5,3.7);
  \node[popmuted] at (8.5,3) {context switches};
\end{tikzpicture}
\legende{Interleaving of the machine instructions of two cooperating processes sharing the variable \texttt{balance}. The scheduler switches between P1 and P2 between their read, modify and write steps.}
\end{popfigure}

The figure shows a possible interleaving:
\begin{enumerate}
  \item P1 reads balance (50\,000).
  \item P2 reads balance (50\,000).
  \item P1 computes $50\,000 + 10\,000 = 60\,000$ and writes it.
  \item P2 computes $50\,000 - 5\,000 = 45\,000$ and writes it.
\end{enumerate}
The final balance is 45\,000 FCFA — the father's deposit of 10\,000 FCFA has been \emph{lost}. This is a classic \cle{race condition}.

\begin{definition}[Race Condition]
A \textbf{race condition} occurs when multiple processes access and manipulate shared data concurrently, and the final result depends on the particular order in which the accesses take place.
\end{definition}

\begin{attention}[Lost Update — The Classic Trap]
The example above is called a \emph{lost update}. It happens because the sequence \texttt{read} $\to$ \texttt{modify} $\to$ \texttt{write} is not atomic: the scheduler can switch processes between any two machine instructions. Never assume that a high-level statement like \texttt{balance = balance + amount} executes in one indivisible step.
\end{attention}

\courssub{The Critical Section Problem}

To avoid race conditions, we must ensure that when one process is executing the sensitive part of its code that accesses shared data, no other process enters its own sensitive part for the same data. This sensitive part is called a \cle{critical section}.

\begin{definition}[Critical Section]
A \textbf{critical section} is a code segment in which a process accesses shared variables (or shared resources such as files or printers) and which must not be executed concurrently with another process's critical section for the same resource.
\end{definition}

Each process can therefore be seen as a loop of four parts: an \emph{entry section} (requesting permission to enter), the \emph{critical section}, an \emph{exit section} (releasing the resource), and a \emph{remainder section} (all the other code). A correct solution to the critical-section problem must satisfy three requirements:

\begin{propriete}[Requirements of a Critical-Section Solution]
\begin{enumerate}
  \item \textbf{Mutual Exclusion:} If process $P_i$ is executing in its critical section, no other process $P_j$ ($j \neq i$) may be in its critical section for the same resource.
  \item \textbf{Progress:} If no process is currently in its critical section and some processes wish to enter theirs, then the choice of the next process to enter cannot be postponed indefinitely.
  \item \textbf{Bounded Waiting:} There exists a bound on the number of times other processes are allowed to enter their critical sections after a process has requested entry and before that request is granted. (No starvation.)
\end{enumerate}
\end{propriete}

\courssub{Software Approaches (Brief Overview)}

Before hardware support existed, computer scientists devised pure software algorithms.

\begin{itemize}
  \item \textbf{Strict Alternation:} A shared variable \texttt{turn} indicates whose turn it is. Process $i$ busy-waits with \texttt{while (turn != i);}. This satisfies mutual exclusion but violates \emph{progress}: it forces strict alternation even when the other process has no work to do.
  \item \textbf{Peterson's Solution (1981):} For two processes, it uses two boolean flags \texttt{interested[2]} and a variable \texttt{turn}. It satisfies all three requirements on paper, but relies on \emph{busy waiting} and assumes that reads and writes to shared variables are atomic — an assumption that fails on modern multi-core processors without memory barriers.
\end{itemize}

\begin{attention}[GCE Note]
Software solutions like Peterson's are \textbf{not} required in detail for the GCE examination. You should know they exist historically, but the examinable solutions are based on \textbf{hardware atomic instructions} and \textbf{semaphores}.
\end{attention}

\courssub{Hardware Support for Synchronisation}

Modern processors provide atomic instructions that cannot be interrupted once started.

\begin{itemize}
  \item \textbf{Disabling Interrupts:} On a single-core system, a process can disable interrupts before entering its critical section and re-enable them after leaving, which prevents any context switch. \emph{Drawbacks:} it only works on uniprocessors, it gives too much power to user processes, and it blocks all I/O handling.
  \item \textbf{Test-and-Set (TSL) Instruction:} An atomic hardware instruction that reads a boolean variable and sets it to \texttt{true} in one indivisible operation:
\end{itemize}

\begin{center}
\begin{tabular}{|l|}
\hline
\texttt{boolean test\_and\_set(boolean *target) \{} \\
\texttt{\quad boolean rv = *target;} \\
\texttt{\quad *target = true;} \\
\texttt{\quad return rv;} \\
\texttt{\}} \\
\hline
\end{tabular}
\end{center}

A spinlock using TSL then looks like this:

\begin{center}
\begin{tabular}{|l|}
\hline
\texttt{while (test\_and\_set(\&lock)) ; // busy wait} \\
\texttt{// critical section} \\
\texttt{lock = false; // release} \\
\hline
\end{tabular}
\end{center}

This works on multiprocessors but still uses busy waiting, which wastes CPU cycles while a process spins on the lock.

\courssub{Semaphores: The Classic Synchronisation Tool}

A \cle{semaphore} is an integer variable that, apart from initialisation, is accessed only through two standard atomic operations: \texttt{wait()} (originally \texttt{P}, from the Dutch \emph{proberen} — to test) and \texttt{signal()} (originally \texttt{V}, from \emph{verhogen} — to increment). The atomicity of these operations is guaranteed by the operating system with hardware support.

\begin{definition}[Semaphore]
A \textbf{semaphore} $S$ is an integer variable that can only be accessed via two indivisible (atomic) operations:
\begin{align*}
\texttt{wait(S):} &\quad \text{while } (S \le 0)\ \text{wait;}\quad S = S - 1; \\
\texttt{signal(S):} &\quad S = S + 1;
\end{align*}
In a good implementation, the waiting inside \texttt{wait} is done by \emph{blocking} the process (putting it to sleep in a queue) rather than by busy waiting, so no CPU time is wasted.
\end{definition}

\begin{itemize}
  \item \textbf{Binary semaphore (mutex):} Takes only the values 0 or 1. Used for mutual exclusion.
  \item \textbf{Counting semaphore:} Takes any non-negative integer value. Used to control access to a pool of several identical resources (for example, 3 printers $\to$ semaphore initialised to 3).
\end{itemize}

\begin{propriete}[Atomicity of wait and signal]
The operations \texttt{wait} and \texttt{signal} are executed atomically: once started, they run to completion without interruption. This is the key property that makes semaphores a correct synchronisation primitive.
\end{propriete}

\courssub{Protecting a Critical Section with a Binary Semaphore (Mutex)}

The standard pattern for mutual exclusion uses a binary semaphore \texttt{mutex} initialised to 1.

\begin{methode}[Protecting a Critical Section with a Mutex]
\begin{enumerate}
  \item Declare a binary semaphore \texttt{mutex} and initialise it to \texttt{1}.
  \item At the \textbf{entry} of the critical section, execute \texttt{wait(mutex)}.
  \item Execute the critical section code.
  \item At the \textbf{exit} of the critical section, execute \texttt{signal(mutex)}.
\end{enumerate}
\end{methode}

\begin{center}
\begin{tabular}{|l|}
\hline
\texttt{wait(mutex);} \\
\texttt{// critical section: access shared data} \\
\texttt{signal(mutex);} \\
\texttt{// remainder section} \\
\hline
\end{tabular}
\end{center}

\begin{popfigure}
\begin{tikzpicture}[scale=0.85, every node/.style={font=\small}]
  % Time axis
  \draw[->, thick] (0,0) -- (14,0) node[right]{Time};
  % Mutex value
  \draw[popblue, thick] (0,3) node[left]{mutex} -- (14,3);
  \foreach \x/\val in {0.2/1, 2/0, 5/0, 7.5/1, 9/0, 12/1} {
    \draw[popblue] (\x,2.8) -- (\x,3.2) node[above]{\val};
  }
  % Process P1
  \draw[poporange, thick] (0,1.5) node[left]{P1} -- (14,1.5);
  \draw[poporange, fill=popgold] (1,1.2) rectangle (4,1.8) node[midway]{CS1};
  \draw[poporange, fill=popgold] (8,1.2) rectangle (11,1.8) node[midway]{CS2};
  % Process P2
  \draw[popgreen, thick] (0,0.5) node[left]{P2} -- (14,0.5);
  \draw[popgreen, fill=popgreenL] (5,0.2) rectangle (7.5,0.8) node[midway]{CS};
  % wait/signal markers
  \draw[->, popmuted] (1,1.8) -- (1,2.3) node[above]{wait};
  \draw[->, popmuted] (4,1.8) -- (4,2.3) node[above]{signal};
  \draw[->, popmuted] (5,0.8) -- (5,1.3) node[above]{wait};
  \draw[->, popmuted] (7.5,0.8) -- (7.5,1.3) node[above]{signal};
  \draw[->, popmuted] (8,1.8) -- (8,2.3) node[above]{wait};
  \draw[->, popmuted] (11,1.8) -- (11,2.3) node[above]{signal};
\end{tikzpicture}
\legende{Timeline showing mutual exclusion enforced by a mutex semaphore. While P1 is in its critical section (CS1), mutex = 0, so P2 blocks at \texttt{wait()} until P1 executes \texttt{signal()}.}
\end{popfigure}

\begin{attention}[Common Mistakes with Semaphores]
\begin{enumerate}
  \item \textbf{Wrong order:} Executing \texttt{signal} before \texttt{wait} (or omitting one of them) breaks mutual exclusion or causes deadlock.
  \item \textbf{Forgetting \texttt{signal}:} If a process crashes or returns early without signalling, the semaphore stays at 0 and \emph{all} other processes block forever.
  \item \textbf{Initialisation errors:} \texttt{mutex} must start at 1; \texttt{empty} at the buffer size $N$; \texttt{full} at 0. Wrong initial values cause immediate deadlock or buffer overflow.
\end{enumerate}
\end{attention}

\courssub{Classical Problem: Producer–Consumer (Bounded Buffer)}

A \cle{producer} process generates data items and puts them into a shared buffer of fixed size $N$; a \cle{consumer} process removes items from the buffer. The producer must not add an item when the buffer is full, the consumer must not remove an item when the buffer is empty, and the two must never access the buffer's internal pointers simultaneously.

We need three semaphores:
\begin{itemize}
  \item \texttt{mutex} (binary, initialised to 1) — protects the buffer data structure itself.
  \item \texttt{empty} (counting, initialised to $N$) — counts empty slots; the producer waits on it.
  \item \texttt{full} (counting, initialised to 0) — counts filled slots; the consumer waits on it.
\end{itemize}

\begin{popfigure}
\begin{tikzpicture}[scale=0.9, every node/.style={font=\small}]
  % Buffer
  \draw[popdark, thick] (2,1) rectangle (10,3);
  \foreach \i in {0,...,7} {
    \draw[popline] (2+\i,1) -- (2+\i,3);
    \node at (2.5+\i,2) {\i};
  }
  \node at (6,3.5) {Shared buffer (size $N=8$)};
  % Producer
  \node[draw, fill=popblueL, rounded corners, thick] (prod) at (1,5) {Producer};
  \draw[->, thick, popblue] (prod.south) -- (3,3.1) node[midway, left, align=center]{wait(empty)\\wait(mutex)};
  % Consumer
  \node[draw, fill=popgold, rounded corners, thick] (cons) at (11,5) {Consumer};
  \draw[->, thick, poporange] (cons.south) -- (9,3.1) node[midway, right, align=center]{wait(full)\\wait(mutex)};
  % Semaphore boxes
  \node[draw, fill=popgreenL, rounded corners] at (4,-0.5) {mutex = 1};
  \node[draw, fill=popgreenL, rounded corners] at (7,-0.5) {empty = $N$};
  \node[draw, fill=popgreenL, rounded corners] at (10,-0.5) {full = 0};
\end{tikzpicture}
\legende{Bounded buffer with a producer, a consumer and the three semaphores. The producer waits on \texttt{empty} before inserting; the consumer waits on \texttt{full} before removing; \texttt{mutex} protects the buffer itself.}
\end{popfigure}

\begin{exemplebox}[Producer–Consumer Pseudocode]
\textbf{Shared declarations:}
\begin{center}
\begin{tabular}{|l|}
\hline
\texttt{semaphore mutex = 1;} \\
\texttt{semaphore empty = N;} \\
\texttt{semaphore full = 0;} \\
\texttt{item buffer[N];} \\
\texttt{int in = 0, out = 0;} \\
\hline
\end{tabular}
\end{center}
\textbf{Producer:}
\begin{center}
\begin{tabular}{|l|}
\hline
\texttt{while (true) \{} \\
\texttt{\quad item next = produce\_item();} \\
\texttt{\quad wait(empty);} \\
\texttt{\quad wait(mutex);} \\
\texttt{\quad buffer[in] = next;} \\
\texttt{\quad in = (in + 1) \%\ N;} \\
\texttt{\quad signal(mutex);} \\
\texttt{\quad signal(full);} \\
\texttt{\}} \\
\hline
\end{tabular}
\end{center}
\textbf{Consumer:}
\begin{center}
\begin{tabular}{|l|}
\hline
\texttt{while (true) \{} \\
\texttt{\quad wait(full);} \\
\texttt{\quad wait(mutex);} \\
\texttt{\quad item next = buffer[out];} \\
\texttt{\quad out = (out + 1) \%\ N;} \\
\texttt{\quad signal(mutex);} \\
\texttt{\quad signal(empty);} \\
\texttt{\quad consume\_item(next);} \\
\texttt{\}} \\
\hline
\end{tabular}
\end{center}
\end{exemplebox}

\begin{attention}[Order of the wait() Calls]
The order \texttt{wait(empty)} \emph{then} \texttt{wait(mutex)} in the producer (and \texttt{wait(full)} \emph{then} \texttt{wait(mutex)} in the consumer) is crucial. If we reversed them (mutex first), a deadlock could occur: the producer holds \texttt{mutex}, finds the buffer full, and waits on \texttt{empty}; the consumer needs \texttt{mutex} to remove an item — a circular wait.
\end{attention}

\courssub{Other Classical Synchronisation Problems (Exam Awareness)}

The GCE syllabus expects you to \emph{name} and briefly \emph{describe} these problems; you are not required to write full semaphore code for them.

\begin{itemize}
  \item \textbf{Readers–Writers Problem:} Multiple readers may read a shared database simultaneously, but a writer must have exclusive access. Variants: \emph{readers' preference} (a continuous stream of readers can starve writers) and \emph{writers' preference} (writers can starve readers).
  \item \textbf{Dining Philosophers Problem:} Five philosophers sit around a table with five forks, one between each pair of plates. Each philosopher needs \emph{two} forks to eat. The challenge is to design a protocol that avoids deadlock (for example, if all five pick up their left fork simultaneously, nobody can ever pick up the right one) and starvation.
\end{itemize}

\courssub{Monitors: A Higher-Level Construct}

Semaphores are low-level and error-prone (a forgotten \texttt{signal}, a wrong order of calls). A \cle{monitor} is a higher-level synchronisation construct, supported by some programming languages or provided by the operating system, that encapsulates shared data together with the operations on it and guarantees that only one process can be active inside the monitor at any time.

\begin{definition}[Monitor]
A \textbf{monitor} is a synchronisation construct that consists of:
\begin{itemize}
  \item shared data (variables);
  \item a set of procedures (methods) that operate on that data;
  \item initialisation code.
\end{itemize}
\textbf{Key property:} mutual exclusion is \textbf{automatic} — only one process can execute a monitor procedure at a time. For waiting on conditions, monitors provide \emph{condition variables} with the operations \texttt{wait()}, \texttt{signal()} and \texttt{broadcast()}.
\end{definition}

\begin{savaistu}[Monitors in Java]
In Java, every object has an intrinsic lock (a monitor). The \texttt{synchronized} keyword on a method or block provides the mutual exclusion, while \texttt{wait()}, \texttt{notify()} and \texttt{notifyAll()} on the object act as condition-variable operations. This is the practical realisation of the monitor concept that you will meet in programming.
\end{savaistu}

\courssub{Worked Example: Tracing a Producer–Consumer Execution}

\begin{exoresolu}[Trace of a Bounded Buffer with N = 3]
\textbf{Scenario:} Buffer size $N = 3$. Initial state: \texttt{empty = 3}, \texttt{full = 0}, \texttt{mutex = 1}. The producer produces items A then B; the consumer then consumes one item. Show the semaphore values after each step.
\tcblower
We trace the execution assuming the producer runs first, then the consumer.

\begin{center}
\begin{tabular}{|c|c|l|c|c|c|l|}
\hline
\rowcolor{popblueL}
Step & Process & Operation & empty & full & mutex & Buffer (in/out) \\
\hline
0 & --- & Initial & 3 & 0 & 1 & [\ ,\ ,\ ] (0/0) \\
\hline
1 & P & wait(empty) & 2 & 0 & 1 & \\
\hline
2 & P & wait(mutex) & 2 & 0 & 0 & \\
\hline
3 & P & put A & 2 & 0 & 0 & [A,\ ,\ ] (1/0) \\
\hline
4 & P & signal(mutex) & 2 & 0 & 1 & \\
\hline
5 & P & signal(full) & 2 & 1 & 1 & \\
\hline
6 & P & wait(empty) & 1 & 1 & 1 & \\
\hline
7 & P & wait(mutex) & 1 & 1 & 0 & \\
\hline
8 & P & put B & 1 & 1 & 0 & [A, B,\ ] (2/0) \\
\hline
9 & P & signal(mutex) & 1 & 1 & 1 & \\
\hline
10 & P & signal(full) & 1 & 2 & 1 & \\
\hline
11 & C & wait(full) & 1 & 1 & 1 & \\
\hline
12 & C & wait(mutex) & 1 & 1 & 0 & \\
\hline
13 & C & take A & 1 & 1 & 0 & [\ , B,\ ] (2/1) \\
\hline
14 & C & signal(mutex) & 1 & 1 & 1 & \\
\hline
15 & C & signal(empty) & 2 & 1 & 1 & \\
\hline
\end{tabular}
\end{center}

\textbf{Final state:} \texttt{empty = 2}, \texttt{full = 1}, \texttt{mutex = 1}; the buffer contains B at index 1, with \texttt{in = 2} and \texttt{out = 1}. Mutual exclusion held throughout: \texttt{mutex} never became negative and the two processes were never inside the buffer code at the same time.
\end{exoresolu}

\courssub{Consolidation Exercises}

\begin{exercice}[Exercise 1 — Race Condition Identification]
Two processes share an integer \texttt{x} initially equal to 0.
Process P1 executes \texttt{x = x + 1;} one thousand times in a loop; process P2 executes \texttt{x = x - 1;} one thousand times in a loop. Both run concurrently without any synchronisation.
\begin{enumerate}
  \item What is the expected final value of \texttt{x} if every operation were atomic?
  \item Explain why the actual final value may be different. Give a concrete interleaving of the machine instructions that produces a wrong result.
\end{enumerate}
\end{exercice}

\begin{exercice}[Exercise 2 — Semaphore Initialisation]
A system has 4 identical tape drives. Processes request a tape drive, use it, then release it.
\begin{enumerate}
  \item What type of semaphore should be used, and with what initial value?
  \item Write the \texttt{wait}/\texttt{signal} code for a process that needs one tape drive.
\end{enumerate}
\end{exercice}

\begin{exercice}[Exercise 3 — Producer–Consumer Deadlock]
In the bounded-buffer solution, suppose a programmer mistakenly writes the producer as:
\begin{center}
\begin{tabular}{|l|}
\hline
\texttt{wait(mutex);} \\
\texttt{wait(empty);} \\
\texttt{// ... put item ...} \\
\texttt{signal(mutex);} \\
\texttt{signal(full);} \\
\hline
\end{tabular}
\end{center}
Explain how this can lead to deadlock. Illustrate with a scenario where $N = 1$.
\end{exercice}

\begin{exercice}[Exercise 4 — Critical Section Requirements]
For each of the following scenarios, state which of the three critical-section requirements (mutual exclusion, progress, bounded waiting) is violated:
\begin{enumerate}
  \item A process enters its critical section and never leaves (infinite loop inside it).
  \item Two processes are in their critical sections simultaneously.
  \item Process P1 wants to enter its critical section, but P2 repeatedly enters and leaves its own critical section, so P1 never gets in.
\end{enumerate}
\end{exercice}

\begin{corrige}[Exercise 1]
\begin{enumerate}
  \item Expected final value: 0, since the 1000 increments and the 1000 decrements cancel out.
  \item The statement \texttt{x = x + 1} compiles into at least three machine instructions: \texttt{load x} into a register, \texttt{add 1}, \texttt{store} the register back to \texttt{x}. A losing interleaving:
    \begin{enumerate}
      \item P1 loads x (value 0) into its register.
      \item P1 increments its register to 1.
      \item The scheduler switches to P2, which loads x (still 0, because P1 has not stored yet).
      \item P2 decrements its register to $-1$ and stores $-1$ into x.
      \item P1 resumes and stores 1 into x, overwriting P2's result.
    \end{enumerate}
    One decrement has been lost. Repeated over many iterations, the final value can be far from 0.
\end{enumerate}
\end{corrige}

\begin{corrige}[Exercise 2]
\begin{enumerate}
  \item A \textbf{counting semaphore}, initialised to 4 (the number of available tape drives).
  \item The code for one process:
\end{enumerate}
\begin{center}
\begin{tabular}{|l|}
\hline
\texttt{wait(tape\_drives); // acquire one drive} \\
\texttt{// use the tape drive} \\
\texttt{signal(tape\_drives); // release it} \\
\hline
\end{tabular}
\end{center}
\end{corrige}

\begin{corrige}[Exercise 3]
With $N = 1$, \texttt{empty} is initialised to 1.
\begin{enumerate}
  \item Producer executes \texttt{wait(mutex)} $\to$ \texttt{mutex} = 0 (producer now holds the mutex).
  \item Producer executes \texttt{wait(empty)} $\to$ \texttt{empty} = 0 (the slot is available, so it proceeds).
  \item Producer puts the item, then \texttt{signal(mutex)} $\to$ \texttt{mutex} = 1, and \texttt{signal(full)} $\to$ \texttt{full} = 1.
  \item The producer loops again: \texttt{wait(mutex)} $\to$ \texttt{mutex} = 0 (it holds the mutex again).
  \item The producer executes \texttt{wait(empty)}: \texttt{empty} is 0, so the producer \textbf{blocks inside wait(empty) while still holding the mutex}.
  \item The consumer runs: \texttt{wait(full)} succeeds (\texttt{full} goes from 1 to 0), then \texttt{wait(mutex)}: \texttt{mutex} is 0 (held by the blocked producer), so the consumer blocks too.
  \item Both processes are blocked: the producer waits for \texttt{empty} (which only the consumer can signal), and the consumer waits for \texttt{mutex} (held by the producer). This is a \textbf{deadlock}.
\end{enumerate}
\end{corrige}

\begin{corrige}[Exercise 4]
\begin{enumerate}
  \item \textbf{Progress} is violated: no other process can ever enter because one process hogs the critical section indefinitely.
  \item \textbf{Mutual exclusion} is violated.
  \item \textbf{Bounded waiting} is violated: P1 waits indefinitely while P2 enters again and again (starvation of P1).
\end{enumerate}
\end{corrige}

\courssub{Key Points to Remember}

\begin{aretenir}[Process Synchronisation — Summary for the GCE]
\begin{itemize}
  \item \textbf{Race condition:} the result depends on the interleaving of instructions; classic example: the lost update on a shared bank balance.
  \item \textbf{Critical section:} code accessing a shared resource; a solution must guarantee \textbf{mutual exclusion}, \textbf{progress} and \textbf{bounded waiting}.
  \item \textbf{Semaphore:} an integer accessed only through the atomic operations \texttt{wait} (P) and \texttt{signal} (V). A binary semaphore (mutex) gives mutual exclusion; a counting semaphore manages a pool of resources.
  \item \textbf{Mutex pattern:} \texttt{wait(mutex);} critical section; \texttt{signal(mutex);} with \texttt{mutex = 1} initially.
  \item \textbf{Producer–Consumer:} three semaphores — \texttt{mutex = 1}, \texttt{empty = N}, \texttt{full = 0}. Always wait on the condition semaphore (\texttt{empty}/\texttt{full}) \emph{before} waiting on \texttt{mutex}.
  \item \textbf{Classical problems:} Readers–Writers and Dining Philosophers — know their statements.
  \item \textbf{Monitors:} a high-level construct with automatic mutual exclusion and condition variables (Java \texttt{synchronized}).
  \item \textbf{Common errors:} wrong order of \texttt{wait}/\texttt{signal}, a missing \texttt{signal}, wrong initial values.
\end{itemize}
\end{aretenir}

\coursec{Section 5: Deadlock Handling and Integration (Lessons 49--50)}

In Section 4 we saw that synchronisation mechanisms (semaphores, critical sections, monitors) allow processes to share resources safely. But synchronisation has a dark side: when several processes hold some resources and wait for others, the system can freeze completely. This situation, called a \cle{deadlock}, is one of the most feared problems in operating systems. In this final section we define deadlock precisely, model it with graphs, state the four conditions that make it possible, and study the four families of strategies to handle it: prevention, avoidance, detection and recovery --- plus the surprising ``do nothing'' strategy. We close the chapter with an integration activity combining scheduling, synchronisation and deadlock analysis.

\courssub{What is a Deadlock?}

\textbf{A concrete situation.} Imagine the computer room of a school in Yaoundé with one printer and one scanner. Process $P_1$ (a student scanning a document) has grabbed the scanner and now requests the printer to print the scanned page. At the same time, process $P_2$ has grabbed the printer and now requests the scanner to scan the page it wants to print next. $P_1$ waits for the printer, held by $P_2$; $P_2$ waits for the scanner, held by $P_1$. Neither can ever proceed, and neither will ever release what it holds. The two processes are blocked \emph{forever}.

\begin{definition}[Deadlock]
A \cle{deadlock} is a situation in which a set of processes is permanently blocked because each process in the set holds at least one resource and is waiting for a resource held by another process in the set. No process can proceed, release its resources, or be awakened without external intervention.
\end{definition}

\begin{savaistu}[The traffic jam analogy]
Deadlock also occurs on roads: at a crossroads in Douala, if four cars each enter the junction and each waits for the car on its right to move, all four are stuck forever --- a real-world circular wait. Traffic lights are a \emph{prevention} mechanism; a traffic policeman who orders one car to reverse is a \emph{recovery} mechanism.
\end{savaistu}

\courssub{System Model and the Resource Allocation Graph}

To reason rigorously about deadlocks we need a model. The system contains \cle{resource types} (printer, scanner, memory region, file, semaphore\ldots), each type having one or more identical \cle{instances} (e.g.\ 2 printers). Every process uses a resource through a strict protocol:

\begin{enumerate}
\item \textbf{Request}: the process asks for the resource; if it is not free, the process \emph{waits}.
\item \textbf{Use}: the process works with the resource.
\item \textbf{Release}: the process gives the resource back.
\end{enumerate}

A deadlock can only arise inside this request--use--release cycle. The state of the system is pictured with a \cle{Resource Allocation Graph} (RAG):

\begin{itemize}
\item a \textbf{circle} represents a process;
\item a \textbf{rectangle} represents a resource type, with one \textbf{dot} per instance;
\item a \textbf{request edge} $P_i \to R_j$ means $P_i$ is waiting for an instance of $R_j$;
\item an \textbf{assignment edge} $R_j \to P_i$ (drawn from a dot) means one instance of $R_j$ is held by $P_i$.
\end{itemize}

\begin{popfigure}
\begin{center}
\begin{tikzpicture}[scale=1.05, every node/.style={font=\small}]
% processes
\node[circle, draw=popblue, thick, minimum size=0.9cm] (p1) at (0,0) {$P_1$};
\node[circle, draw=poppurple, thick, minimum size=0.9cm] (p2) at (6,0) {$P_2$};
% resources
\node[draw=poporange, thick, minimum width=1.1cm, minimum height=0.8cm] (r1) at (3,1.6) {};
\node at (3,2.15) {Printer};
\fill[poporange] (3,1.6) circle (0.07);
\node[draw=popgreen, thick, minimum width=1.1cm, minimum height=0.8cm] (r2) at (3,-1.6) {};
\node at (3,-2.15) {Scanner};
\fill[popgreen] (3,-1.6) circle (0.07);
% edges
\draw[->, very thick, popblue] (p1) -- (r1.west) node[midway, above, sloped, font=\footnotesize]{requests};
\draw[->, very thick, poporange] (r1.east) -- (p2) node[midway, above, sloped, font=\footnotesize]{held by};
\draw[->, very thick, poppurple] (p2) -- (r2.east) node[midway, below, sloped, font=\footnotesize]{requests};
\draw[->, very thick, popgreen] (r2.west) -- (p1) node[midway, below, sloped, font=\footnotesize]{held by};
\end{tikzpicture}
\end{center}
\legende{RAG of the printer--scanner deadlock: the cycle $P_1 \to \text{Printer} \to P_2 \to \text{Scanner} \to P_1$ with one instance per resource means a certain deadlock.}
\end{popfigure}

\begin{methode}[Analysing a Resource Allocation Graph]
\begin{enumerate}
\item Draw all processes, resources (with one dot per instance) and all request/assignment edges.
\item Look for a \textbf{cycle} in the graph.
\item If there is \textbf{no cycle}: no deadlock is possible.
\item If there is a cycle and \textbf{every resource in the cycle has exactly one instance}: deadlock is \emph{certain}.
\item If there is a cycle but some resources have \textbf{several instances}: deadlock is only \emph{possible} --- check whether another process outside the cycle can release an instance and break it.
\end{enumerate}
\end{methode}

\begin{attention}[A cycle does not always mean deadlock]
A very common exam mistake is to write ``cycle $\Rightarrow$ deadlock'' without qualification. This is only true when each resource involved has \textbf{one instance}. With several instances, a cycle is a \emph{necessary but not sufficient} condition: an instance held by a process \emph{outside} the cycle may be released and unblock the cycle.
\end{attention}

\courssub{The Four Necessary Conditions (Coffman Conditions)}

\begin{propriete}[Coffman necessary conditions for deadlock]
A deadlock can occur only if \textbf{all four} of the following conditions hold \emph{simultaneously}:
\begin{enumerate}
\item \textbf{Mutual exclusion}: at least one resource is non-sharable (only one process can use it at a time).
\item \textbf{Hold and wait}: a process holds at least one resource while waiting for another.
\item \textbf{No preemption}: a resource cannot be taken away from a process by force; it is released only voluntarily.
\item \textbf{Circular wait}: there exists a cycle of processes $P_1, P_2, \dots, P_n$ such that each waits for a resource held by the next (and $P_n$ waits for $P_1$).
\end{enumerate}
The four conditions are \emph{necessary}, not individually sufficient: removing any one of them makes deadlock impossible. (The justification is direct from the definition and is examinable as a short essay question.)
\end{propriete}

\courssub{Deadlock Prevention}

\cle{Deadlock prevention} consists in designing the system so that at least one Coffman condition can \emph{never} hold:

\begin{itemize}
\item \textbf{Negate mutual exclusion}: make resources sharable where possible (e.g.\ spooled printer queue). Often impossible by nature of the device.
\item \textbf{Negate hold and wait}: force each process to \textbf{request all its resources at once} before starting; it runs only when everything is available. Alternatively, a process must release everything before requesting a new resource.
\item \textbf{Negate no preemption}: if a process holding resources is denied a new request, forcibly take its resources away (works for memory or CPU, not for a printer).
\item \textbf{Negate circular wait}: impose a \textbf{total ordering} on resource types (printer $= 1$, scanner $= 2$, disk $= 3$\ldots) and require every process to request resources in \emph{increasing order}. Then a circular wait is impossible, because a cycle would require a resource number to be strictly smaller than itself.
\end{itemize}

\begin{attention}[Prevention has a cost]
Breaking ``hold and wait'' by requesting all resources at once is safe but wasteful: a process that needs the printer only at the end of its execution keeps it idle from the start, so \textbf{resource utilisation drops} and starvation may appear. Every deadlock strategy trades safety against performance --- a classic GCE discussion point.
\end{attention}

\courssub{Deadlock Avoidance: Safe States and the Banker's Algorithm}

Prevention is rigid. \cle{Avoidance} is more flexible: the system grants requests dynamically, but only when the request leaves the system in a \cle{safe state}. Avoidance requires that each process declares its \textbf{maximum claim} (Max) for each resource type in advance.

\begin{definition}[Safe and unsafe states]
A state is \textbf{safe} if there exists at least one ordering of the processes (a \emph{safe sequence}) such that each process can obtain its remaining needs from the currently available resources plus the resources released by the processes that finish before it. A state that is not safe is \textbf{unsafe}. An unsafe state does not guarantee a deadlock, but a deadlock is only reachable from an unsafe state.
\end{definition}

The \cle{Banker's algorithm} (Dijkstra) implements avoidance. Like a careful banker, the OS never grants a loan (resource) that could make it impossible to satisfy all its clients. For each process it maintains:

\begin{itemize}
\item \textbf{Max} (claim): the maximum number of instances of each resource the process may need;
\item \textbf{Allocation}: what the process currently holds;
\item \textbf{Need} $=$ Max $-$ Allocation: what it may still request;
\item \textbf{Available}: the free instances in the system (a single vector).
\end{itemize}

\begin{methode}[Banker's safety algorithm]
\begin{enumerate}
\item Compute $\text{Need} = \text{Max} - \text{Allocation}$ for every process.
\item Set Work $=$ Available. Mark all processes ``unfinished''.
\item Find an unfinished process $P_i$ with $\text{Need}_i \le \text{Work}$ (component by component). If none exists, go to step 5.
\item Pretend $P_i$ runs to completion and releases everything: Work $=$ Work $+$ Allocation$_i$; mark $P_i$ finished. Return to step 3.
\item If all processes are finished, the state is \textbf{safe} (the order found is a safe sequence); otherwise it is \textbf{unsafe}.
\end{enumerate}
\end{methode}

\begin{exemplebox}[Safety check with 3 processes and 2 resources]
Consider resources $A$ (10 instances) and $B$ (7 instances):

\begin{center}
\begin{tabular}{|l|cc|cc|cc|}
\hline
\rowcolor{popblueL}
 & \multicolumn{2}{c|}{Allocation} & \multicolumn{2}{c|}{Max} & \multicolumn{2}{c|}{Need} \\
\rowcolor{popblueL}
 & $A$ & $B$ & $A$ & $B$ & $A$ & $B$ \\
\hline
$P_1$ & 2 & 1 & 5 & 3 & 3 & 2 \\
$P_2$ & 3 & 2 & 4 & 4 & 1 & 2 \\
$P_3$ & 2 & 1 & 6 & 3 & 4 & 2 \\
\hline
\end{tabular}
\end{center}

Available $= (10-7,\; 7-4) = (3, 3)$.
\begin{itemize}
\item Work $= (3,3)$: $P_2$ has Need $(1,2) \le (3,3)$ $\Rightarrow$ run $P_2$; Work $= (3,3)+(3,2) = (6,5)$.
\item $P_1$ has Need $(3,2) \le (6,5)$ $\Rightarrow$ run $P_1$; Work $= (8,6)$.
\item $P_3$ has Need $(4,2) \le (8,6)$ $\Rightarrow$ run $P_3$; Work $= (10,7)$.
\end{itemize}
All processes finish: the state is \textbf{safe}, with safe sequence $\langle P_2, P_1, P_3 \rangle$.
\end{exemplebox}

When a new request arrives, the OS \emph{pretends} to grant it (updates Allocation, Need, Available), then runs the safety check: if the resulting state is safe, the request is granted; otherwise the process waits.

\courssub{Deadlock Detection and Recovery}

If the system neither prevents nor avoids deadlocks, it must \cle{detect} them after the fact:

\begin{itemize}
\item \textbf{Single instance per resource}: periodically search for a \textbf{cycle} in the RAG (or in its simplified form, the \emph{wait-for graph} where nodes are processes and an edge $P_i \to P_j$ means $P_i$ waits for a resource held by $P_j$).
\item \textbf{Multiple instances}: run a \textbf{detection algorithm} identical in spirit to the safety check, but using current \emph{requests} instead of maximum needs; any process that can never be marked finished is deadlocked.
\end{itemize}

Once a deadlock is detected, \cle{recovery} options are:

\begin{enumerate}
\item \textbf{Process termination}: kill one process at a time (or all deadlocked processes) until the cycle breaks --- choose victims by priority, work already done, resources held.
\item \textbf{Resource preemption}: forcibly take resources from a victim and give them to others; the victim is \textbf{rolled back} to a previous checkpoint and restarted later.
\end{enumerate}

\begin{savaistu}[The ostrich algorithm]
Many real operating systems (Windows, Linux) simply \textbf{ignore} deadlocks --- the so-called \emph{ostrich algorithm}. Deadlocks are rare, and prevention/avoidance mechanisms slow the system down permanently to guard against an event that may never happen. Engineers therefore accept a rare freeze (fixed by a reboot) rather than pay a constant performance cost. It is a deliberate engineering trade-off, not an oversight.
\end{savaistu}

\begin{popfigure}
\begin{center}
\begin{tikzpicture}[node distance=0.55cm, font=\small,
box/.style={draw=popblue, thick, rounded corners, align=center, text width=3.1cm, minimum height=0.9cm},
dec/.style={draw=poporange, thick, diamond, aspect=2.2, align=center, text width=2.6cm, inner sep=1pt},
arr/.style={->, thick, popdark}]
\node[box, fill=popblueL] (start) {Deadlock problem};
\node[dec, below=of start] (q) {Choose a strategy};
\node[box, below left=0.9cm and 2.6cm of q, fill=popgreenL, align=center] (prev){\textbf{Prevention}\\ negate one Coffman condition};
\node[box, below=1.1cm of q, fill=popgreenL, align=center] (avoid){\textbf{Avoidance}\\ Banker's algorithm, safe states};
\node[box, below right=0.9cm and 2.6cm of q, fill=popgreenL, align=center] (det){\textbf{Detection}\\ cycle search in RAG};
\node[box, below=of det, fill=popblueL, align=center] (rec){\textbf{Recovery}\\ kill, preempt, rollback};
\node[box, below=of avoid, fill=popblueL, align=center] (ostrich){\textbf{Ostrich}\\ ignore; reboot if it happens};
\draw[arr] (start) -- (q);
\draw[arr] (q) -| (prev);
\draw[arr] (q) -- (avoid);
\draw[arr] (q) -| (det);
\draw[arr] (det) -- (rec);
\draw[arr] (avoid) -- (ostrich);
\end{tikzpicture}
\end{center}
\legende{The four families of deadlock handling strategies.}
\end{popfigure}

\courssub{Integration Activity (Lesson 50): A Complete Case Study}

\begin{exoresolu}[Scheduling, synchronisation and deadlock combined]
A school's server runs three processes. $P_1$ prints report cards (needs the printer, then the scanner), $P_2$ scans archives (needs the scanner, then the printer), $P_3$ computes statistics (CPU only). Burst times: $P_1 = 6$ ms, $P_2 = 4$ ms, $P_3 = 2$ ms; all arrive at time $0$ in the order $P_1, P_2, P_3$.
\begin{enumerate}
\item Compute average waiting time under FCFS and under SJF (non-preemptive).
\item The printer and scanner are protected by semaphores \texttt{S1} and \texttt{S2}. $P_1$ executes \texttt{wait(S1)} then \texttt{wait(S2)}; $P_2$ executes \texttt{wait(S2)} then \texttt{wait(S1)}. What can happen?
\item Propose a prevention solution and justify it with the Coffman conditions.
\end{enumerate}
\tcblower
\textbf{1. Scheduling.} FCFS order: $P_1, P_2, P_3$. Waiting times: $P_1 = 0$, $P_2 = 6$, $P_3 = 6+4 = 10$; average $= (0+6+10)/3 \approx 5.33$ ms. SJF order: $P_3 (2)$, $P_2 (4)$, $P_1 (6)$. Waiting times: $P_3 = 0$, $P_2 = 2$, $P_1 = 2+4 = 6$; average $= (0+2+6)/3 \approx 2.67$ ms. SJF is better, as expected.

\textbf{2. Synchronisation.} If $P_1$ acquires \texttt{S1} (printer) and is then preempted, and $P_2$ acquires \texttt{S2} (scanner), each process blocks forever on its second \texttt{wait}: this is exactly the printer--scanner \textbf{deadlock} of the RAG above. The four Coffman conditions hold: mutual exclusion (one printer, one scanner), hold and wait, no preemption of semaphores, circular wait.

\textbf{3. Prevention.} Impose a \textbf{resource ordering}: every process must request \texttt{S1} before \texttt{S2}. Then $P_2$ must also lock the printer first; whichever process gets the printer first can complete, and the circular wait condition is destroyed, so deadlock becomes impossible. (Equivalently, both processes could request both semaphores atomically, breaking hold and wait --- at the cost of lower utilisation.)
\end{exoresolu}

\begin{exercice}[Deadlock analysis]
A system has one resource type with 4 instances and three processes with Max $= (3, 3, 2)$ and Allocation $= (1, 2, 1)$.
\begin{enumerate}
\item Compute the Need matrix and the Available vector.
\item Is the state safe? If yes, give a safe sequence.
\item $P_1$ now requests 1 more instance. Should the request be granted according to the Banker's algorithm?
\end{enumerate}
\end{exercice}

\begin{corrige}[Exercise 1]
\textbf{1.} Need $= \text{Max} - \text{Allocation} = (3-1,\; 3-2,\; 2-1) = (2, 1, 1)$. Available $= 4 - (1+2+1) = 0$.
\textbf{2.} Work $= 0$: no process has Need $\le 0$ (all needs are positive), so no process can be guaranteed to finish. The state is \textbf{unsafe} (not yet a deadlock, but one bad request away).
\textbf{3.} Pretend to grant: Allocation becomes $(2,2,1)$, Need becomes $(1,1,1)$, Available becomes $-1$ --- impossible; the request exceeds Available, so it cannot even be pretended. The request must be \textbf{refused / delayed}; granting it would be invalid, and the system is already unsafe.
\end{corrige}

\begin{aretenir}[Key points of Section 5]
\begin{itemize}
\item A \cle{deadlock} is a permanent mutual blocking of processes, each holding resources and waiting for others.
\item Deadlock requires \textbf{all four Coffman conditions}: mutual exclusion, hold and wait, no preemption, circular wait.
\item In a RAG, \textbf{no cycle $\Rightarrow$ no deadlock}; a cycle with one instance per resource $\Rightarrow$ certain deadlock; with several instances, a cycle is only a warning.
\item \textbf{Prevention} negates one condition (request-all-at-once, resource ordering); \textbf{avoidance} keeps the system in safe states (Banker's algorithm: Need $=$ Max $-$ Allocation, find Need $\le$ Work, release, repeat); \textbf{detection} finds cycles after the fact, with recovery by killing, preemption or rollback; the \textbf{ostrich algorithm} deliberately does nothing.
\item Every strategy trades safety against resource utilisation and performance.
\end{itemize}
\end{aretenir}

\newpage

\coursec{Integration activity}

\courssub{Aims of the activity}
This integration activity mobilises the core concepts of process management covered in Lessons 44–49: process scheduling (both non-preemptive and preemptive), process synchronisation using semaphores, and deadlock handling. By working on a single realistic scenario — a cyber-café server in Yaoundé — you will:
\begin{itemize}
    \item Compute and compare scheduling metrics (waiting time, turnaround time) for FCFS, SJF and Round Robin, and justify a policy choice for customer satisfaction.
    \item Design a synchronisation mechanism (semaphore) to protect a shared printer and explain the race condition that occurs without it.
    \item Model a deadlock situation with a Resource Allocation Graph, verify the four necessary conditions, and propose both a prevention and a recovery strategy.
    \item Present your findings as a professional recommendation to the café owner, demonstrating the ability to integrate theory into practice — a key skill for the GCE Advanced Level examination.
\end{itemize}

\courssub{Situation}
\textbf{Yaoundé Cyber-Café ``Digital Hub''} operates a single server that handles customer jobs. At 08:00 the server receives five jobs with the following characteristics (arrival time in minutes after 08:00, burst time in minutes):

\begin{center}
\begin{tabular}{|c|l|c|c|}
\hline
\rowcolor{popblueL}
\textbf{Job ID} & \textbf{Description} & \textbf{Arrival Time} & \textbf{Burst Time} \\
\hline
J1 & Printing a 20-page document & 0 & 8 \\
\hline
J2 & Scanning a batch of photos & 1 & 4 \\
\hline
J3 & Browsing session (video streaming) & 2 & 9 \\
\hline
J4 & Video rendering (short clip) & 3 & 5 \\
\hline
J5 & File download (large dataset) & 4 & 2 \\
\hline
\end{tabular}
\end{center}

The café has only \textbf{one printer} and \textbf{one scanner}, both connected to the server. Print jobs are spooled to the printer; scan jobs are sent to the scanner. The server runs a simple job scheduler. The owner wants to minimise average waiting time and turnaround time to keep customers happy, but also needs to ensure that print jobs never interleave (which would produce garbled pages) and that the system never deadlocks when customers simultaneously request the printer and scanner.

Later that morning, two customers (C1 and C2) arrive almost simultaneously:
\begin{itemize}
    \item C1 starts a print job (acquires the printer) and then requests the scanner to scan a receipt.
    \item C2 starts a scan job (acquires the scanner) and then requests the printer to print a confirmation.
\end{itemize}
Both requests are made while the other resource is still held.

\courssub{Tasks}
\begin{enumerate}
    \item \textbf{Scheduling analysis}
    \begin{enumerate}
        \item Draw the Gantt chart for each of the following scheduling policies: First-Come-First-Served (FCFS), Shortest Job First (SJF) \emph{non-preemptive}, and Round Robin with time quantum $q = 3$ minutes.
        \item For each policy, compute the \emph{waiting time} and \emph{turnaround time} of every job, then the average waiting time and average turnaround time.
        \item Based on the computed averages and the nature of the jobs (interactive vs. batch), recommend the best scheduling policy for customer satisfaction. Justify your choice.
    \end{enumerate}
    \item \textbf{Synchronisation of the printer}
    \begin{enumerate}
        \item Write clear pseudocode (using \texttt{wait(S)} and \texttt{signal(S)} operations on a binary semaphore) that protects the critical section where a print job sends data to the printer.
        \item Explain, in terms of a \emph{race condition}, what would happen if two print jobs entered the critical section concurrently (describe the garbled output).
    \end{enumerate}
    \item \textbf{Deadlock analysis}
    \begin{enumerate}
        \item Draw the Resource Allocation Graph (RAG) for the situation where C1 holds the printer and requests the scanner, while C2 holds the scanner and requests the printer.
        \item Verify that all four necessary conditions for deadlock (Mutual Exclusion, Hold and Wait, No Preemption, Circular Wait) are satisfied.
        \item Propose \textbf{one deadlock prevention} technique (choose from: eliminating Hold and Wait, allowing Preemption, or imposing a Resource Ordering) and explain how it would be applied in this café.
        \item Propose \textbf{one deadlock recovery} technique (choose from: Process Termination or Resource Preemption) and describe the steps the server would take to recover.
    \end{enumerate}
    \item \textbf{Presentation} Prepare a 3-minute oral summary (as if advising the café owner) that integrates your scheduling recommendation, synchronisation solution, and deadlock strategy.
\end{enumerate}

\begin{methode}[How to compute scheduling metrics]
For every job, once the Gantt chart is drawn:
\begin{enumerate}
    \item Read the \cle{completion time} $CT$: the instant at which the job finishes.
    \item Compute the \cle{turnaround time}: $TT = CT - \text{arrival time}$.
    \item Compute the \cle{waiting time}: $WT = TT - \text{burst time}$ (total time spent in the ready queue).
    \item Average over all jobs: $\overline{WT} = \dfrac{\sum WT_i}{n}$ and $\overline{TT} = \dfrac{\sum TT_i}{n}$.
\end{enumerate}
\end{methode}

\courssub{Model answer}

\textbf{1. Scheduling analysis}\par

\textbf{FCFS (First-Come-First-Served).} Jobs are served strictly in order of arrival: J1, J2, J3, J4, J5. Since J1 arrives at time 0 and runs for 8 minutes, all other jobs queue behind it.

\begin{popfigure}
\begin{tikzpicture}[scale=0.7, transform shape]
\draw[thick] (0,0) -- (28,0);
\foreach \x in {0,2,...,28} \draw (\x,0.1) -- (\x,-0.1) node[below,font=\tiny] {\x};
\fill[popblue!30] (0,0.2) rectangle (8,0.8) node[midway,font=\small] {J1};
\fill[popgreen!30] (8,0.2) rectangle (12,0.8) node[midway,font=\small] {J2};
\fill[poporange!30] (12,0.2) rectangle (21,0.8) node[midway,font=\small] {J3};
\fill[poppurple!30] (21,0.2) rectangle (26,0.8) node[midway,font=\small] {J4};
\fill[poppink!30] (26,0.2) rectangle (28,0.8) node[midway,font=\small] {J5};
\end{tikzpicture}
\legende{Gantt chart for FCFS}
\end{popfigure}

Completion times: $CT_{J1}=8$, $CT_{J2}=12$, $CT_{J3}=21$, $CT_{J4}=26$, $CT_{J5}=28$.

Turnaround times ($TT = CT - \text{arrival}$):
\begin{align*}
TT_{J1} &= 8 - 0 = 8, & TT_{J2} &= 12 - 1 = 11, & TT_{J3} &= 21 - 2 = 19, \\
TT_{J4} &= 26 - 3 = 23, & TT_{J5} &= 28 - 4 = 24. \\
\text{Average } TT &= \frac{8+11+19+23+24}{5} = \frac{85}{5} = 17.0 \text{ min}.
\end{align*}

Waiting times ($WT = TT - \text{burst}$):
\begin{align*}
WT_{J1} &= 8 - 8 = 0, & WT_{J2} &= 11 - 4 = 7, & WT_{J3} &= 19 - 9 = 10, \\
WT_{J4} &= 23 - 5 = 18, & WT_{J5} &= 24 - 2 = 22. \\
\text{Average } WT &= \frac{0+7+10+18+22}{5} = \frac{57}{5} = 11.4 \text{ min}.
\end{align*}

\textbf{SJF (Non-preemptive).} At each decision point, choose among the \emph{arrived} jobs the one with the shortest burst time; once started, a job runs to completion.
\begin{itemize}
    \item Time 0: only J1 has arrived $\Rightarrow$ J1 runs 0--8.
    \item Time 8: arrived jobs are J2 (4), J3 (9), J4 (5), J5 (2). Shortest is J5 $\Rightarrow$ runs 8--10.
    \item Time 10: remaining J2 (4), J4 (5), J3 (9). Shortest is J2 $\Rightarrow$ runs 10--14.
    \item Time 14: remaining J4 (5), J3 (9). Shortest is J4 $\Rightarrow$ runs 14--19.
    \item Time 19: only J3 remains $\Rightarrow$ runs 19--28.
\end{itemize}

\begin{popfigure}
\begin{tikzpicture}[scale=0.7, transform shape]
\draw[thick] (0,0) -- (28,0);
\foreach \x in {0,2,...,28} \draw (\x,0.1) -- (\x,-0.1) node[below,font=\tiny] {\x};
\fill[popblue!30] (0,0.2) rectangle (8,0.8) node[midway,font=\small] {J1};
\fill[poppink!30] (8,0.2) rectangle (10,0.8) node[midway,font=\small] {J5};
\fill[popgreen!30] (10,0.2) rectangle (14,0.8) node[midway,font=\small] {J2};
\fill[poppurple!30] (14,0.2) rectangle (19,0.8) node[midway,font=\small] {J4};
\fill[poporange!30] (19,0.2) rectangle (28,0.8) node[midway,font=\small] {J3};
\end{tikzpicture}
\legende{Gantt chart for SJF (non-preemptive)}
\end{popfigure}

Completion times: $CT_{J1}=8$, $CT_{J5}=10$, $CT_{J2}=14$, $CT_{J4}=19$, $CT_{J3}=28$.

Turnaround times:
\begin{align*}
TT_{J1} &= 8 - 0 = 8, & TT_{J2} &= 14 - 1 = 13, & TT_{J3} &= 28 - 2 = 26, \\
TT_{J4} &= 19 - 3 = 16, & TT_{J5} &= 10 - 4 = 6. \\
\text{Average } TT &= \frac{8+13+26+16+6}{5} = \frac{69}{5} = 13.8 \text{ min}.
\end{align*}

Waiting times:
\begin{align*}
WT_{J1} &= 0, & WT_{J2} &= 13 - 4 = 9, & WT_{J3} &= 26 - 9 = 17, \\
WT_{J4} &= 16 - 5 = 11, & WT_{J5} &= 6 - 2 = 4. \\
\text{Average } WT &= \frac{0+9+17+11+4}{5} = \frac{41}{5} = 8.2 \text{ min}.
\end{align*}

\textbf{Round Robin (q = 3).} The ready queue is managed in arrival order (FCFS); each job receives at most 3 minutes of CPU per turn, then goes to the back of the queue if not finished. New arrivals join the \emph{tail} of the queue. Careful trace:
\begin{itemize}
    \item Time 0: J1 runs 0--3 (remaining 5). Queue: J2, J3, J4, J1. (J5 arrives at time 4 and joins the tail.)
    \item Time 3: J2 runs 3--6 (remaining 1). Queue: J3, J4, J1, J5, J2.
    \item Time 6: J3 runs 6--9 (remaining 6). Queue: J4, J1, J5, J2, J3.
    \item Time 9: J4 runs 9--12 (remaining 2). Queue: J1, J5, J2, J3, J4.
    \item Time 12: J1 runs 12--15 (remaining 2). Queue: J5, J2, J3, J4, J1.
    \item Time 15: J5 runs 15--17 (burst 2 $<$ q, finishes). Queue: J2, J3, J4, J1.
    \item Time 17: J2 runs 17--18 (remaining 1, finishes). Queue: J3, J4, J1.
    \item Time 18: J3 runs 18--21 (remaining 3). Queue: J4, J1, J3.
    \item Time 21: J4 runs 21--23 (remaining 2, finishes). Queue: J1, J3.
    \item Time 23: J1 runs 23--25 (remaining 2, finishes). Queue: J3.
    \item Time 25: J3 runs 25--28 (finishes).
\end{itemize}

\begin{popfigure}
\begin{tikzpicture}[scale=0.6, transform shape]
\draw[thick] (0,0) -- (28,0);
\foreach \x in {0,2,...,28} \draw (\x,0.1) -- (\x,-0.1) node[below,font=\tiny] {\x};
\fill[popblue!30] (0,0.2) rectangle (3,0.8) node[midway,font=\small] {J1};
\fill[popgreen!30] (3,0.2) rectangle (6,0.8) node[midway,font=\small] {J2};
\fill[poporange!30] (6,0.2) rectangle (9,0.8) node[midway,font=\small] {J3};
\fill[poppurple!30] (9,0.2) rectangle (12,0.8) node[midway,font=\small] {J4};
\fill[popblue!30] (12,0.2) rectangle (15,0.8) node[midway,font=\small] {J1};
\fill[poppink!30] (15,0.2) rectangle (17,0.8) node[midway,font=\small] {J5};
\fill[popgreen!30] (17,0.2) rectangle (18,0.8) node[midway,font=\small] {J2};
\fill[poporange!30] (18,0.2) rectangle (21,0.8) node[midway,font=\small] {J3};
\fill[poppurple!30] (21,0.2) rectangle (23,0.8) node[midway,font=\small] {J4};
\fill[popblue!30] (23,0.2) rectangle (25,0.8) node[midway,font=\small] {J1};
\fill[poporange!30] (25,0.2) rectangle (28,0.8) node[midway,font=\small] {J3};
\end{tikzpicture}
\legende{Gantt chart for Round Robin (q = 3)}
\end{popfigure}

Completion times: $CT_{J1}=25$, $CT_{J2}=18$, $CT_{J3}=28$, $CT_{J4}=23$, $CT_{J5}=17$.

Turnaround times:
\begin{align*}
TT_{J1} &= 25 - 0 = 25, & TT_{J2} &= 18 - 1 = 17, & TT_{J3} &= 28 - 2 = 26, \\
TT_{J4} &= 23 - 3 = 20, & TT_{J5} &= 17 - 4 = 13. \\
\text{Average } TT &= \frac{25+17+26+20+13}{5} = \frac{101}{5} = 20.2 \text{ min}.
\end{align*}

Waiting times:
\begin{align*}
WT_{J1} &= 25 - 8 = 17, & WT_{J2} &= 17 - 4 = 13, & WT_{J3} &= 26 - 9 = 17, \\
WT_{J4} &= 20 - 5 = 15, & WT_{J5} &= 13 - 2 = 11. \\
\text{Average } WT &= \frac{17+13+17+15+11}{5} = \frac{73}{5} = 14.6 \text{ min}.
\end{align*}

\begin{attention}[Common mistakes in Round Robin traces]
\begin{itemize}
    \item Do not forget a job that arrives \emph{while} another is running: J5 arrives at time 4 and must join the tail of the queue.
    \item A job whose remaining burst is smaller than the quantum finishes \emph{before} the quantum expires (J5 runs only 2 minutes, J2 only 1 minute at the end).
    \item A new arrival goes to the \emph{tail} of the ready queue, never in front of jobs already waiting.
\end{itemize}
\end{attention}

\textbf{Comparison and recommendation.}\par

\begin{center}
\begin{tabular}{|l|c|c|}
\hline
\rowcolor{popblueL}
\textbf{Policy} & \textbf{Avg Waiting Time (min)} & \textbf{Avg Turnaround Time (min)} \\
\hline
FCFS & 11.4 & 17.0 \\
\hline
SJF (non-preemptive) & 8.2 & 13.8 \\
\hline
Round Robin (q = 3) & 14.6 & 20.2 \\
\hline
\end{tabular}
\end{center}

SJF gives the lowest average waiting time (8.2 min) and average turnaround time (13.8 min) — this is expected, since SJF is provably optimal for average waiting time among non-preemptive policies. However, SJF can cause \cle{starvation} of long jobs (J3 waits 17 minutes) and requires knowing burst times in advance. Round Robin guarantees fairness and good \emph{response time} for interactive jobs (browsing, scanning), but its averages are worse here because of frequent context switches. FCFS is simple but penalises short jobs stuck behind long ones (the \emph{convoy effect}: J5 waits 22 minutes behind J1).

\textbf{Recommendation:} since the owner's stated goal is to minimise average waiting and turnaround times, \textbf{SJF (non-preemptive)} is the best single policy for the batch-like jobs (printing, rendering, downloads). In practice, a hybrid is ideal: SJF for batch jobs and Round Robin with a small quantum for interactive sessions (browsing, scanning), so that interactive customers never feel the server is frozen. The owner should simply be warned that pure SJF may starve very long jobs.

\textbf{2. Synchronisation of the printer}\par

\textbf{Semaphore pseudocode.} A binary semaphore \texttt{printerSem}, initialised to 1 (printer free), guarantees mutual exclusion on the printer:
\begin{itemize}
    \item \texttt{var printerSem : semaphore := 1}  \emph{(1 = printer free, 0 = printer busy)}
    \item \textbf{Procedure} \texttt{PrintJob(jobData)}:
    \begin{enumerate}
        \item \texttt{wait(printerSem)}  \emph{(acquire the printer; block if busy)}
        \item \texttt{sendToPrinter(jobData)}  \emph{(critical section: one job at a time)}
        \item \texttt{signal(printerSem)}  \emph{(release the printer; wake a waiting job)}
    \end{enumerate}
\end{itemize}
Because \texttt{wait} and \texttt{signal} are \emph{atomic} (indivisible) operations, at most one process can be inside the critical section at any instant.

\textbf{Race condition without the semaphore.} If two print jobs (e.g., J1 and J4) enter \texttt{sendToPrinter} concurrently, their data streams interleave in the printer's buffer: a fragment of J1's page 1, then a fragment of J4's page 1, then J1 again, and so on. The printed pages are \emph{garbled} — mixed text and images from both documents — and unusable. This is a classic \cle{race condition}: the final result depends on the unpredictable relative timing of the processes accessing the shared resource without mutual exclusion. The semaphore removes the race by serialising access.

\textbf{3. Deadlock analysis}\par

\textbf{Resource Allocation Graph (RAG).} Circles represent processes, rectangles represent resources. An \emph{assignment edge} (resource $\to$ process) means the resource is held by the process; a \emph{request edge} (process $\to$ resource) means the process is waiting for the resource.
\begin{itemize}
    \item Assignment edges: Printer $\to$ C1, Scanner $\to$ C2.
    \item Request edges: C1 $\to$ Scanner, C2 $\to$ Printer.
\end{itemize}

\begin{popfigure}
\begin{tikzpicture}[node distance=2.5cm, >=stealth, thick]
\node[circle, draw, minimum size=1cm] (C1) {C1};
\node[circle, draw, minimum size=1cm, right of=C1] (C2) {C2};
\node[rectangle, draw, minimum size=1cm, above of=C1] (P) {Printer};
\node[rectangle, draw, minimum size=1cm, above of=C2] (S) {Scanner};
\draw[->] (P) -- (C1) node[midway, left, font=\small] {held};
\draw[->] (S) -- (C2) node[midway, right, font=\small] {held};
\draw[->] (C1) -- (S) node[midway, left, font=\small] {request};
\draw[->] (C2) -- (P) node[midway, right, font=\small] {request};
\end{tikzpicture}
\legende{Resource Allocation Graph: the cycle C1 $\to$ Scanner $\to$ C2 $\to$ Printer $\to$ C1 reveals the deadlock}
\end{popfigure}

The graph contains a cycle: C1 $\to$ Scanner $\to$ C2 $\to$ Printer $\to$ C1. Since each resource has a single instance, a cycle in the RAG is a \emph{sufficient} condition for deadlock: neither process can ever proceed.

\textbf{Verification of the four necessary conditions.}
\begin{enumerate}
    \item \textbf{Mutual Exclusion:} the printer and the scanner are non-shareable; only one job can use each at a time. Satisfied.
    \item \textbf{Hold and Wait:} C1 holds the printer while waiting for the scanner; C2 holds the scanner while waiting for the printer. Satisfied.
    \item \textbf{No Preemption:} resources cannot be forcibly taken away; they are released only voluntarily by the holder. Satisfied.
    \item \textbf{Circular Wait:} C1 waits for the scanner held by C2, and C2 waits for the printer held by C1 — a circular chain exists. Satisfied.
\end{enumerate}
All four conditions hold simultaneously, so the system is deadlocked.

\textbf{Deadlock prevention: Resource Ordering.} Impose a global ordering on resources, e.g., Printer $<$ Scanner, and require every process to request resources in increasing order only. In the café, any job needing both devices must request the printer \emph{first}, then the scanner. C1 (print then scan) already follows the order; C2 would be forced to request the printer first as well. Then a circular wait becomes impossible: a cycle would require a process to hold a high-numbered resource while waiting for a lower-numbered one, which the rule forbids. The Circular Wait condition is broken, so deadlock can never occur.

\textbf{Deadlock recovery: Process Termination.} The server detects the deadlock (e.g., via a timeout on blocked jobs or a periodic cycle-detection algorithm on the RAG). It then \textbf{terminates one of the deadlocked processes} — say C2, the scan job. The scanner is released and allocated to C1, which completes its work and releases both resources. C2's job is restarted later (the customer re-submits the scan request). The system recovers at the cost of losing C2's progress; the victim is usually chosen to minimise cost (youngest job, least work done, lowest priority).

\textbf{4. Presentation outline (3 minutes)}\par
\begin{itemize}
    \item \textbf{Opening (30 s):} ``Good morning. I have analysed the server workload of Digital Hub. Here are my recommendations to maximise customer satisfaction and system reliability.''
    \item \textbf{Scheduling (60 s):} show the comparison table. ``SJF gives the lowest average waiting time (8.2 min) and turnaround time (13.8 min). For your mix of jobs, I recommend SJF for background jobs such as printing and downloads, combined with Round Robin (q = 3) for interactive browsing and scanning sessions. This keeps averages low while guaranteeing responsiveness.''
    \item \textbf{Printer synchronisation (45 s):} ``Without protection, two simultaneous print jobs interleave and produce garbled pages. A binary semaphore around the printer driver enforces mutual exclusion — a simple, proven solution with negligible overhead.''
    \item \textbf{Deadlock strategy (45 s):} ``The Resource Allocation Graph shows a cycle: all four deadlock conditions hold. Prevention by resource ordering — always request the printer before the scanner — eliminates circular wait entirely. Should a deadlock ever occur, the server will terminate the younger job and restart it, minimising disruption.''
    \item \textbf{Closing (15 s):} ``These three measures will make Digital Hub faster, fairer and crash-free. Thank you.''
\end{itemize}

\begin{aretenir}[Key takeaways for the exam]
\begin{itemize}
    \item Always draw the Gantt chart first, then compute $TT = CT - \text{arrival}$ and $WT = TT - \text{burst}$; show the chart in your answer.
    \item SJF minimises average waiting time but can starve long processes; Round Robin guarantees fairness and good response time at the cost of more context switches.
    \item A binary semaphore (mutex) with atomic \texttt{wait}/\texttt{signal} operations solves the critical-section problem for a single shared resource.
    \item Deadlock requires all four conditions (Mutual Exclusion, Hold and Wait, No Preemption, Circular Wait); prevention attacks one condition, recovery breaks the deadlock after it occurs.
    \item In an integration question, link every answer back to the scenario (the café owner's goals) and justify each recommendation with computed values.
\end{itemize}
\end{aretenir}

\newpage

\coursec{Methods and worked examples}

\courssub{Method 1 — Drawing a Gantt Chart and Computing Performance Criteria}

\begin{methode}[Scheduling a Set of Processes with a Gantt Chart]
\begin{enumerate}
\item \textbf{Read the data table carefully.} For each process note its \cle{arrival time} (AT) and \cle{burst time} (BT, the CPU time it needs). Never confuse the two.
\item \textbf{Identify the algorithm} (FCFS, SJF, SRTF, Round Robin, Priority\ldots) and whether it is \cle{preemptive} or not.
\item \textbf{Build the ready queue step by step.} At each decision instant (an arrival, or the end of a burst/quantum), list the processes present and apply the algorithm's rule to choose the next one.
\item \textbf{Draw the Gantt chart}: a horizontal time axis, one block per execution interval, labelled with the process name, with the start and end times written below.
\item \textbf{Compute for each process}:
\begin{itemize}
\item \cle{Completion time} CT = time at which the process finishes;
\item \cle{Turnaround time} TAT $=$ CT $-$ AT;
\item \cle{Waiting time} WT $=$ TAT $-$ BT;
\item \cle{Response time} RT $=$ (first time the process gets the CPU) $-$ AT.
\end{itemize}
\item \textbf{Average the results} over all processes. Give answers as fractions or decimals with a point (e.g. $8.25$~ms).
\item \textbf{Check}: WT $=$ TAT $-$ BT must never be negative; the total length of the Gantt chart must equal the sum of burst times (if the CPU is never idle).
\end{enumerate}
\end{methode}

\begin{attention}[Standard Assumptions in GCE Scheduling Problems]
Unless the question says otherwise, assume that: the cost of a \cle{context switch} is negligible (0~ms); a process arriving at time $t$ is available for scheduling \emph{at} time $t$; and all times are integers. State these assumptions briefly if you use them.
\end{attention}

\begin{exoresolu}[FCFS Scheduling — Cyber Café in Buea]
The manager of a cyber café in Buea observes four print jobs arriving at the single printer server. Modelled as CPU processes, they are:

\begin{center}
\begin{tabular}{|c|c|c|}
\hline
\rowcolor{popblueL} Process & Arrival time (ms) & Burst time (ms) \\
\hline
$P_1$ & 0 & 8 \\
\hline
$P_2$ & 1 & 4 \\
\hline
$P_3$ & 2 & 9 \\
\hline
$P_4$ & 3 & 5 \\
\hline
\end{tabular}
\end{center}

Using \textbf{First-Come First-Served (FCFS)}, draw the Gantt chart and compute the average turnaround time and average waiting time.

\tcblower
\textbf{Step 1 — Order of execution.} FCFS serves processes in arrival order: $P_1$ (arrives at 0), then $P_2$ (1), then $P_3$ (2), then $P_4$ (3). Since FCFS is \emph{non-preemptive}, each process runs to completion once started.

\textbf{Step 2 — Gantt chart.}

\begin{popfigure}
\begin{center}
\begin{tikzpicture}[x=0.42cm,y=0.8cm]
\draw[fill=popblueL] (0,0) rectangle (8,1);
\draw[fill=popgreenL] (8,0) rectangle (12,1);
\draw[fill=popblueL!60] (12,0) rectangle (21,1);
\draw[fill=popgreenL!60] (21,0) rectangle (26,1);
\node at (4,0.5) {$P_1$};
\node at (10,0.5) {$P_2$};
\node at (16.5,0.5) {$P_3$};
\node at (23.5,0.5) {$P_4$};
\foreach \x in {0,8,12,21,26} \node[below] at (\x,0) {\x};
\end{tikzpicture}
\legende{Gantt chart for FCFS.}
\end{center}
\end{popfigure}

$P_1$ runs from 0 to 8; $P_2$ from 8 to 12; $P_3$ from 12 to 21; $P_4$ from 21 to 26.

\textbf{Step 3 — Criteria table.} Recall: TAT $=$ CT $-$ AT and WT $=$ TAT $-$ BT.

\begin{center}
\begin{tabular}{|c|c|c|c|c|c|}
\hline
\rowcolor{popblueL} Process & AT & BT & CT & TAT $=$ CT $-$ AT & WT $=$ TAT $-$ BT \\
\hline
$P_1$ & 0 & 8 & 8 & $8-0=8$ & $8-8=0$ \\
\hline
$P_2$ & 1 & 4 & 12 & $12-1=11$ & $11-4=7$ \\
\hline
$P_3$ & 2 & 9 & 21 & $21-2=19$ & $19-9=10$ \\
\hline
$P_4$ & 3 & 5 & 26 & $26-3=23$ & $23-5=18$ \\
\hline
\end{tabular}
\end{center}

\textbf{Step 4 — Averages.}
\[ \text{Average TAT} = \frac{8+11+19+23}{4} = \frac{61}{4} = 15.25\ \text{ms} \]
\[ \text{Average WT} = \frac{0+7+10+18}{4} = \frac{35}{4} = 8.75\ \text{ms} \]

\textbf{Conclusion.} With FCFS, the average turnaround time is $15.25$~ms and the average waiting time is $8.75$~ms. Note how the long job $P_3$ penalises $P_4$: this is the \cle{convoy effect}, a classic weakness of FCFS.
\end{exoresolu}

\courssub{Method 2 — Applying Non-preemptive SJF}

\begin{methode}[Shortest Job First (Non-preemptive)]
\begin{enumerate}
\item At the moment the CPU becomes free, look \textbf{only at the processes that have already arrived}.
\item Choose the one with the \textbf{smallest burst time}. In case of a tie, use the arrival order (FCFS tie-break).
\item Once started, a process runs \textbf{to completion} (non-preemptive): do not interrupt it even if a shorter job arrives meanwhile.
\item If the CPU is free but no process has arrived yet, leave it \textbf{idle} until the next arrival.
\item Draw the Gantt chart, then compute CT, TAT, WT as in Method 1.
\item \textbf{Reflex}: SJF minimises the average waiting time among non-preemptive algorithms — a result the examiner may ask you to state.
\end{enumerate}
\end{methode}

\begin{exoresolu}[SJF versus FCFS]
Using the same data as the previous exercise:

\begin{center}
\begin{tabular}{|c|c|c|}
\hline
\rowcolor{popblueL} Process & Arrival time (ms) & Burst time (ms) \\
\hline
$P_1$ & 0 & 8 \\
\hline
$P_2$ & 1 & 4 \\
\hline
$P_3$ & 2 & 9 \\
\hline
$P_4$ & 3 & 5 \\
\hline
\end{tabular}
\end{center}

Apply non-preemptive \textbf{Shortest Job First}, draw the Gantt chart and compute the average waiting time. Compare with FCFS.

\tcblower
\textbf{Step 1 — Decision at time 0.} Only $P_1$ has arrived, so $P_1$ runs from 0 to 8 (non-preemptive: it cannot be interrupted).

\textbf{Step 2 — Decision at time 8.} $P_2$ (BT 4), $P_3$ (BT 9) and $P_4$ (BT 5) are all waiting. The shortest is $P_2$ (4~ms): it runs from 8 to 12.

\textbf{Step 3 — Decision at time 12.} Remaining: $P_3$ (9) and $P_4$ (5). The shortest is $P_4$: it runs from 12 to 17.

\textbf{Step 4 — Decision at time 17.} Only $P_3$ remains: it runs from 17 to 26.

\begin{popfigure}
\begin{center}
\begin{tikzpicture}[x=0.42cm,y=0.8cm]
\draw[fill=popblueL] (0,0) rectangle (8,1);
\draw[fill=popgreenL] (8,0) rectangle (12,1);
\draw[fill=popgreenL!60] (12,0) rectangle (17,1);
\draw[fill=popblueL!60] (17,0) rectangle (26,1);
\node at (4,0.5) {$P_1$};
\node at (10,0.5) {$P_2$};
\node at (14.5,0.5) {$P_4$};
\node at (21.5,0.5) {$P_3$};
\foreach \x in {0,8,12,17,26} \node[below] at (\x,0) {\x};
\end{tikzpicture}
\legende{Gantt chart for non-preemptive SJF.}
\end{center}
\end{popfigure}

\textbf{Step 5 — Criteria table.}

\begin{center}
\begin{tabular}{|c|c|c|c|c|c|}
\hline
\rowcolor{popblueL} Process & AT & BT & CT & TAT & WT \\
\hline
$P_1$ & 0 & 8 & 8 & 8 & 0 \\
\hline
$P_2$ & 1 & 4 & 12 & 11 & 7 \\
\hline
$P_4$ & 3 & 5 & 17 & 14 & 9 \\
\hline
$P_3$ & 2 & 9 & 26 & 24 & 15 \\
\hline
\end{tabular}
\end{center}

\[ \text{Average WT} = \frac{0+7+9+15}{4} = \frac{31}{4} = 7.75\ \text{ms} \]

\textbf{Comparison.} SJF gives an average waiting time of $7.75$~ms against $8.75$~ms for FCFS: SJF is better, as expected, since it is provably optimal for average waiting time. However, the long process $P_3$ is served last: SJF can cause \cle{starvation} of long jobs.
\end{exoresolu}

\courssub{Method 3 — Preemptive Scheduling: SRTF}

\begin{methode}[Shortest Remaining Time First]
\begin{enumerate}
\item SRTF is the \textbf{preemptive} version of SJF: the CPU always goes to the process with the smallest \cle{remaining} burst time.
\item \textbf{Re-examine the situation at every arrival}: if the new process's burst time is smaller than the remaining time of the running process, \textbf{preempt} the running process and put it back in the ready queue with its updated remaining time.
\item Keep a small table of \textbf{remaining times}, updated at each event (arrival or completion).
\item On ties in remaining time, prefer the process that arrived first (or keep the running one to avoid useless context switches).
\item When computing the response time, remember it is measured at the \textbf{first} allocation of the CPU, which may be long before completion.
\end{enumerate}
\end{methode}

\begin{exoresolu}[SRTF with Preemption]
Consider the following processes:

\begin{center}
\begin{tabular}{|c|c|c|}
\hline
\rowcolor{popblueL} Process & Arrival time (ms) & Burst time (ms) \\
\hline
$P_1$ & 0 & 7 \\
\hline
$P_2$ & 2 & 4 \\
\hline
$P_3$ & 4 & 1 \\
\hline
$P_4$ & 5 & 4 \\
\hline
\end{tabular}
\end{center}

Schedule them with \textbf{SRTF}, draw the Gantt chart and compute the average turnaround time.

\tcblower
\textbf{Step 1 — Trace the events.}
\begin{itemize}
\item \textbf{$t=0$:} only $P_1$ is present; it starts. Remaining: $P_1 = 7$.
\item \textbf{$t=2$:} $P_2$ arrives (BT 4). $P_1$ has run 2~ms, so its remaining time is $7-2=5$. Since $4 < 5$, $P_2$ \textbf{preempts} $P_1$.
\item \textbf{$t=4$:} $P_3$ arrives (BT 1). $P_2$ has remaining $4-2=2$. Since $1 < 2$, $P_3$ preempts $P_2$.
\item \textbf{$t=5$:} $P_3$ finishes (it needed only 1~ms). $P_4$ arrives (BT 4). Ready queue: $P_2$ (remaining 2), $P_1$ (remaining 5), $P_4$ (4). The smallest is $P_2$: it resumes.
\item \textbf{$t=7$:} $P_2$ finishes. Remaining: $P_4$ (4) and $P_1$ (5). $P_4$ runs from 7 to 11.
\item \textbf{$t=11$:} $P_4$ finishes. $P_1$ resumes with 5~ms remaining, from 11 to 16.
\end{itemize}

\begin{popfigure}
\begin{center}
\begin{tikzpicture}[x=0.5cm,y=0.8cm]
\draw[fill=popblueL] (0,0) rectangle (2,1);
\draw[fill=popgreenL] (2,0) rectangle (4,1);
\draw[fill=popblueL!60] (4,0) rectangle (5,1);
\draw[fill=popgreenL] (5,0) rectangle (7,1);
\draw[fill=popgreenL!60] (7,0) rectangle (11,1);
\draw[fill=popblueL] (11,0) rectangle (16,1);
\node at (1,0.5) {$P_1$};
\node at (3,0.5) {$P_2$};
\node at (4.5,0.5) {$P_3$};
\node at (6,0.5) {$P_2$};
\node at (9,0.5) {$P_4$};
\node at (13.5,0.5) {$P_1$};
\foreach \x in {0,2,4,5,7,11,16} \node[below] at (\x,0) {\x};
\end{tikzpicture}
\legende{Gantt chart for SRTF: note the two preemptions.}
\end{center}
\end{popfigure}

\textbf{Step 2 — Criteria table.}

\begin{center}
\begin{tabular}{|c|c|c|c|c|c|}
\hline
\rowcolor{popblueL} Process & AT & BT & CT & TAT $=$ CT $-$ AT & WT $=$ TAT $-$ BT \\
\hline
$P_1$ & 0 & 7 & 16 & 16 & 9 \\
\hline
$P_2$ & 2 & 4 & 7 & 5 & 1 \\
\hline
$P_3$ & 4 & 1 & 5 & 1 & 0 \\
\hline
$P_4$ & 5 & 4 & 11 & 6 & 2 \\
\hline
\end{tabular}
\end{center}

\[ \text{Average TAT} = \frac{16+5+1+6}{4} = \frac{28}{4} = 7\ \text{ms} \]

\textbf{Conclusion.} SRTF gives an average turnaround time of $7$~ms. (For comparison, non-preemptive SJF on the same data gives an average TAT of $8$~ms: $P_1$ runs 0--7, then $P_3$ 7--8, $P_2$ 8--12, $P_4$ 12--16, so TAT $= 7, 10, 4, 11$.) The gain comes at the price of two context switches. Long processes like $P_1$ may be delayed repeatedly: the starvation risk is even stronger than with SJF.
\end{exoresolu}

\courssub{Method 4 — Round Robin Scheduling}

\begin{methode}[Round Robin with Quantum $q$]
\begin{enumerate}
\item Maintain a \textbf{FIFO ready queue}. New arrivals are added at the \textbf{tail} of the queue.
\item Give the CPU to the process at the head for at most \cle{quantum} $q$ time units.
\item If the process finishes before the quantum expires, remove it and take the next one immediately.
\item If the quantum expires first, \textbf{preempt} the process and put it back at the \textbf{tail} of the queue (after any process that arrived during the quantum — apply the convention stated in the question; if none is stated, add new arrivals first, then the preempted process).
\item Repeat until all processes complete. A process's remaining time decreases by $q$ each time it uses a full quantum.
\item \textbf{Reflex}: a very large $q$ makes RR behave like FCFS; a very small $q$ wastes time in context switches.
\end{enumerate}
\end{methode}

\begin{exoresolu}[Round Robin with $q = 3$ ms]
Three processes arrive at a single-CPU system:

\begin{center}
\begin{tabular}{|c|c|c|}
\hline
\rowcolor{popblueL} Process & Arrival time (ms) & Burst time (ms) \\
\hline
$P_1$ & 0 & 5 \\
\hline
$P_2$ & 1 & 9 \\
\hline
$P_3$ & 2 & 3 \\
\hline
\end{tabular}
\end{center}

Apply \textbf{Round Robin} with quantum $q = 3$~ms. Draw the Gantt chart and compute the average waiting time and the response time of each process.

\tcblower
\textbf{Step 1 — Simulate the ready queue.} Convention: when a quantum expires, newly arrived processes join the queue before the preempted process.

\begin{itemize}
\item $t=0$: queue $= [P_1]$. $P_1$ runs 0--3 (remaining 2). During this quantum $P_2$ (at 1) and $P_3$ (at 2) arrive. Queue $= [P_2, P_3, P_1]$.
\item $t=3$: $P_2$ runs 3--6 (remaining 6). Queue $= [P_3, P_1, P_2]$.
\item $t=6$: $P_3$ runs 6--9; it needs exactly 3, so it \textbf{finishes} at 9. Queue $= [P_1, P_2]$.
\item $t=9$: $P_1$ runs 9--11; it needs only 2, so it \textbf{finishes} at 11. Queue $= [P_2]$.
\item $t=11$: $P_2$ runs 11--14 (remaining 3). Queue $= [P_2]$.
\item $t=14$: $P_2$ runs 14--17 and \textbf{finishes}.
\end{itemize}

\begin{popfigure}
\begin{center}
\begin{tikzpicture}[x=0.5cm,y=0.8cm]
\draw[fill=popblueL] (0,0) rectangle (3,1);
\draw[fill=popgreenL] (3,0) rectangle (6,1);
\draw[fill=popblueL!60] (6,0) rectangle (9,1);
\draw[fill=popblueL] (9,0) rectangle (11,1);
\draw[fill=popgreenL] (11,0) rectangle (14,1);
\draw[fill=popgreenL] (14,0) rectangle (17,1);
\node at (1.5,0.5) {$P_1$};
\node at (4.5,0.5) {$P_2$};
\node at (7.5,0.5) {$P_3$};
\node at (10,0.5) {$P_1$};
\node at (12.5,0.5) {$P_2$};
\node at (15.5,0.5) {$P_2$};
\foreach \x in {0,3,6,9,11,14,17} \node[below] at (\x,0) {\x};
\end{tikzpicture}
\legende{Gantt chart for Round Robin, $q = 3$~ms.}
\end{center}
\end{popfigure}

\textbf{Step 2 — Criteria table.} Response time RT $=$ first CPU allocation $-$ AT.

\begin{center}
\begin{tabular}{|c|c|c|c|c|c|c|}
\hline
\rowcolor{popblueL} Process & AT & BT & CT & TAT & WT & RT \\
\hline
$P_1$ & 0 & 5 & 11 & 11 & 6 & 0 \\
\hline
$P_2$ & 1 & 9 & 17 & 16 & 7 & 2 \\
\hline
$P_3$ & 2 & 3 & 9 & 7 & 4 & 4 \\
\hline
\end{tabular}
\end{center}

For example, for $P_1$: WT $=$ TAT $-$ BT $= 11 - 5 = 6$~ms (it waited from 3 to 9, i.e. 6~ms — consistent). For $P_2$: RT $= 3 - 1 = 2$~ms (first run starts at $t=3$); for $P_3$: RT $= 6 - 2 = 4$~ms.

\[ \text{Average WT} = \frac{6+7+4}{3} = \frac{17}{3} \approx 5.67\ \text{ms} \]

\textbf{Conclusion.} Round Robin guarantees every process a share of the CPU within a bounded delay (good response times, no starvation), which is why it is the basis of time-sharing systems. The cost is more context switches than FCFS or SJF.
\end{exoresolu}

\courssub{Method 5 — Priority Scheduling}

\begin{methode}[Priority Scheduling (Preemptive or Not)]
\begin{enumerate}
\item Check the \textbf{convention}: in most GCE problems, the \textbf{smallest number} means the \textbf{highest} priority — but always read the statement.
\item If \textbf{non-preemptive}: when the CPU is free, pick the waiting process with the highest priority; let it finish.
\item If \textbf{preemptive}: at every arrival, compare priorities; a newly arrived process with a higher priority takes the CPU immediately.
\item Break ties with FCFS (earliest arrival first).
\item Mention the danger: \cle{starvation} of low-priority processes, and its remedy, \cle{ageing} (gradually raising the priority of processes that wait long).
\end{enumerate}
\end{methode}

\begin{exoresolu}[Preemptive Priority Scheduling]
A hospital management system in Yaoundé handles four tasks (1 = highest priority):

\begin{center}
\begin{tabular}{|c|c|c|c|}
\hline
\rowcolor{popblueL} Process & Arrival (ms) & Burst (ms) & Priority \\
\hline
$P_1$ & 0 & 10 & 3 \\
\hline
$P_2$ & 1 & 1 & 1 \\
\hline
$P_3$ & 2 & 2 & 4 \\
\hline
$P_4$ & 3 & 1 & 2 \\
\hline
\end{tabular}
\end{center}

Apply \textbf{preemptive priority scheduling}, draw the Gantt chart and compute the average turnaround time.

\tcblower
\textbf{Step 1 — Trace the events.}
\begin{itemize}
\item $t=0$: only $P_1$ (priority 3) is present; it starts.
\item $t=1$: $P_2$ arrives with priority 1 (highest). It \textbf{preempts} $P_1$ (remaining 9). $P_2$ runs 1--2 and finishes.
\item $t=2$: $P_3$ arrives (priority 4, lower than $P_1$'s 3). $P_1$ resumes.
\item $t=3$: $P_4$ arrives (priority 2, higher than $P_1$'s 3). $P_4$ \textbf{preempts} $P_1$ (remaining 8). $P_4$ runs 3--4 and finishes.
\item $t=4$: ready queue: $P_1$ (priority 3, remaining 8) and $P_3$ (priority 4). $P_1$ has the higher priority: it runs 4--12 and finishes.
\item $t=12$: $P_3$ runs 12--14 and finishes.
\end{itemize}

\begin{popfigure}
\begin{center}
\begin{tikzpicture}[x=0.55cm,y=0.8cm]
\draw[fill=popblueL] (0,0) rectangle (1,1);
\draw[fill=popgreenL] (1,0) rectangle (2,1);
\draw[fill=popblueL] (2,0) rectangle (3,1);
\draw[fill=popgreenL!60] (3,0) rectangle (4,1);
\draw[fill=popblueL] (4,0) rectangle (12,1);
\draw[fill=popblueL!60] (12,0) rectangle (14,1);
\node at (0.5,0.5) {$P_1$};
\node at (1.5,0.5) {$P_2$};
\node at (2.5,0.5) {$P_1$};
\node at (3.5,0.5) {$P_4$};
\node at (8,0.5) {$P_1$};
\node at (13,0.5) {$P_3$};
\foreach \x in {0,1,2,3,4,12,14} \node[below] at (\x,0) {\x};
\end{tikzpicture}
\legende{Gantt chart for preemptive priority scheduling.}
\end{center}
\end{popfigure}

\textbf{Step 2 — Criteria table.}

\begin{center}
\begin{tabular}{|c|c|c|c|c|}
\hline
\rowcolor{popblueL} Process & AT & BT & CT & TAT $=$ CT $-$ AT \\
\hline
$P_1$ & 0 & 10 & 12 & 12 \\
\hline
$P_2$ & 1 & 1 & 2 & 1 \\
\hline
$P_3$ & 2 & 2 & 14 & 12 \\
\hline
$P_4$ & 3 & 1 & 4 & 1 \\
\hline
\end{tabular}
\end{center}

\[ \text{Average TAT} = \frac{12+1+12+1}{4} = \frac{26}{4} = 6.5\ \text{ms} \]

\textbf{Conclusion.} Urgent tasks ($P_2$, $P_4$) are served almost immediately — exactly what a hospital system needs. But $P_3$, with the lowest priority, waits 10~ms before starting (it arrives at $t=2$ and first runs at $t=12$): without \cle{ageing}, such a process could starve indefinitely if high-priority tasks keep arriving.
\end{exoresolu}

\courssub{Method 6 — Synchronising Processes with Semaphores}

\begin{methode}[Solving a Synchronisation Problem with Semaphores]
\begin{enumerate}
\item \textbf{Identify the critical sections}: the code regions where shared resources (variables, files, printer) are accessed and which must not be executed by two processes at once.
\item \textbf{Choose the semaphores}: a \cle{binary semaphore} (or mutex) initialised to 1 for mutual exclusion; a \cle{counting semaphore} initialised to $n$ for a resource with $n$ identical instances (e.g. 3 printers).
\item \textbf{Protect} each critical section: \texttt{wait(S)} (or \texttt{P(S)}) before entering, \texttt{signal(S)} (or \texttt{V(S)}) after leaving.
\item Recall the semantics: \texttt{wait(S)} decrements $S$; if $S < 0$ the process \textbf{blocks}. \texttt{signal(S)} increments $S$; if processes are blocked, one is \textbf{woken up}.
\item \textbf{Check the invariants}: the semaphore's initial value must equal the number of processes allowed simultaneously in the critical section; every \texttt{wait} must be matched by a \texttt{signal} on \emph{every} execution path.
\item For producer--consumer problems, use two counting semaphores (\texttt{empty} $= N$, \texttt{full} $= 0$) plus a mutex for the buffer itself.
\end{enumerate}
\end{methode}

\begin{exoresolu}[Mutual Exclusion in a School Fees System]
Two clerks of a college in Bamenda simultaneously update the same shared variable \texttt{balance} (school fees account), each executing: read \texttt{balance}, add the payment, write \texttt{balance}.

\textbf{1)} Explain what can go wrong without synchronisation. \textbf{2)} Write pseudocode using a semaphore that guarantees a correct update. \textbf{3)} Trace the semaphore values if both processes try to enter at the same time.

\tcblower
\textbf{1) The race condition.} Suppose \texttt{balance} $= 50\,000$ FCFA, clerk A adds $10\,000$ and clerk B adds $25\,000$. Without protection the operations may interleave:
\begin{itemize}
\item A reads 50000; B reads 50000;
\item A computes $50000 + 10000 = 60000$ and writes 60000;
\item B computes $50000 + 25000 = 75000$ and writes 75000.
\end{itemize}
Final value: $75\,000$ instead of $85\,000$ — A's update is \textbf{lost}. This is a \cle{race condition}: the result depends on the unpredictable order of execution. The three instructions form a \cle{critical section}.

\textbf{2) Solution with a binary semaphore.} Declare \texttt{mutex}, a binary semaphore initialised to 1. Each clerk executes:

\begin{itemize}
\item[] \texttt{wait(mutex);}
\item[] \texttt{read balance;}
\item[] \texttt{balance := balance + payment;}
\item[] \texttt{write balance;}
\item[] \texttt{signal(mutex);}
\end{itemize}

The \texttt{wait}/\texttt{signal} pair guarantees \cle{mutual exclusion}: at most one process is inside the critical section at any time.

\textbf{3) Trace of the semaphore.} Recall: \texttt{wait} decrements and blocks the caller if the value becomes negative; \texttt{signal} increments and wakes a blocked process if any.

\begin{center}
\begin{tabular}{|c|l|c|l|}
\hline
\rowcolor{popblueL} Step & Action & \texttt{mutex} & Effect \\
\hline
1 & Initial state & 1 & critical section free \\
\hline
2 & A executes \texttt{wait(mutex)} & 0 & A enters the section \\
\hline
3 & B executes \texttt{wait(mutex)} & $-1$ & B \textbf{blocks} \\
\hline
4 & A updates and executes \texttt{signal} & 0 & B is woken up, enters \\
\hline
5 & B executes \texttt{signal(mutex)} & 1 & section free again \\
\hline
\end{tabular}
\end{center}

Now the updates happen one after the other: A brings the balance to $60\,000$, then B reads $60\,000$ and writes $85\,000$. The final value is correct.

\textbf{Remark.} Forgetting the \texttt{signal} (for example on an error path) would leave \texttt{mutex} at 0 forever: every other process would block indefinitely — a classic programming error that leads to deadlock.
\end{exoresolu}

\courssub{Method 7 — Detecting and Handling Deadlock}

\begin{methode}[Analysing a Deadlock Situation]
\begin{enumerate}
\item Recall the \textbf{four necessary conditions} (Coffman): \cle{mutual exclusion}, \cle{hold and wait}, \cle{no preemption}, \cle{circular wait}. A deadlock requires \emph{all four}.
\item Model the situation with a \cle{resource allocation graph} (RAG): circles for processes, rectangles for resources; an arrow resource $\to$ process means ``allocated to'', an arrow process $\to$ resource means ``requests''.
\item \textbf{Detect}: with one instance per resource, a deadlock exists if and only if the graph contains a \textbf{cycle}.
\item \textbf{Prevent}: design the system so that at least one condition can never hold (e.g. impose a global ordering of resources to break circular wait).
\item \textbf{Avoid}: before granting a request, check with the \cle{Banker's algorithm} that the resulting state is \emph{safe} (an order exists in which all processes can finish).
\item \textbf{Recover}: if deadlock occurs, kill a victim process or preempt its resources, then roll back.
\end{enumerate}
\end{methode}

\begin{exoresolu}[Deadlock Detection with a Resource Graph]
A system has two processes $P_1$, $P_2$ and two single-instance resources $R_1$ (printer) and $R_2$ (scanner). The current state is: $R_1$ is allocated to $P_1$; $R_2$ is allocated to $P_2$; $P_1$ requests $R_2$; $P_2$ requests $R_1$.

\textbf{1)} Draw the resource allocation graph. \textbf{2)} Is the system deadlocked? Justify using the four conditions. \textbf{3)} Propose one prevention strategy and one recovery action.

\tcblower
\textbf{1) Resource allocation graph.}

\begin{popfigure}
\begin{center}
\begin{tikzpicture}[scale=1.1, every node/.style={font=\small}]
\node[circle, draw=popblue, thick, minimum size=0.9cm] (p1) at (0,2) {$P_1$};
\node[circle, draw=popblue, thick, minimum size=0.9cm] (p2) at (4,2) {$P_2$};
\node[rectangle, draw=popgreen, thick, minimum width=1cm, minimum height=0.8cm] (r1) at (4,0) {$R_1$};
\node[rectangle, draw=popgreen, thick, minimum width=1cm, minimum height=0.8cm] (r2) at (0,0) {$R_2$};
\draw[->, thick, popdark] (r1) -- (p2) node[midway, right] {\scriptsize allocated};
\draw[->, thick, popdark] (r2) -- (p1) node[midway, left] {\scriptsize allocated};
\draw[->, thick, poporange] (p1) -- (r1) node[midway, above] {\scriptsize requests};
\draw[->, thick, poporange] (p2) -- (r2) node[midway, above] {\scriptsize requests};
\end{tikzpicture}
\legende{Resource allocation graph showing a cycle.}
\end{center}
\end{popfigure}

\textbf{2) Deadlock analysis.} The graph contains the cycle
\[ P_1 \to R_1 \to P_2 \to R_2 \to P_1. \]
Since each resource has a single instance, a cycle is a \textbf{sufficient} condition for deadlock. Let us verify the four Coffman conditions:
\begin{itemize}
\item \textbf{Mutual exclusion}: a printer or scanner can serve only one process at a time. \checkmark
\item \textbf{Hold and wait}: $P_1$ holds $R_1$ while waiting for $R_2$; $P_2$ holds $R_2$ while waiting for $R_1$. \checkmark
\item \textbf{No preemption}: the OS cannot forcefully take the printer from $P_1$. \checkmark
\item \textbf{Circular wait}: the cycle above shows each process waits for a resource held by the other. \checkmark
\end{itemize}
All four conditions hold: the system \textbf{is deadlocked}. Neither process can ever proceed.

\textbf{3) Handling strategies.}
\begin{itemize}
\item \textbf{Prevention} (break circular wait): impose a global order on resources, e.g. every process must request the printer ($R_1$) \emph{before} the scanner ($R_2$). Then $P_2$ could not hold $R_2$ while requesting $R_1$, and the cycle becomes impossible.
\item \textbf{Recovery}: choose a victim, e.g. abort $P_2$ (or preempt $R_2$ from it). $P_1$ then obtains $R_2$, finishes and releases $R_1$; $P_2$ is restarted later. The victim is usually chosen to minimise cost (lowest priority, least work done).
\end{itemize}
\end{exoresolu}

\courssub{Method 8 — Integration: Comparing Algorithms and Choosing One}

\begin{methode}[Choosing and Justifying a Scheduling Policy (Integration)]
\begin{enumerate}
\item \textbf{Compute} the average waiting and turnaround times for each candidate algorithm on the given data — never compare algorithms from intuition alone.
\item \textbf{Match the algorithm to the context}: interactive systems need good response time (Round Robin); batch systems need throughput (SJF); urgent tasks need priorities; real-time systems need deadlines guarantees.
\item \textbf{Discuss side effects}: starvation (SJF, priority), convoy effect (FCFS), context-switch overhead (small quantum in RR).
\item \textbf{Conclude in writing}: state the chosen algorithm, the numerical evidence, and one limitation — this three-part argument is what GCE examiners reward.
\end{enumerate}
\end{methode}

\begin{exoresolu}[Integration Activity — The GCE Board Results Server]
The results-processing server of an examination board receives three batch jobs at the same instant $t = 0$:

\begin{center}
\begin{tabular}{|c|c|}
\hline
\rowcolor{popblueL} Job & Burst time (s) \\
\hline
$J_1$ & 24 \\
\hline
$J_2$ & 3 \\
\hline
$J_3$ & 3 \\
\hline
\end{tabular}
\end{center}

\textbf{1)} Compute the average waiting time under FCFS (order $J_1, J_2, J_3$) and under SJF. \textbf{2)} The board also runs an interactive terminal service on the same machine. Which algorithm would you recommend for the combined workload, and why? \textbf{3)} State one risk of your recommendation and how to limit it.

\tcblower
\textbf{1) Numerical comparison.} All jobs arrive at $t=0$, so WT of the first job is 0 and each following job waits for the sum of the preceding bursts.

\emph{FCFS, order $J_1, J_2, J_3$:}
\begin{itemize}
\item $J_1$: runs 0--24, WT $= 0$;
\item $J_2$: runs 24--27, WT $= 24$;
\item $J_3$: runs 27--30, WT $= 27$.
\end{itemize}
\[ \text{Average WT (FCFS)} = \frac{0+24+27}{3} = \frac{51}{3} = 17\ \text{s}. \]

\emph{SJF:} shortest first: $J_2$ and $J_3$ (3~s each, tie broken by arrival order), then $J_1$.
\begin{itemize}
\item $J_2$: runs 0--3, WT $= 0$;
\item $J_3$: runs 3--6, WT $= 3$;
\item $J_1$: runs 6--30, WT $= 6$.
\end{itemize}
\[ \text{Average WT (SJF)} = \frac{0+3+6}{3} = \frac{9}{3} = 3\ \text{s}. \]

SJF reduces the average waiting time from 17~s to 3~s — a dramatic improvement, caused by eliminating the \cle{convoy effect} (short jobs stuck behind the 24~s job).

\textbf{2) Recommendation.} For a \emph{combined} workload (batch jobs \emph{plus} interactive terminals), I recommend \textbf{Round Robin} (or a multilevel feedback queue). Justification:
\begin{itemize}
\item Interactive users need a \textbf{response time} of a fraction of a second; under SJF or FCFS, a terminal user arriving behind the 24~s job would wait up to 24~s before the first response — unacceptable.
\item Round Robin gives every process the CPU within at most $(n-1)\times q$ time units, so terminals stay responsive while batch jobs progress steadily.
\item Numerically, SJF is optimal for average waiting time, but it optimises the wrong criterion for interactive work.
\end{itemize}

\textbf{3) Risk and mitigation.} Round Robin's risk is \textbf{context-switch overhead}: if the quantum is too small (say $q = 1$~ms for jobs measured in seconds), the CPU spends a significant fraction of its time switching instead of computing. The remedy is to choose a quantum large compared with the switch cost but small compared with typical interactive bursts (commonly 10--100~ms), or to use a multilevel feedback queue that gives short quanta to interactive processes and long ones to batch jobs.

\textbf{Final argument (as expected at the GCE).} SJF is numerically the best for pure batch (3~s vs 17~s average wait), but the presence of interactive users makes response time the dominant criterion; therefore Round Robin with a well-chosen quantum is the appropriate policy, its overhead being controlled by the choice of $q$.
\end{exoresolu}

\newpage

\coursec{Exercises}

\courssub{I check my knowledge}

\begin{exercice}[Multiple choice questions]
For each question, choose the single correct answer.
\begin{enumerate}
\item A \cle{process} is:
\begin{enumerate}
\item a program stored on the hard disk;
\item a program in execution, with its own state;
\item a hardware component of the CPU;
\item a file containing source code.
\end{enumerate}
\item The information about a process (state, program counter, registers, open files) is stored by the operating system in:
\begin{enumerate}
\item the cache memory;
\item the \cle{Process Control Block} (PCB);
\item the page table only;
\item the boot sector.
\end{enumerate}
\item Which scheduling algorithm can cause \cle{starvation} of long processes?
\begin{enumerate}
\item FCFS;
\item Round Robin;
\item Shortest Job First (SJF);
\item none of them.
\end{enumerate}
\item A deadlock can occur only if the following number of Coffman conditions hold simultaneously:
\begin{enumerate}
\item one; \item two; \item three; \item four.
\end{enumerate}
\end{enumerate}
\end{exercice}

\begin{exercice}[True or false — justify]
Say whether each statement is true or false, and justify your answer in one or two sentences.
\begin{enumerate}
\item In a non-preemptive scheduler, a running process can be forced to leave the CPU before it terminates or blocks.
\item A context switch consumes CPU time during which no useful user work is done.
\item A binary semaphore initialised to $1$ can be used to guarantee mutual exclusion on a critical section.
\item If a Resource Allocation Graph contains a cycle, then a deadlock necessarily exists, whatever the number of instances of each resource.
\end{enumerate}
\end{exercice}

\begin{exercice}[Definitions]
Give a precise definition of each of the following terms, in your own words:
\begin{enumerate}
\item \cle{turnaround time} of a process;
\item \cle{waiting time} of a process;
\item \cle{response time} of a process;
\item \cle{race condition};
\item \cle{deadlock}.
\end{enumerate}
\end{exercice}

\begin{exercice}[Direct calculation]
Three processes arrive at time $0$ in the order $P_1, P_2, P_3$, with burst times $8$, $4$ and $2$ ms respectively.
\begin{enumerate}
\item Compute the waiting time and the turnaround time of each process under FCFS scheduling.
\item Compute the average turnaround time.
\item Repeat the two previous questions assuming the burst times are executed in the order $P_3, P_2, P_1$. What do you observe?
\end{enumerate}
\end{exercice}

\courssub{I practise}

\begin{exercice}[FCFS and SJF Gantt charts]
Four processes are submitted to a single-CPU system:
\begin{center}
\begin{tabular}{|c|c|c|}
\hline
\rowcolor{popblueL} Process & Arrival time (ms) & Burst time (ms) \\
\hline
$P_1$ & 0 & 7 \\
\hline
$P_2$ & 2 & 4 \\
\hline
$P_3$ & 4 & 1 \\
\hline
$P_4$ & 5 & 4 \\
\hline
\end{tabular}
\end{center}
\begin{enumerate}
\item Draw the Gantt chart for FCFS scheduling and compute the average waiting time.
\item Draw the Gantt chart for non-preemptive SJF scheduling and compute the average waiting time.
\item Which algorithm performs better on this example? Justify.
\end{enumerate}
\end{exercice}

\begin{exercice}[Round Robin]
Five processes arrive at time $0$ in the order $P_1, \dots, P_5$ with burst times $10$, $3$, $6$, $8$ and $5$ ms. The system uses Round Robin with a quantum $q = 4$ ms.
\begin{enumerate}
\item Draw the Gantt chart.
\item Compute the response time of each process.
\item Compute the average turnaround time.
\item Without redoing the whole computation, discuss the effect of halving the quantum to $q = 2$ ms on the response time and on the overhead due to context switches.
\end{enumerate}
\end{exercice}

\begin{exercice}[The lost increment]
Two processes $A$ and $B$ share a variable \texttt{count} initialised to $0$. Each executes, once, the high-level instruction \texttt{count := count + 1}, which the machine carries out in three steps: \texttt{LOAD} (read \texttt{count} into a register), \texttt{ADD 1}, \texttt{STORE} (write the register back).
\begin{enumerate}
\item Exhibit an interleaving of the six machine instructions for which the final value of \texttt{count} is $1$ instead of $2$.
\item Name the phenomenon observed.
\item Correct the situation using a binary semaphore \texttt{mutex}: give its initial value and the pseudocode of each process.
\end{enumerate}
\end{exercice}

\begin{exercice}[Producer–consumer with a bounded buffer]
A producer process and a consumer process share a circular buffer of $5$ slots. The producer must never write into a full buffer and the consumer must never read from an empty buffer.
\begin{enumerate}
\item Give the semaphores needed and their initial values.
\item Write the complete pseudocode of the producer and of the consumer using the operations \texttt{wait} and \texttt{signal}.
\item Explain why the order of the two \texttt{wait} operations in each process is critical: what happens if the mutual-exclusion semaphore is taken before the counting semaphore?
\end{enumerate}
\end{exercice}

\courssub{I prepare for the GCE}

\begin{exercice}[Scheduling comparison — GCE style]
A computer centre in Yaoundé runs a batch system with one CPU. Four jobs are submitted as follows:
\begin{center}
\begin{tabular}{|c|c|c|}
\hline
\rowcolor{popblueL} Job & Arrival time (ms) & Burst time (ms) \\
\hline
$J_1$ & 0 & 8 \\
\hline
$J_2$ & 1 & 4 \\
\hline
$J_3$ & 2 & 9 \\
\hline
$J_4$ & 3 & 5 \\
\hline
\end{tabular}
\end{center}
\begin{enumerate}
\item Define the terms \emph{waiting time} and \emph{turnaround time} of a process. \hfill (2 marks)
\item Draw the Gantt chart and compute the average turnaround time and the average waiting time for each of the following algorithms:
\begin{enumerate}
\item FCFS; \hfill (4 marks)
\item non-preemptive SJF; \hfill (4 marks)
\item preemptive SJF (Shortest Remaining Time First). \hfill (5 marks)
\end{enumerate}
\item State which algorithm gives the best average waiting time on this example, and justify why this algorithm is optimal with respect to average waiting time. \hfill (2 marks)
\item Give one reason why this optimal algorithm is difficult to use in practice. \hfill (1 mark)
\item The centre decides to switch to Round Robin with quantum $q = 3$ ms. Draw the new Gantt chart and compute the response time of each job. \hfill (4 marks)
\end{enumerate}
\end{exercice}

\begin{exercice}[Banker's algorithm and Resource Allocation Graph — GCE style]
A system has three resource types $R_1, R_2, R_3$ and four processes. The current state is:
\begin{center}
\begin{tabular}{|c|ccc|ccc|}
\hline
\rowcolor{popblueL} & \multicolumn{3}{c|}{Allocation} & \multicolumn{3}{c|}{Max} \\
\rowcolor{popblueL} Process & $R_1$ & $R_2$ & $R_3$ & $R_1$ & $R_2$ & $R_3$ \\
\hline
$P_1$ & 0 & 1 & 0 & 7 & 5 & 3 \\
\hline
$P_2$ & 2 & 0 & 0 & 3 & 2 & 2 \\
\hline
$P_3$ & 3 & 0 & 2 & 9 & 0 & 2 \\
\hline
$P_4$ & 2 & 1 & 1 & 2 & 2 & 2 \\
\hline
\end{tabular}
\end{center}
The vector of available resources is $\text{Available} = (3, 3, 2)$.
\begin{enumerate}
\item Compute the \emph{Need} matrix. \hfill (2 marks)
\item Apply the Banker's safety algorithm to determine whether the system is in a safe state. If it is, exhibit a safe sequence; if it is not, justify carefully. \hfill (6 marks)
\item $P_2$ now requests one additional instance of $R_1$ and one of $R_2$. Using the Banker's algorithm, decide whether this request can be granted immediately. \hfill (4 marks)
\item State the four Coffman conditions necessary for a deadlock. \hfill (4 marks)
\item In another system, $P_1$ holds one instance of $R_1$ and requests $R_2$; $P_2$ holds one instance of $R_2$ and requests $R_3$; $P_3$ holds one instance of $R_3$ and requests $R_1$. Each resource type has exactly one instance. Draw the Resource Allocation Graph, state whether a deadlock exists, and justify your answer using the Coffman conditions. \hfill (4 marks)
\end{enumerate}
\end{exercice}

\begin{exercice}[Structured essay — processes, synchronisation and scheduling trade-offs]
\begin{enumerate}
\item Define the following terms: \emph{process}, \emph{Process Control Block}, \emph{context switch}. \hfill (3 marks)
\item Explain why a context switch is considered pure overhead, and state two consequences of this overhead on the design of a scheduling policy. \hfill (3 marks)
\item A bank in Douala equips its branches with interactive customer terminals connected to a central server. Discuss the trade-offs between preemptive and non-preemptive scheduling for such a system. Your discussion should address: responsiveness to interactive users, overhead, starvation, and predictability for long batch jobs such as end-of-day account processing. Conclude by recommending a scheduling discipline and justifying your choice. \hfill (8 marks)
\item The server runs two cooperating processes that both update a shared counter recording the number of transactions of the day. Explain, with an example of instruction interleaving, how the final value of the counter can be wrong, and show how a binary semaphore solves the problem. \hfill (6 marks)
\end{enumerate}
\end{exercice}

\newpage

\coursec{Solutions to the exercises}

\begin{corrige}[Exercise 1 — Multiple choice questions]
\begin{enumerate}
\item \textbf{Answer (b).} A process is a \emph{program in execution}, together with its current state (program counter, registers, stack, open files). A program on disk (a) is a passive entity; it becomes a process only when loaded and executed.
\item \textbf{Answer (b).} All the information the operating system needs about a process is grouped in the \cle{Process Control Block} (PCB), one per process. The PCB is what makes a context switch possible: the state of the suspended process is saved in its PCB and restored later.
\item \textbf{Answer (c).} Under \cle{SJF}, a long job can be postponed indefinitely if short jobs keep arriving: this is \cle{starvation}. FCFS serves jobs in arrival order and Round Robin gives every process a quantum, so neither can starve a process permanently.
\item \textbf{Answer (d).} A deadlock requires the \emph{four} Coffman conditions to hold simultaneously: mutual exclusion, hold and wait, no preemption, and circular wait. If any single one is broken, deadlock is impossible.
\end{enumerate}
\end{corrige}

\begin{corrige}[Exercise 2 — True or false]
\begin{enumerate}
\item \textbf{False.} By definition, a \emph{non-preemptive} scheduler lets the running process keep the CPU until it terminates or blocks voluntarily (e.g.\ for an I/O request). Forcing it out before that is precisely what \emph{preemption} means.
\item \textbf{True.} During a \cle{context switch} the CPU saves the PCB of the old process and loads the PCB of the new one; no user instruction is executed. It is pure overhead, which is why the quantum in Round Robin must not be too small.
\item \textbf{True.} A binary semaphore \texttt{mutex} initialised to $1$ guarantees mutual exclusion: the first process that executes \texttt{wait(mutex)} enters its critical section, and any other process executing \texttt{wait(mutex)} is blocked until \texttt{signal(mutex)} is called.
\item \textbf{False.} A cycle in a Resource Allocation Graph is a \emph{necessary} condition for deadlock only when each resource type has a single instance. If resources have several instances, a cycle may exist without deadlock, because a process in the cycle may be completed using an instance held by a process \emph{outside} the cycle.
\end{enumerate}
\end{corrige}

\begin{corrige}[Exercise 3 — Definitions]
\begin{enumerate}
\item \textbf{Turnaround time:} the total time elapsed between the submission (arrival) of a process and its completion: $T_{\text{turn}} = t_{\text{completion}} - t_{\text{arrival}}$. It includes waiting, CPU execution and I/O.
\item \textbf{Waiting time:} the total time a process spends in the ready queue, i.e.\ waiting for the CPU: $T_{\text{wait}} = T_{\text{turn}} - T_{\text{burst}}$ (for a process that does no I/O).
\item \textbf{Response time:} the time between the arrival of a process and the \emph{first} moment it receives the CPU: $T_{\text{resp}} = t_{\text{first run}} - t_{\text{arrival}}$. It measures the reactivity perceived by an interactive user.
\item \textbf{Race condition:} a situation in which several processes access shared data concurrently and the final result depends on the exact order (interleaving) of their machine instructions, so that some executions give an incorrect result.
\item \textbf{Deadlock:} a situation in which a set of processes are all blocked, each one waiting for a resource held by another process of the set, so that none of them can ever proceed without external intervention.
\end{enumerate}
\end{corrige}

\begin{corrige}[Exercise 4 — Direct calculation]
\textbf{(a) FCFS in the order $P_1, P_2, P_3$.} All processes arrive at time $0$, so waiting time of a process is the sum of the burst times of the processes executed before it.
\begin{center}
\begin{tabular}{|c|c|c|c|}
\hline
\rowcolor{popblueL} Process & Burst & Waiting time & Turnaround time \\
\hline
$P_1$ & 8 & $0$ & $0+8=8$ \\
\hline
$P_2$ & 4 & $8$ & $8+4=12$ \\
\hline
$P_3$ & 2 & $8+4=12$ & $12+2=14$ \\
\hline
\end{tabular}
\end{center}
\textbf{(b) Average turnaround time:}
\[ \bar{T}_{\text{turn}} = \frac{8+12+14}{3} = \frac{34}{3} \approx 11.33\ \text{ms}. \]
\textbf{(c) Order $P_3, P_2, P_1$:}
\begin{center}
\begin{tabular}{|c|c|c|c|}
\hline
\rowcolor{popblueL} Process & Burst & Waiting time & Turnaround time \\
\hline
$P_3$ & 2 & $0$ & $2$ \\
\hline
$P_2$ & 4 & $2$ & $6$ \\
\hline
$P_1$ & 8 & $2+4=6$ & $14$ \\
\hline
\end{tabular}
\end{center}
\[ \bar{T}_{\text{turn}} = \frac{2+6+14}{3} = \frac{22}{3} \approx 7.33\ \text{ms}. \]
\textbf{Observation:} the average turnaround time drops from $11.33$ ms to $7.33$ ms simply by running the shortest jobs first. This is exactly the idea behind the \cle{SJF} algorithm: serving short jobs early reduces the time that many processes spend waiting behind a long job.
\end{corrige}

\begin{corrige}[Exercise 5 — FCFS and SJF Gantt charts]
\textbf{(a) FCFS.} Processes are served in arrival order $P_1, P_2, P_3, P_4$:
\begin{center}
\begin{tabular}{|c|c|c|c|}
\hline
\rowcolor{popblueL} $P_1$ & $P_2$ & $P_3$ & $P_4$ \\
\hline
$0$ -- $7$ & $7$ -- $11$ & $11$ -- $12$ & $12$ -- $16$ \\
\hline
\end{tabular}
\end{center}
Waiting times: $P_1: 0$; $P_2: 7-2=5$; $P_3: 11-4=7$; $P_4: 12-5=7$.
\[ \bar{T}_{\text{wait}} = \frac{0+5+7+7}{4} = \frac{19}{4} = 4.75\ \text{ms}. \]
\textbf{(b) Non-preemptive SJF.} At $t=0$ only $P_1$ has arrived, so it runs first (non-preemptive: it cannot be interrupted). At $t=7$, the ready processes are $P_2$ (burst 4), $P_3$ (burst 1), $P_4$ (burst 4): the shortest is $P_3$, then $P_2$ and $P_4$ (tie broken by arrival order).
\begin{center}
\begin{tabular}{|c|c|c|c|}
\hline
\rowcolor{popblueL} $P_1$ & $P_3$ & $P_2$ & $P_4$ \\
\hline
$0$ -- $7$ & $7$ -- $8$ & $8$ -- $12$ & $12$ -- $16$ \\
\hline
\end{tabular}
\end{center}
Waiting times: $P_1: 0$; $P_3: 7-4=3$; $P_2: 8-2=6$; $P_4: 12-5=7$.
\[ \bar{T}_{\text{wait}} = \frac{0+3+6+7}{4} = \frac{16}{4} = 4\ \text{ms}. \]
\textbf{(c)} SJF performs better here ($4$ ms against $4.75$ ms average waiting time). This illustrates the general result: SJF gives the minimum average waiting time for a given set of jobs, because every job that runs early adds its burst to the waiting time of fewer processes.
\end{corrige}

\begin{corrige}[Exercise 6 — Round Robin]
\textbf{(a) Gantt chart} (quantum $q=4$ ms; a process that finishes its quantum with remaining burst goes to the back of the queue):
\begin{center}
\begin{tabular}{|c|c|c|c|c|c|c|c|c|c|}
\hline
\rowcolor{popblueL} $P_1$ & $P_2$ & $P_3$ & $P_4$ & $P_5$ & $P_1$ & $P_3$ & $P_4$ & $P_1$ & $P_5$ \\
\hline
$0$--$4$ & $4$--$7$ & $7$--$11$ & $11$--$15$ & $15$--$19$ & $19$--$23$ & $23$--$25$ & $25$--$29$ & $29$--$31$ & $31$--$32$ \\
\hline
\end{tabular}
\end{center}
(Remaining bursts: after the first round $P_1\!:6$, $P_2\!:0$, $P_3\!:2$, $P_4\!:4$, $P_5\!:1$; then $P_1$ uses $4$ more, $P_3$ finishes with $2$, $P_4$ finishes with $4$, $P_1$ finishes with $2$, $P_5$ finishes with $1$.)
\textbf{(b) Response times} (time of first allocation of the CPU, since all arrive at $0$):
\begin{center}
\begin{tabular}{|c|c|c|c|c|c|}
\hline
\rowcolor{popblueL} Process & $P_1$ & $P_2$ & $P_3$ & $P_4$ & $P_5$ \\
\hline
Response time (ms) & $0$ & $4$ & $7$ & $11$ & $15$ \\
\hline
\end{tabular}
\end{center}
\textbf{(c) Turnaround times} (completion times, since arrival is at $0$): $P_1\!:31$, $P_2\!:7$, $P_3\!:25$, $P_4\!:29$, $P_5\!:32$.
\[ \bar{T}_{\text{turn}} = \frac{31+7+25+29+32}{5} = \frac{124}{5} = 24.8\ \text{ms}. \]
\textbf{(d) Effect of $q=2$ ms.} Each process is reached by the scheduler sooner after its arrival, so the \emph{response time decreases}: the system feels more reactive. However, the number of quanta (and therefore of \cle{context switches}) roughly doubles; since each context switch is pure overhead, a larger fraction of CPU time is wasted, which increases the total completion time of the batch. The quantum must therefore be chosen as a compromise: large enough to keep overhead small, small enough to keep the system responsive.
\end{corrige}

\begin{corrige}[Exercise 7 — The lost increment]
\textbf{(a)} Let $R_A$ and $R_B$ be the registers of $A$ and $B$. Consider the interleaving:
\begin{enumerate}
\item $A$: \texttt{LOAD} $R_A \leftarrow \texttt{count}$ \quad ($R_A = 0$)
\item $B$: \texttt{LOAD} $R_B \leftarrow \texttt{count}$ \quad ($R_B = 0$)
\item $A$: \texttt{ADD 1} \quad ($R_A = 1$)
\item $A$: \texttt{STORE} $\texttt{count} \leftarrow R_A$ \quad ($\texttt{count} = 1$)
\item $B$: \texttt{ADD 1} \quad ($R_B = 1$)
\item $B$: \texttt{STORE} $\texttt{count} \leftarrow R_B$ \quad ($\texttt{count} = 1$)
\end{enumerate}
Both processes read the value $0$ before either writes back, so the final value is $1$ instead of $2$: one increment is lost.
\textbf{(b)} This is a \cle{race condition} (more precisely a \emph{lost update}): the result depends on the interleaving of the instructions of the two processes.
\textbf{(c)} Use a binary semaphore \texttt{mutex} initialised to $1$; the update becomes a critical section:
\begin{itemize}
\item \texttt{Process A (identically B):}
\item \texttt{wait(mutex);}
\item \texttt{count := count + 1;}
\item \texttt{signal(mutex);}
\end{itemize}
Now whichever process executes \texttt{wait(mutex)} first completes all three machine steps before the other can read \texttt{count}, so the final value is always $2$.
\end{corrige}

\begin{corrige}[Exercise 8 — Producer–consumer with a bounded buffer]
\textbf{(a) Semaphores:}
\begin{itemize}
\item \texttt{empty} $= 5$: counts the free slots (initially the whole buffer is empty);
\item \texttt{full} $= 0$: counts the occupied slots;
\item \texttt{mutex} $= 1$: binary semaphore protecting the buffer itself (mutual exclusion).
\end{itemize}
\textbf{(b) Pseudocode:}
\begin{itemize}
\item \texttt{Producer: loop forever}
\item \texttt{\quad produce an item;}
\item \texttt{\quad wait(empty); \quad\quad\quad\ \% wait for a free slot}
\item \texttt{\quad wait(mutex); \quad\quad\quad\ \% enter critical section}
\item \texttt{\quad insert item into buffer;}
\item \texttt{\quad signal(mutex);}
\item \texttt{\quad signal(full); \quad\quad\quad\ \% one more occupied slot}
\item \texttt{Consumer: loop forever}
\item \texttt{\quad wait(full); \quad\quad\quad\quad \% wait for an occupied slot}
\item \texttt{\quad wait(mutex);}
\item \texttt{\quad remove item from buffer;}
\item \texttt{\quad signal(mutex);}
\item \texttt{\quad signal(empty); \quad\quad\quad \% one more free slot}
\item \texttt{\quad consume the item;}
\end{itemize}
\textbf{(c) Why the order matters.} Suppose the producer took \texttt{mutex} \emph{before} \texttt{empty}. If the buffer is full, the producer blocks on \texttt{wait(empty)} \emph{while still holding} \texttt{mutex}. The consumer, which is the only process able to free a slot, then blocks on \texttt{wait(mutex)}. Each process waits for the other forever: a \cle{deadlock}. The rule is therefore: \emph{always wait on the resource-counting (synchronisation) semaphore first, and on the mutual-exclusion semaphore last}.
\end{corrige}

\begin{corrige}[Exercise 9 — Scheduling comparison (GCE style)]
\textbf{(a) Definitions (2 marks — 1 per definition).} \emph{Waiting time}: total time a process spends in the ready queue waiting for the CPU. \emph{Turnaround time}: total time from arrival to completion, $T_{\text{turn}} = t_{\text{completion}} - t_{\text{arrival}}$.

\textbf{(b) (i) FCFS (4 marks).} Order $J_1, J_2, J_3, J_4$:
\begin{center}
\begin{tabular}{|c|c|c|c|}
\hline
\rowcolor{popblueL} $J_1$ & $J_2$ & $J_3$ & $J_4$ \\
\hline
$0$--$8$ & $8$--$12$ & $12$--$21$ & $21$--$26$ \\
\hline
\end{tabular}
\end{center}
Turnaround: $J_1\!:8$, $J_2\!:12-1=11$, $J_3\!:21-2=19$, $J_4\!:26-3=23$; average $= 61/4 = 15.25$ ms.
Waiting: $0,\ 7,\ 10,\ 18$; average $= 35/4 = 8.75$ ms.

\textbf{(ii) Non-preemptive SJF (4 marks).} At $t=0$ only $J_1$ is present, so it runs to completion. At $t=8$: $J_2$ (4), $J_4$ (5), $J_3$ (9) are ready, in that SJF order.
\begin{center}
\begin{tabular}{|c|c|c|c|}
\hline
\rowcolor{popblueL} $J_1$ & $J_2$ & $J_4$ & $J_3$ \\
\hline
$0$--$8$ & $8$--$12$ & $12$--$17$ & $17$--$26$ \\
\hline
\end{tabular}
\end{center}
Turnaround: $J_1\!:8$, $J_2\!:11$, $J_4\!:14$, $J_3\!:24$; average $= 57/4 = 14.25$ ms.
Waiting: $0,\ 7,\ 9,\ 15$; average $= 31/4 = 7.75$ ms.

\textbf{(iii) Preemptive SJF / SRTF (5 marks).} At $t=1$, $J_2$ (burst $4$) arrives while $J_1$ has $7$ ms remaining: $J_2$ preempts. At $t=2$ and $t=3$, $J_3$ (9) and $J_4$ (5) arrive but $J_2$'s remaining time ($3$, then $2$) stays the smallest. At $t=5$, $J_2$ finishes; the remaining times are $J_4\!:5$, $J_1\!:7$, $J_3\!:9$, so $J_4$ runs to completion, then $J_1$, then $J_3$.
\begin{center}
\begin{tabular}{|c|c|c|c|c|}
\hline
\rowcolor{popblueL} $J_1$ & $J_2$ & $J_4$ & $J_1$ & $J_3$ \\
\hline
$0$--$1$ & $1$--$5$ & $5$--$10$ & $10$--$17$ & $17$--$26$ \\
\hline
\end{tabular}
\end{center}
Turnaround: $J_1\!:17$, $J_2\!:5-1=4$, $J_3\!:26-2=24$, $J_4\!:10-3=7$; average $= 52/4 = 13$ ms.
Waiting: $J_1\!:17-8=9$, $J_2\!:0$, $J_3\!:24-9=15$, $J_4\!:7-5=2$; average $= 26/4 = 6.5$ ms.

\textbf{(c) (2 marks).} SRTF gives the best average waiting time ($6.5$ ms). Justification: SJF/SRTF is \emph{optimal} with respect to average waiting time because always running the job with the smallest (remaining) burst minimises the sum of the waiting times — any exchange that places a longer job before a shorter one increases the total.

\textbf{(d) (1 mark).} In practice the CPU burst of a job is \emph{not known in advance} (it can only be estimated); moreover long jobs may suffer \cle{starvation} if short jobs keep arriving.

\textbf{(e) Round Robin, $q=3$ (4 marks).} Ready queue discipline: newly arrived jobs join the queue, then the preempted job is requeued behind them.
\begin{center}
\begin{tabular}{|c|c|c|c|c|c|c|c|c|}
\hline
\rowcolor{popblueL} $J_1$ & $J_2$ & $J_3$ & $J_4$ & $J_1$ & $J_2$ & $J_3$ & $J_4$ & $J_1$ \\
\hline
$0$--$3$ & $3$--$6$ & $6$--$9$ & $9$--$12$ & $12$--$15$ & $15$--$16$ & $16$--$19$ & $19$--$21$ & $21$--$23$ \\
\hline
\end{tabular}
\end{center}
then $J_3$ runs $23$--$26$ and finishes. (Remaining bursts: after first pass $J_1\!:5$, $J_2\!:1$, $J_3\!:6$, $J_4\!:2$.)

Response times ($t_{\text{first run}} - t_{\text{arrival}}$): $J_1\!: 0-0 = 0$ ms; $J_2\!: 3-1 = 2$ ms; $J_3\!: 6-2 = 4$ ms; $J_4\!: 9-3 = 6$ ms.

\textbf{Examiner's expectations:} a correct Gantt chart with labelled time boundaries (essential for the method marks), waiting/turnaround computed \emph{from the chart}, correct averages with units, and for (c) the keyword \emph{optimal} linked to shortest-job-first reasoning.
\end{corrige}

\begin{corrige}[Exercise 10 — Banker's algorithm and RAG (GCE style)]
\textbf{(a) Need matrix (2 marks).} $\text{Need} = \text{Max} - \text{Allocation}$:
\begin{center}
\begin{tabular}{|c|ccc|}
\hline
\rowcolor{popblueL} Process & $R_1$ & $R_2$ & $R_3$ \\
\hline
$P_1$ & 7 & 4 & 3 \\
\hline
$P_2$ & 1 & 2 & 2 \\
\hline
$P_3$ & 6 & 0 & 0 \\
\hline
$P_4$ & 0 & 1 & 1 \\
\hline
\end{tabular}
\end{center}
\textbf{(b) Safety algorithm (6 marks).} Work $= (3,3,2)$.
\begin{enumerate}
\item $P_2$: Need $(1,2,2) \le (3,3,2)$ \checkmark \quad Work $\leftarrow (3,3,2)+(2,0,0) = (5,3,2)$.
\item $P_4$: Need $(0,1,1) \le (5,3,2)$ \checkmark \quad Work $\leftarrow (5,3,2)+(2,1,1) = (7,4,3)$.
\item $P_1$: Need $(7,4,3) \le (7,4,3)$ \checkmark \quad Work $\leftarrow (7,4,3)+(0,1,0) = (7,5,3)$.
\item $P_3$: Need $(6,0,0) \le (7,5,3)$ \checkmark \quad Work $\leftarrow (7,5,3)+(3,0,2) = (10,5,5)$.
\end{enumerate}
All processes can finish: the system is in a \cle{safe state}, with safe sequence $\langle P_2, P_4, P_1, P_3 \rangle$ (other safe sequences exist, e.g.\ $\langle P_4, P_2, P_1, P_3 \rangle$).

\textbf{(c) Request $(1,1,0)$ by $P_2$ (4 marks).}
\begin{enumerate}
\item Request $(1,1,0) \le \text{Need}_{P_2} = (1,2,2)$ \checkmark \ (the request is legal).
\item Request $(1,1,0) \le \text{Available} = (3,3,2)$ \checkmark \ (resources exist).
\item \emph{Pretend} to grant: Available $= (2,2,2)$; $\text{Alloc}_{P_2} = (3,1,0)$; $\text{Need}_{P_2} = (0,1,2)$.
\item Re-run safety: $P_2$: $(0,1,2)\le(2,2,2)$ \checkmark, Work $=(5,3,2)$; $P_4$: $(0,1,1)\le(5,3,2)$ \checkmark, Work $=(7,4,3)$; $P_1$: $(7,4,3)\le(7,4,3)$ \checkmark, Work $=(7,5,3)$; $P_3$: $(6,0,0)\le(7,5,3)$ \checkmark, Work $=(10,5,5)$.
\end{enumerate}
The resulting state is still safe (sequence $\langle P_2, P_4, P_1, P_3 \rangle$), so the request \textbf{can be granted immediately}.

\textbf{(d) Coffman conditions (4 marks — 1 each).} (i) \emph{Mutual exclusion}: at least one resource is held in a non-sharable mode; (ii) \emph{Hold and wait}: a process holds resources while requesting others; (iii) \emph{No preemption}: resources cannot be forcibly taken, only released voluntarily; (iv) \emph{Circular wait}: a cycle of processes each waiting for a resource held by the next.

\textbf{(e) Resource Allocation Graph (4 marks).}
\begin{popfigure}
\begin{tikzpicture}[scale=1.1, every node/.style={font=\small}]
\node[draw, circle, minimum size=7mm] (p1) at (0,2) {$P_1$};
\node[draw, circle, minimum size=7mm] (p2) at (4,2) {$P_2$};
\node[draw, circle, minimum size=7mm] (p3) at (2,0) {$P_3$};
\node[draw, rectangle, minimum size=6mm] (r1) at (2,3.4) {$R_1$};
\node[draw, rectangle, minimum size=6mm] (r2) at (4.6,0.6) {$R_2$};
\node[draw, rectangle, minimum size=6mm] (r3) at (-0.6,0.6) {$R_3$};
\draw[->, thick, popblue] (r1) -- node[above left, font=\scriptsize]{held} (p1);
\draw[->, thick, poporange] (p1) -- node[above, font=\scriptsize]{requests} (r2);
\draw[->, thick, popblue] (r2) -- (p2);
\draw[->, thick, poporange] (p2) -- (r3);
\draw[->, thick, popblue] (r3) -- (p3);
\draw[->, thick, poporange] (p3) -- (r1);
\end{tikzpicture}
\legende{Resource Allocation Graph: assignment edges $R_i \to P_i$ (blue) and request edges $P_i \to R_j$ (orange).}
\end{popfigure}
The graph contains the cycle $P_1 \to R_2 \to P_2 \to R_3 \to P_3 \to R_1 \to P_1$. Since each resource type has \emph{exactly one instance}, a cycle here is a \emph{sufficient} condition: a \textbf{deadlock exists}. Justification via Coffman: each resource is non-sharable (mutual exclusion); each process holds one resource while requesting another (hold and wait); resources cannot be preempted (no preemption); and the cycle above is a circular wait. All four conditions hold simultaneously, so none of $P_1, P_2, P_3$ can ever proceed.

\textbf{Examiner's expectations:} the Need matrix computed as Max $-$ Allocation (not the reverse); the safety algorithm shown \emph{step by step} with the evolving Work vector; for (c), the three checks (request $\le$ Need, request $\le$ Available, safety of the hypothetical state) — simply saying ``resources are available'' earns no credit; for (e), correct orientation of the two edge types.
\end{corrige}

\begin{corrige}[Exercise 11 — Structured essay]
\textbf{(a) Definitions (3 marks — 1 each).} A \emph{process} is a program in execution, together with its current state (program counter, registers, stack, data, open files). The \emph{Process Control Block} (PCB) is the data structure, maintained by the operating system, that stores all the information about one process: its identifier, state, program counter, register values, memory and open-file information. A \emph{context switch} is the operation by which the CPU passes from one process to another: the state of the outgoing process is saved in its PCB and the state of the incoming process is restored from its own PCB.

\textbf{(b) Context switch as overhead (3 marks).} During a context switch the CPU executes only operating-system code (saving and restoring registers, updating queues); \emph{no user program progresses}: it is pure overhead (1 mark). Consequences on scheduling design (any two, 1 mark each): the Round Robin quantum must not be too small, otherwise the proportion of time wasted in switches becomes unacceptable; a scheduler should avoid unnecessary preemptions, so non-preemptive or weakly preemptive policies may be preferred for batch work; the number of processes competing for the CPU (degree of multiprogramming) must be controlled.

\textbf{(c) Preemptive vs non-preemptive scheduling for the bank (8 marks).} Indicative mark scheme: 2 marks per correctly discussed criterion.
\begin{itemize}
\item \emph{Responsiveness.} Interactive tellers and customers at terminals need answers within a second or so. A non-preemptive scheduler would let a long job (e.g.\ printing statements) monopolise the CPU and freeze all terminals; a preemptive scheduler (Round Robin or priority with preemption) interrupts long jobs and gives each interactive session frequent slices of CPU, guaranteeing a short \cle{response time}.
\item \emph{Overhead.} Preemption costs a context switch each time; with a small quantum the server spends a noticeable fraction of its power switching instead of computing. Non-preemptive scheduling has almost no such overhead.
\item \emph{Starvation.} Under preemptive priority or SJF-like policies, long batch jobs may be perpetually postponed by arriving interactive requests (starvation); a remedy is \emph{aging} (gradually raising the priority of waiting processes). FCFS non-preemptive scheduling has no starvation but terrible response times (the \emph{convoy effect}).
\item \emph{Predictability.} End-of-day account processing must finish before the next morning; non-preemptive execution of these batch jobs (e.g.\ at night) gives predictable completion times, whereas preemption makes their duration depend on the interactive load.
\end{itemize}
\emph{Conclusion (justified recommendation):} a \textbf{preemptive} discipline such as Round Robin with a moderate quantum (or multilevel queues: interactive jobs in a high-priority RR queue, batch jobs in a lower FCFS queue scheduled mainly at night) is recommended, because for a banking terminal system the perceived responsiveness is the critical criterion, and the modest context-switch overhead is an acceptable price; aging protects the batch jobs from starvation.

\textbf{(d) Shared counter and semaphore (6 marks).} The update \texttt{counter := counter + 1} compiles to \texttt{LOAD}, \texttt{ADD 1}, \texttt{STORE}. With \texttt{counter} $= 100$, consider the interleaving: process $A$ executes \texttt{LOAD} (reads $100$); the scheduler preempts $A$; process $B$ executes \texttt{LOAD} (reads $100$), \texttt{ADD 1}, \texttt{STORE} (writes $101$); $A$ resumes with \texttt{ADD 1} (its register becomes $101$) and \texttt{STORE} (writes $101$). Two transactions were recorded but the counter increased only by $1$: a \cle{race condition} / lost update (3 marks: 1 for the three machine steps, 2 for a correct interleaving with explanation). Solution: a binary semaphore \texttt{mutex} initialised to $1$; each process executes \texttt{wait(mutex)} before the update and \texttt{signal(mutex)} after it (2 marks). The critical section is then executed atomically: the second process cannot read the counter until the first has written it back, so the final value is always correct (1 mark).
\end{corrige}

\newpage

\coursec{Summary sheet}

\begin{popfigure}
\begin{center}\tcbox[colback=popink,colframe=popink,coltext=white,arc=9pt,boxsep=4pt,left=6pt,right=6pt]{\bfseries\small Managing Processes (CSC10)}\end{center}
\vspace{-2pt}
\begin{tcbraster}[raster columns=3,raster equal height=rows,raster column skip=6pt,raster row skip=6pt,enhanced,blankest]
\begin{tcolorbox}[enhanced,colback=popblue,colframe=popblue,coltext=white,boxrule=0.9pt,arc=3pt,left=4pt,right=4pt,top=3pt,bottom=3pt,boxsep=1pt,halign=left]\bfseries\scriptsize Process basics: definition, PCB, states, transitions\end{tcolorbox}
\begin{tcolorbox}[enhanced,colback=popgreen,colframe=popgreen,coltext=white,boxrule=0.9pt,arc=3pt,left=4pt,right=4pt,top=3pt,bottom=3pt,boxsep=1pt,halign=left]\bfseries\scriptsize Scheduling: goals, criteria, queues\end{tcolorbox}
\begin{tcolorbox}[enhanced,colback=poporange,colframe=poporange,coltext=white,boxrule=0.9pt,arc=3pt,left=4pt,right=4pt,top=3pt,bottom=3pt,boxsep=1pt,halign=left]\bfseries\scriptsize Non-preemptive: FCFS, SJF, priority\end{tcolorbox}
\begin{tcolorbox}[enhanced,colback=poppurple,colframe=poppurple,coltext=white,boxrule=0.9pt,arc=3pt,left=4pt,right=4pt,top=3pt,bottom=3pt,boxsep=1pt,halign=left]\bfseries\scriptsize Preemptive: SRTF, Round Robin\end{tcolorbox}
\begin{tcolorbox}[enhanced,colback=popgold,colframe=popgold,coltext=white,boxrule=0.9pt,arc=3pt,left=4pt,right=4pt,top=3pt,bottom=3pt,boxsep=1pt,halign=left]\bfseries\scriptsize Synchronisation: race conditions, semaphores\end{tcolorbox}
\begin{tcolorbox}[enhanced,colback=popdark,colframe=popdark,coltext=white,boxrule=0.9pt,arc=3pt,left=4pt,right=4pt,top=3pt,bottom=3pt,boxsep=1pt,halign=left]\bfseries\scriptsize Deadlock: 4 conditions, prevention, avoidance\end{tcolorbox}
\end{tcbraster}

\legende{Mind map of the chapter: from the life of a process to deadlock handling.}
\end{popfigure}

\begin{bilan}

\textbf{1. Key definitions}
\begin{itemize}
\item A \cle{process} is a program in execution, with its own code, data, stack and context. A program is passive; a process is active.
\item The \cle{Process Control Block} (PCB) is the data structure the operating system keeps for each process: process ID, state, program counter, register values, memory information, open files, scheduling data.
\item The main \cle{process states} are: New, Ready, Running, Blocked (Waiting), Terminated. Transitions: admit (New $\to$ Ready), dispatch (Ready $\to$ Running), interrupt/time-out (Running $\to$ Ready), I/O request (Running $\to$ Blocked), I/O completion (Blocked $\to$ Ready), exit (Running $\to$ Terminated).
\item A \cle{context switch} saves the state of the current process in its PCB and loads the state of the next one; it is pure overhead.
\item \cle{CPU scheduling} decides which ready process gets the processor. \cle{Non-preemptive}: a running process keeps the CPU until it finishes or blocks. \cle{Preemptive}: the OS may take the CPU back (interrupt, higher priority, end of time quantum).
\item A \cle{race condition} occurs when several processes access shared data concurrently and the result depends on the order of execution. A \cle{critical section} is the code region accessing the shared resource; it must run in \cle{mutual exclusion}.
\item A \cle{semaphore} is an integer variable accessed only through the atomic operations \texttt{wait(P)} and \texttt{signal(V)}. A binary semaphore (mutex) values 0 or 1.
\item A \cle{deadlock} is a situation where a set of processes each wait for a resource held by another member of the set, so none can ever proceed.
\end{itemize}

\textbf{2. Formulas and results to know}
\begin{itemize}
\item Turnaround time: $T_{\text{turn}} = T_{\text{completion}} - T_{\text{arrival}}$.
\item Waiting time: $T_{\text{wait}} = T_{\text{turn}} - T_{\text{burst}}$.
\item Response time: $T_{\text{resp}} = T_{\text{first run}} - T_{\text{arrival}}$.
\item CPU utilisation $= \dfrac{\text{busy time}}{\text{total time}} \times 100\ \%$; throughput $= \dfrac{\text{number of processes completed}}{\text{time}}$.
\item FCFS: average waiting time depends heavily on arrival order (convoy effect). SJF/SRTF give the minimum average waiting time for a given set of bursts (no proof required at GCE, but you must be able to exploit it).
\item Round Robin with quantum $q$: a process with burst $B$ needs $\lceil B/q \rceil$ slices; small $q$ $\Rightarrow$ good response but many context switches; large $q$ $\Rightarrow$ behaves like FCFS.
\item Coffman's four \cle{necessary conditions} for deadlock: mutual exclusion, hold and wait, no preemption, circular wait. All four must hold simultaneously.
\end{itemize}

\textbf{3. Methods to master}
\begin{enumerate}
\item \textbf{Gantt chart method}: list processes with arrival and burst times; draw the CPU timeline; read off completion times; compute turnaround, waiting and response times; average over all processes.
\item \textbf{Choosing the running process}: FCFS $\to$ earliest arrival; SJF $\to$ shortest burst among arrived processes; SRTF $\to$ shortest \emph{remaining} time, rechecked at every arrival; RR $\to$ circular ready queue, quantum $q$; priority $\to$ highest priority (state the convention: small number $=$ high priority).
\item \textbf{Protecting a critical section}: \texttt{wait(mutex)} before entering, critical section, \texttt{signal(mutex)} after leaving. Producer--consumer: semaphores \texttt{empty} (initialised to $N$), \texttt{full} (initialised to $0$), plus \texttt{mutex}.
\item \textbf{Deadlock handling}: prevention (break one Coffman condition, e.g.\ impose a global resource ordering), avoidance (Banker's algorithm: grant only if the resulting state is \emph{safe}), detection and recovery (resource allocation graph, kill or roll back a victim).
\item \textbf{Banker's safety check}: a state is safe if there exists an order in which every process can obtain its remaining need ($\text{Need} = \text{Max} - \text{Allocation}$) and finish.
\end{enumerate}

\textbf{4. Common mistakes (examiner's warnings)}
\begin{itemize}
\item Confusing waiting time and turnaround time: always subtract the burst for waiting time.
\item In SRTF, forgetting to re-examine remaining times when a new process arrives.
\item In Round Robin, forgetting that a process arriving \emph{during} a quantum joins the back of the queue, and that a process finishing exactly at the quantum boundary leaves before the new arrival is enqueued (state your convention clearly).
\item Believing SJF is practical: the burst time is unknown in advance; it is only estimated.
\item Writing \texttt{wait}/\texttt{signal} in the wrong order, or placing them around the wrong code region: this destroys mutual exclusion or causes deadlock.
\item Thinking one Coffman condition suffices: deadlock requires \emph{all four}.
\end{itemize}

\textbf{5. GCE tips}
\begin{itemize}
\item Always draw the Gantt chart first: every mark for averages follows from it. Show the formula, the substitution, then the result.
\item Read the question for the tie-breaking rule (usually FCFS among equals) and the priority convention.
\item Keep one decimal place (or the requested precision) in averages, e.g.\ $6.5$ ms, not $6.499\ldots$
\item For synchronisation questions, name the shared resource, identify the critical section, then justify the semaphore initial values.
\item For deadlock questions, quote the four Coffman conditions explicitly and say which one your proposed solution removes.
\item Manage your time: a full scheduling table (FCFS, SJF, RR) is a classic 8--10 mark question; practise until a Gantt chart takes under five minutes.
\end{itemize}

\end{bilan}

\findecours

\end{document}
